\section{The volume theorem}\label{sec:pf:volthm}

\providecommand{\volthmdeg}{\deg_{\mathrm{Fl}}}
\providecommand{\volthmlin}{\mathrm{lin}}
\providecommand{\volthmfd}{\mathrm{fdeg}}
\providecommand{\volthmII}{\mathcal{I}}
\providecommand{\volthmhh}{h_{\mathrm{Fl}}}
\providecommand{\volthmSh}{\mathrm{Sh}}
\providecommand{\volthmTot}{\mathrm{Tot}}
\providecommand{\volthmPP}{\mathrm{PP}}
\providecommand{\volthmgap}{\mathrm{gap}}
\providecommand{\volthmR}{\mathcal{R}}
\providecommand{\volthmcoef}{\varpi}

This section proves the \emph{volume theorem} $\sfname{VolumeTheoremX}$ of \cref{def:covolume:tower}
from a list of geometric and arithmetic inputs (\cref{pf:volthm:main}). Together with the existence of
good towers it makes up \cref{thm:tower}; the theorem proved here is used, with its hypotheses discharged, in
\lref{GlobalTheory.CompanionTheory.Ar.Top.goodTower_and_volume}.

\emph{What is proved.} A root tower $T$ attaches to every cutoff $B\in\N$ and every component $c$ a
\emph{root model} $M_{B,c}$: an arithmetic variety of dimension $n+1$, where $n$ is the dimension of the
flag variety of type $(\Psi,N)$, carrying Hermitian lines $x_t(\tau,i)$ (``Chern roots'', depending on a
regularisation parameter $t>0$) and a reference line $A$. Arithmetic intersection numbers of these lines
define, in the limit $t\to0^+$, a linear functional $\volthmPP_{B,c}$ on polynomials in the Chern roots, and
geometric degrees define a functional $\mathrm{geo}_{B,c}$. The volume theorem asserts, uniformly over abelian
CM fields, levels and signatures $(p,q)$ with $p\ge2q+1$: proportionality of $\mathrm{geo}_{B,c}$ with the
degree on the flag variety, equality of the volumes of the components, limit formulas as the regularisation
is removed, an asymptotic formula for $\sum_c\volthmPP_{B,c}$ with error $O(\log B/B)$, and a lower bound
\begin{equation}\label{pf:volthm:eq:intro}
\frac1{\#\mathrm{Comp}}\sum_c\mathrm{slope}(B,c,u)\;\ge\;-C\bigl(\Lambda(B)\,\|u\|_1+\min(j,N-j)\,\mathrm{spread}(u)\bigr),
\end{equation}
where $\Lambda(B)=\log(B\,\mathrm{rd}_E)+\bigl|\Re\tfrac{L'}{L}(1,\chi_d)\bigr|$,
for the mean arithmetic slope of automorphic line bundles of regular weight $u$, with a constant $C$ independent
of all data.

\emph{Structure of the proof.} The proof has two layers.
\begin{enumerate}[label=(\arabic*)]
\item An abstract \emph{tower theorem} (\cref{pf:volthm:thm:tower}) deduces the conclusions from sixteen
  conditions on a root tower. It is a statement about linear functionals on a polynomial ring and their
  dependence on $B$, and its proof is given in full in
  \cref{pf:volthm:sec:poly,pf:volthm:sec:models,pf:volthm:sec:weil,pf:volthm:sec:asym,pf:volthm:sec:lower}; the bound for
  the remainder of the test polynomial is reduced in \cref{pf:volthm:sec:remainder} to a combinatorial bound for its
  terms.
  The steps are: a \emph{Weil-type bound} on walls (on weights with a vanishing gap the functional is
  $O(\log B/B)$; this comes from squeezing arithmetic self-intersections between Zhang's lower and upper
  bounds by heights of points, while the geometric degree vanishes on the wall); an \emph{antisymmetry}
  argument, which turns the wall bound into the asymptotic formula; \emph{Galois averaging}, which reduces
  the asymptotic formula to two scalar parameters; an \emph{Euler product identity} relating the cutoffs $B_0$
  and $B$, which in the limit $B\to\infty$ gives an exact identity at every cutoff, with two constants
  $A_0,b_0$ and convergent prime sums as remainders; and a \emph{test polynomial} on the Hodge locus, whose
  arithmetic degree is bounded below by a lower bound for heights (the slope condition) and above by the
  arithmetic height of the flag variety, which bounds $A_0$ from below. The lower bound
  \eqref{pf:volthm:eq:intro} then follows from the identity at the cutoff $B$.
\item The \emph{verification of the conditions} for root towers with good origins
  (\cref{pf:volthm:sec:conditions}). Each condition is the conclusion of a theorem of the development; we
  give the reductions and state precisely the geometric results they rest on. The main geometric results are
  proved at their top levels in the following subsections: the linking of the root models at two cutoffs and the
  height change formula behind the Euler product identity (\cref{pf:volthm:sec:linking}), the base degree settings
  on the Hodge locus behind the slope condition (\cref{pf:volthm:sec:basesettings}), the equality of volumes and the
  flat data behind condition $\mathrm{FL}$ (\cref{pf:volthm:sec:voleq}), and the bounds for heights of points
  ($\mathrm{I}_5$, $\mathrm{I}_6$), the antisymmetry identity ($\mathrm{I}_8$) and the semipositivity ($\mathrm{A}_1$)
  (\cref{pf:volthm:sec:deep}). Their deepest ingredients (analytic estimates on the Hodge locus, the construction of
  models and the combinatorics of the remainder terms) are results of the development which we cite.
\end{enumerate}
\Cref{pf:volthm:sec:proof} assembles the proof of \cref{pf:volthm:main}, and
\cref{pf:volthm:sec:remarks} records some features of the formal statements.

% ---------------------------------------------------------------------------
\subsection{Conventions}\label{pf:volthm:sec:conv}

We use the coordinate framework for arithmetic varieties of \cref{def:covolume:hs}. Fix finite index sets
$K$ and $\iota_k$ ($k\in K$) and let $R_\iota=\Z[x_{k,i}]$ with its $\Z^K$-grading. An
\emph{arithmetic variety of dimension $d$} is a relevant multihomogeneous prime $P$ with $P\cap\Z=0$ and
$\dim R_\iota/P=d+|K|$; for a point $y$ of $\bigsqcup_k\iota_k$ in a field $\Omega$,
$P_y$ is the ideal of multihomogeneous relations of $y$, and $y$ is
\emph{birational} if $\Omega$ is generated by quotients of forms of equal multidegree evaluated at $y$. A
\emph{Hermitian line} $L$ is a multidegree $a\in\Z^K$ together with a positive function $h_L$ on the complex
multicone with $h_L(c\cdot z)=\prod_k|c_k|^{a_k}h_L(z)$; it is \emph{smooth}, \emph{conjugation invariant},
\emph{semipositive} on $P$ (the Levi form of $\log h_L$ is $\ge0$ on the tangent cones of the complex points
$P(\C)$ of $P$), or \emph{nef on the fibres} of $P$ (geometric degree $\ge0$ on all vertical curves in $P$)
in the evident sense. For integers $n_j$ and lines $\ell_j$, the line $\bigotimes_j\ell_j^{n_j}$ has multidegree $\sum n_ja_j$ and metric $\prod h_{\ell_j}^{n_j}$.

An \emph{intersection system} $\volthmII$ assigns a real number $\volthmII(P;L_1,\dots,L_m)$ to every
coordinate datum, ideal $P$ and $m$-tuple of Hermitian lines. It is an \emph{intersection family} if it
satisfies the conditions of \cref{def:covolume:hs}(c) (symmetry, additivity, the inductive formula of Bost,
Gillet and Soul\'e \cite{BostGilletSoule1994}), and it is \emph{full} if it is moreover multilinear in each
smooth entry, satisfies the projection formula for block maps of coordinates, and vanishes on
mixed-dimension configurations. We write:
\begin{itemize}
\item $\int_P\log h_{L_1}\cdots\log h_{L_e}$ for the integral over $P(\C)$ of the product of the first Chern
  forms of $L_1,\dots,L_e$ when $P$ has dimension $e+1$, and
  $\deg(P;L)=\int_P(\log h_L)^{e}$ for the geometric degree;
\item $\deg_\C(W)$ for the number of complex points (up to the torus action) of an arithmetic curve $W$, so that $\volthmII(W;L)/\deg_\C(W)$ is the normalised height of $W$ with respect to $L$;
\item $(a\cdot W)$ for the degree of the multidegree $a$ on a vertical curve $W$;
\item a set $F$ of arithmetic curves containing $P$ is \emph{Zariski dense} in $P$ if no multihomogeneous
  $f\notin P$ lies in all $W\in F$.
\end{itemize}
All results below quantify over all intersection systems; ``full'' is always assumed explicitly.

By definition, for a relevant ideal $W$ the degree $(a\cdot W)$ is $\sum_ka_k\,d_k(W)$, where $d_k(W)\ge0$ is the
degree of the divisor cut out on $W$ by a coordinate of block $k$ (the sum over its components $W'$ of the
multiplicity of $W'$ times the length of the dehomogenised coordinate ring of $W'$); for other $W$ it is $0$. Hence:

\begin{lemma}\label{pf:volthm:lem:curvedeg}
If $a_k\ge0$ for all $k$, then $(a\cdot W)\ge0$ for every $W$.
\formallabel{RootFieldNumerics.geoNum_one_nonneg'}
\end{lemma}

\begin{proof}
Every term $a_kd_k(W)$ is $\ge0$.
\end{proof}

% ---------------------------------------------------------------------------
\subsection{Chern-root polynomials and the flag degree}\label{pf:volthm:sec:poly}

Throughout, $\Psi$ is a finite set, $N\ge2$, $[N]=\{0,\dots,N-1\}$, and
\[
n=n(\Psi,N)=|\Psi|\,N(N-1)/2
\]
is the dimension of the flag variety $\mathrm{Fl}=\prod_{\tau\in\Psi}\mathrm{Fl}_N$. Let
$\volthmR=\R[x_{\tau,i}:(\tau,i)\in\Psi\times[N]]$ and $\volthmR^+=\volthmR[x_\infty]$, with the inclusion
$F\mapsto F^+$. A \emph{weight} is a vector $\lambda\in\R^{\Psi\times[N]}$, with
$\|\lambda\|=\max|\lambda_{\tau,i}|$ and $\|\lambda\|_1=\sum|\lambda_{\tau,i}|$; it is \emph{rational} if all
entries are rational. Put
\[
\volthmlin(\lambda)=-\sum_{\tau,i}\lambda_{\tau,i}x_{\tau,i},\qquad
\volthmlin^+(\lambda,s)=\volthmlin(\lambda)^++s\,x_\infty .
\]
For $\sigma=(\sigma_\tau)\in S_N^\Psi$ let $m_\sigma$ be the monomial exponent with
$m_\sigma(\tau,i)=N-1-\sigma_\tau(i)$; it has degree $n$. The \emph{flag degree} is the linear form
\begin{equation}\label{pf:volthm:eq:degfl}
\volthmdeg(F)=(-1)^n\sum_{\sigma\in S_N^\Psi}\Bigl(\prod_\tau\operatorname{sgn}\sigma_\tau\Bigr)\,
\mathrm{coeff}_{m_\sigma}(F),
\end{equation}
the integral over $\mathrm{Fl}$ when $x_{\tau,i}$ are read as Chern roots, and
$\volthmfd(\lambda)=\volthmdeg(\volthmlin(\lambda)^n)$. We also use:
\begin{itemize}
\item the power sums $p_k(\tau)=\sum_ix_{\tau,i}^k$, $p_k^{<j}(\tau)=\sum_{i<j}x_{\tau,i}^k$, and for
  $\sigma$ in $\Psi\sqcup\Psi$, written $\tau^+$ or $\tau^-$, the polynomials
  $\pi_{j,k}(\tau^+)=p^{<j}_k(\tau)$ and $\pi_{j,k}(\tau^-)=\sum_{i\ge j}(-x_{\tau,i})^k$;
\item the \emph{shift} and \emph{total} functionals
  $\volthmSh_j^\tau(F)=\volthmdeg\bigl(\sum_{i<j}\partial F/\partial x_{\tau,i}\bigr)$ and
  $\volthmTot^\tau(F)=\volthmdeg\bigl(\sum_{i}\partial F/\partial x_{\tau,i}\bigr)$;
\item $\volthmgap_j(\lambda,\tau)=\lambda_{\tau,j-1}-\lambda_{\tau,j}$ for $1\le j<N$; $\lambda$ is
  \emph{dominant} (resp.\ \emph{regular dominant}) if $\lambda_{\tau,b}\le\lambda_{\tau,a}$ (resp.\ $<$) for
  $a<b$; it has \emph{compact gaps $\ge c$} if $\lambda_{\tau,i}-\lambda_{\tau,i+1}\ge c$ whenever
  $i+1<N$ and $i+1\ne j$;
\item the simple reflection $s_{\tau,a,b}\lambda$ exchanging $\lambda_{\tau,a}$ and $\lambda_{\tau,b}$, the
  weight $\rho_{\tau,i}=N-1-i$, and the Vandermonde polynomial
  $\Delta=\prod_\tau\prod_{a<b}(x_{\tau,a}-x_{\tau,b})$;
\item the action $F\mapsto\gamma F$ of $\gamma\in\mathrm{Sym}(\Psi)$ by $x_{\tau,i}\mapsto x_{\gamma\tau,i}$.
\end{itemize}
For the lower bounds we need the following quantities, attached to $1\le j<N$, a weight $u$ and $g>0$:
\begin{itemize}
\item $\mathcal N_j(b,\mu)$: $4|\Psi|\,\|\mu\|\le(b-2)\volthmgap_j(\mu,\tau)$ for all $\tau$;
\item $\mathrm{WH}_j(u,g)$: $g>0$, $u$ is regular dominant, $\volthmgap_j(u,\tau)\ge g$ for all $\tau$, and
  $|u_{\tau,a}-u_{\tau,b}|\ge g$ whenever $a\ne b$ and $a$ lies in the smaller block (that is, $a<j$ if
  $j\le N-j$ and $a\ge j$ otherwise);
\item $B_0(u,g)=4|\Psi|N^2(1+\|u\|/g)$ and $\mathrm{spread}(u)=\sum_\tau(\max_iu_{\tau,i}-\min_iu_{\tau,i})$;
\item for reals $r,L$ (later $r=\mathrm{rd}_E$ and $L=L_H$), $\Lambda(B)=\Lambda(B;r,L)=\log(Br)+|L|$;
\item for a prime $p$, $\mathrm{fr}\colon\N\to\Psi\sqcup\Psi$, $k\ge1$ and $Q\in\volthmR$,
  \begin{align}
  \delta_j(p,\tau,\mathrm{fr},k,Q)&=\log p\cdot k\Bigl(\sum_{m\ge0}p^{-k(m+1)}
    \volthmdeg\bigl(\pi_{j,k-1}(\mathrm{fr}(m))Q\bigr)-\frac{\volthmdeg(p^{<j}_{k-1}(\tau)Q)}{p-1}\Bigr),
    \label{pf:volthm:eq:delta}\\
  r_j(p,\mathrm{fr},k,Q)&=\mathbf 1_{k\ge2}\,\log p\cdot k\sum_{m\ge0}p^{-k(m+1)}
    \volthmdeg\bigl(\pi_{j,k-1}(\mathrm{fr}(m))Q\bigr),\label{pf:volthm:eq:rest}
  \end{align}
  and, for a finite set $\mathrm{bad}$, a map $\mathrm{frob}\colon\N\times\Psi\times\N\to\Psi\sqcup\Psi$ and
  $B_0\in\N$, $R_j(B_0,\tau,k,Q)=\sum_{p>B_0,\ p\notin\mathrm{bad}}r_j(p,\mathrm{frob}(p,\tau,\cdot),k,Q)$,
  the sum over primes;
\item a \emph{power-sum decomposition} of a homogeneous $F$ of degree $n+1$ is a family $Q_{\tau,k}$
  ($1\le k\le N$) of homogeneous polynomials of degree $n+1-k$ with
  $F=\sum_\tau\sum_{k=1}^Np_k(\tau)Q_{\tau,k}$;
\item the test data: $d_X=|\Psi|\,j(N-j)$,
  $\mu_j=\sum_\tau\sum_{i\ge j}x_{\tau,i}$ if $N-j\le j$ and $\mu_j=-\sum_\tau p_1^{<j}(\tau)$ otherwise, $\phi_j=\prod_\tau\prod(x_{\tau,b}-x_{\tau,a})$ over $a<b$ on the same side of $j$, the \emph{test polynomial} $\mathrm{tp}_j=\mu_j^{d_X+1}\phi_j$ and
  $\mathrm{td}_j=\volthmdeg(\mu_j^{d_X}\phi_j)$.
\end{itemize}

\begin{remark}\label{pf:volthm:rem:sums}
The sums over $m$ and over $p$ in \eqref{pf:volthm:eq:delta}, \eqref{pf:volthm:eq:rest} and in $R_j$ are
defined formally as sums of possibly non-summable families, which are $0$ if the family is not summable. All
families occurring here are absolutely summable: the terms over $m$ are bounded by a constant times
$p^{-k(m+1)}$, and the terms over $p$ by a constant times $\log p/p^2$ (\cref{pf:volthm:lem:rest-summable}).
\end{remark}

\begin{lemma}[The flag degree]\label{pf:volthm:lem:degfl}
Let $1\le j\le N$ and $\tau,\sigma\in\Psi$.
\begin{enumerate}[label=(\alph*)]
\item $\volthmdeg$, $\volthmSh_j^\tau$ and $\volthmTot^\tau$ are linear, and $\volthmdeg(F)=0$ for $F$
  homogeneous of degree $\ne n$.
\item $\volthmdeg(p_k(\tau)Q)=0$ for every $k\ge1$ and $Q\in\volthmR$.
\item For $k\ge1$: $\volthmSh_j^\sigma(p_k(\tau)Q)=\mathbf 1_{\sigma=\tau}\,k\,\volthmdeg(p^{<j}_{k-1}(\tau)Q)$
  and $\volthmTot^\sigma(p_k(\tau)Q)=\mathbf 1_{\sigma=\tau,\,k=1}\,N\volthmdeg(Q)$.
\item For every weight $u$:
  $\volthmSh_j^\tau(\volthmlin(u)^{n+1})=-(n+1)\volthmfd(u)\sum_{i<j}u_{\tau,i}$ and
  $\volthmTot^\tau(\volthmlin(u)^{n+1})=-(n+1)\volthmfd(u)\sum_iu_{\tau,i}$.
\item $\volthmdeg$ changes sign under the exchange of two variables $x_{\tau,a},x_{\tau,b}$, $a\ne b$; for
  $\gamma\in\mathrm{Sym}(\Psi)$, $\volthmdeg(\gamma F)=\volthmdeg(F)$,
  $\volthmSh_j^{\gamma\tau}(\gamma F)=\volthmSh_j^\tau(F)$ and $\volthmTot^{\gamma\tau}(\gamma F)=\volthmTot^\tau(F)$.
\item There is $Q_1$ homogeneous of degree $n$ with $\volthmdeg(Q_1)=1$; if $1\le j<N$, there is $Q_2$
  homogeneous of degree $n-1$ with $\volthmdeg(p_1^{<j}(\tau)Q_2)=1$.
\end{enumerate}
\formallabel{RootFieldNumerics.degFl_Sh_Tot_linear_and_degFl_eq_zero_off_dim}
\formallabel{GlobalTheory.CompanionTheory.Ar.Top.degFl_Sh_Tot_linear}
\formallabel{GlobalTheory.T.flag_degree_linear_functionals}
\formallabel{Formalism.DegreeFormula.PPowSum.degFlPowerSum}
\formallabel{Formalism.DegreeFormula.degreePowerSumMulEqZero}
\formallabel{GlobalTheory.CompanionTheory.Ar.Top.Model.degFl_pk_mul}
\formallabel{RootFieldNumerics.shift_power_sum_mul}
\formallabel{GlobalTheory.CompanionTheory.Ar.Top.sh_tot_pk_mul}
\formallabel{GlobalTheory.T.flag_shift_total_power_sum}
\formallabel{GlobalTheory.CompanionTheory.Ar.Top.Sh_Tot_linRoot_pow}
\formallabel{GlobalTheory.T.flag_shift_total_lin_root_power}
\formallabel{GlobalTheory.CompanionTheory.Ar.Top.flag_integral_permutations}
\formallabel{GlobalTheory.T.flag_degree_symmetries}
\formallabel{RootFieldNumerics.exists_degree_one_class}
\formallabel{GlobalTheory.CompanionTheory.Ar.Top.testForms_sh_tot}
\formallabel{GlobalTheory.T.exists_flag_degree_one_polys}
\end{lemma}

\begin{proof}
(a) $\volthmdeg$ is a combination of coefficient functionals at the monomials $m_\sigma$, which have degree
$\sum_\tau\sum_i(N-1-\sigma_\tau(i))=n$; $\volthmSh$ and $\volthmTot$ are $\volthmdeg$ composed with
derivations.

(b) Write $\varepsilon(\sigma)=\prod_\tau\operatorname{sgn}\sigma_\tau$. Reindexing the sum over $i$ by
$i=\sigma_\tau^{-1}(v)$,
$\volthmdeg(p_k(\tau)Q)=\pm\sum_{v\in[N]}\sum_\sigma\varepsilon(\sigma)\,
\mathrm{coeff}_{m_\sigma}(x_{\tau,\sigma_\tau^{-1}(v)}^kQ)$, and each inner sum vanishes. Indeed
$\mathrm{coeff}_{m_\sigma}(x_{\tau,i}^kQ)$ with $i=\sigma_\tau^{-1}(v)$ is the coefficient of $Q$ at
$m_\sigma-k\,e_{\tau,i}$, which is $0$ if $v+k>N-1$; otherwise the exponent $N-1-v-k$ at $(\tau,i)$ equals the
exponent at $(\tau,i')$, $i'=\sigma_\tau^{-1}(v+k)$, and replacing $\sigma_\tau$ by its composite with the
transposition of the values $v,v+k$ gives a sign-reversing involution of the index set preserving
$m_\sigma-k\,e_{\tau,\sigma_\tau^{-1}(v)}$.

(c) By the Leibniz rule,
$\sum_{i<j}\partial_{\sigma,i}(p_k(\tau)Q)=\mathbf 1_{\sigma=\tau}\,k\,p^{<j}_{k-1}(\tau)Q+
p_k(\tau)\sum_{i<j}\partial_{\sigma,i}Q$, and the second term is killed by $\volthmdeg$ by (b). For
$\volthmTot$ the first term is $\mathbf 1_{\sigma=\tau}\,k\,p_{k-1}(\tau)Q$; for $k\ge2$ it is killed by (b),
and for $k=1$ it is $N\,Q$.

(d) $\partial_{\tau,i}\volthmlin(u)^{n+1}=-(n+1)u_{\tau,i}\volthmlin(u)^n$.

(e) The exchange of $x_{\tau,a},x_{\tau,b}$ maps $m_\sigma$ to $m_{\sigma'}$ with $\sigma'_\tau=\sigma_\tau\circ
(a\,b)$, and $\gamma$ maps $m_\sigma$ to $m_{\sigma''}$ with $\sigma''_{\gamma\tau}=\sigma_\tau$, preserving
$\varepsilon$; the statements on $\volthmSh$ and $\volthmTot$ follow since $\gamma$ commutes with the relevant
sums of derivations up to the relabelling $\tau\mapsto\gamma\tau$.

(f) $Q_1=(-1)^nx^{m_{\mathrm{id}}}$: among the $m_\sigma$ only $m_{\mathrm{id}}$ occurs in it. For $Q_2$ take
$(-1)^nx^{m}$ with $m=m_{\mathrm{id}}-e_{\tau,j-1}$. In $p_1^{<j}(\tau)x^m=\sum_{i<j}x^{m+e_{\tau,i}}$ the
term $i=j-1$ is $x^{m_{\mathrm{id}}}$, and for $i<j-1$ the exponent $m+e_{\tau,i}$ takes the value $N-1-j$ at
both $(\tau,j-1)$ and $(\tau,j)$, so it is not of the form $m_\sigma$.
\end{proof}

\begin{lemma}[Weyl formula]\label{pf:volthm:lem:weyl}
For every weight $\lambda$,
$\volthmfd(\lambda)=n!\prod_\tau\prod_{i<k}\frac{\lambda_{\tau,i}-\lambda_{\tau,k}}{k-i}$. In particular
$\volthmfd(\lambda)>0$ if $\lambda$ is regular dominant, $\volthmfd(\lambda)\ge0$ if $\lambda$ is dominant,
$\volthmfd(\rho)=n!$, and there is $c>0$ with $\volthmfd(\lambda)=c\,\Delta(\lambda)$ for all $\lambda$.
\formallabel{GlobalTheory.T.flag_degree_weight_formula}
\formallabel{GlobalTheory.CompanionTheory.Ar.Top.flagDeg_eq_delta}
\end{lemma}

\begin{proof}
The product formula and the two sign statements are
\lref{GlobalTheory.CompanionTheory.Ar.Top.flagDeg_eq_weyl}, which states exactly this (it is the
Vandermonde evaluation of \eqref{pf:volthm:eq:degfl}). The last two statements follow with
$c=n!/\prod_\tau\prod_{i<k}(k-i)$.
\end{proof}

\begin{lemma}[Polarisation]\label{pf:volthm:lem:polar}
Let $J$ be a finite set, $d\ge0$ and $\Phi$ a linear form on $\R[x_z:z\in J]$. The function
$c\mapsto\Phi\bigl((-\sum_zc_zx_z)^d\bigr)$ is a homogeneous polynomial $\mathrm{Pl}_\Phi$ of degree $d$ in
$c\in\R^J$, with $\mathrm{coeff}_m(\mathrm{Pl}_\Phi)=(-1)^d\binom{d}{m}\Phi(x^m)$ for $|m|=d$ (multinomial
coefficient). In particular $|\Phi(x^m)|\le|\mathrm{coeff}_m(\mathrm{Pl}_\Phi)|$.
\formallabel{GlobalTheory.CompanionTheory.Ar.Top.polarisation_functional}
\formallabel{GlobalTheory.T.polarization_polynomial_coeff_bound}
\end{lemma}

\begin{proof}
Expand by the multinomial theorem and apply $\Phi$ termwise; multinomial coefficients are $\ge1$.
\end{proof}

\begin{lemma}[Power-sum decomposition]\label{pf:volthm:lem:decomp}
Every homogeneous $F\in\volthmR$ of degree $n+1$ has a power-sum decomposition.
\formallabel{RootFieldNumerics.decomposition_power_sums_top_degree}
\formallabel{RootFieldNumerics.coinv_above_top}
\formallabel{GlobalTheory.CompanionTheory.Ar.Top.form_top_succ_eq_sum_pk}
\formallabel{GlobalTheory.T.homogeneous_power_sum_decomposition}
\end{lemma}

\begin{proof}
This is the statement that the coinvariant algebra of $S_N^\Psi$ vanishes in degrees $>n$. Let $I$ be the
ideal generated by the $p_k(\tau)$, $1\le k\le N$, $\tau\in\Psi$. By Newton's identities (over $\Q$) $I$ is
the ideal generated by the invariants of positive degree of $S_N^\Psi$, so by Chevalley's theorem
\cite{Chevalley1955} the graded algebra $\volthmR/I$ has Hilbert series
$\prod_\tau\prod_{d=1}^N(1+t+\dots+t^{d-1})$, a polynomial of degree $|\Psi|\sum_{d\le N}(d-1)=n$; concretely,
$\volthmR/I$ is spanned by the monomials $x^m$ with $m(\tau,i)\le N-1-i$. Hence every homogeneous $F$ of
degree $n+1$ lies in $I$, and taking homogeneous components of the coefficients gives a decomposition with
$Q_{\tau,k}$ homogeneous of degree $n+1-k$. The development proves the spanning statement directly, by an explicit reduction of monomials
modulo $I$.
\end{proof}

\begin{lemma}[Invariant linear forms]\label{pf:volthm:lem:invlin}
Let $W_j$ be the group of permutations of $\Psi\times[N]$ that preserve $\tau$ and the two sides $\{i<j\}$,
$\{i\ge j\}$. A linear form $l$ invariant under $W_j$ has the shape
$l(x)=\sum_\tau(a_\tau-b_\tau)\sum_{i<j}x_{\tau,i}+\sum_\tau b_\tau\sum_ix_{\tau,i}$ for some
$a,b\in\R^\Psi$.
\formallabel{GlobalTheory.CompanionTheory.Ar.Top.inv_linear_eval}
\end{lemma}

\begin{proof}
Invariance forces the coefficient of $x_{\tau,i}$ to depend only on $\tau$ and on the side of $i$; call
these $a_\tau$ ($i<j$) and $b_\tau$ ($i\ge j$).
\end{proof}

\begin{lemma}[Partial sums of decreasing sequences]\label{pf:volthm:lem:partial}
Let $0\le j\le N$, $N>0$, and $u_0\ge u_1\ge\dots\ge u_{N-1}$ real. Then
$0\le\sum_{i<j}u_i-\frac jN\sum_iu_i\le\sum_i|u_i|$ and $|\sum_iu_i|\le\sum_i|u_i|$.
\formallabel{RootFieldNumerics.partial_sum_ordered_bounds}
\formallabel{GlobalTheory.CompanionTheory.Ar.Top.partialSum_sub_mean_bounds}
\formallabel{GlobalTheory.T.dominant_partial_sum_bounds}
\end{lemma}

\begin{proof}
The average of the $j$ largest terms is at least the total average, which gives the first inequality; the
others are the triangle inequality, using $\sum_{i<j}u_i-\frac jN\sum_iu_i=\frac{N-j}{N}\sum_{i<j}u_i-
\frac jN\sum_{i\ge j}u_i$.
\end{proof}

% ---------------------------------------------------------------------------
\subsection{Root models and root towers}\label{pf:volthm:sec:models}

\begin{definition}[Root models]\label{pf:volthm:def:rootmodel}
A \emph{root model} $M$ of type $(\Psi,N)$ consists of coordinate data $(K,\iota)$; a field $\Omega$ and a
birational point $y$ such that $P_M=P_y$ is an arithmetic variety of dimension $n+1$; multidegrees
$a(r)\in\Z^K$ for $r\in\Psi\times[N]$; for every $t\in\R$ and $r$ a smooth, conjugation-invariant
Hermitian line $x_t(r)$ of multidegree $a(r)$; an integer $k_A\ge1$; a smooth,
conjugation-invariant Hermitian line $A$; and an integer $\delta\ge1$. Put $\ell_t(r)=x_t(r)$ and $\ell_t(\infty)=A$, and for
a weight $\lambda$ and $s\in\R$ let $\volthmcoef(\lambda,s)(r)=-\lambda_r$, $\volthmcoef(\lambda,s)(\infty)=s/k_A$.
\begin{enumerate}[label=(\alph*)]
\item For an exponent $m$ on $(\Psi\times[N])\sqcup\{\infty\}$ of degree $n+1$ let
  $\mathrm{ar}_t(x^m)=\volthmII(P_M;\ell_t^{m})/(\delta k_A^{m(\infty)})$, where $\ell_t^m$ lists each
  $\ell_t(r)$ with multiplicity $m(r)$, and $\mathrm{ar}_t(x^m)=0$ in other degrees; $\mathrm{ar}_t$ is the
  linear extension to $\volthmR^+$. $M$ \emph{converges} if $\lim_{t\to0^+}\mathrm{ar}_t(x^m)$
  exists for every $m$, and $\volthmPP_M$ is the linear form on $\volthmR^+$ whose value at $x^m$ is this limit.
\item For an exponent $m$ of degree $n$ let $\mathrm{geo}_M(x^m)$ be $1/(\delta k_A^{m(\infty)})$ times the integral
  over $P_M$ of the product of the $\log h_{\ell_1(r)}$, each $r$ with multiplicity $m(r)$; let
  $\mathrm{geo}_M(x^m)=0$ in other degrees, and extend linearly.
\item For an arithmetic curve $W\supseteq P_M$ (a \emph{point} of $M$) let
  \[
  \mathrm{ht}_t(\lambda,s;W)=\frac1{\deg_\C(W)}\sum_r\volthmcoef(\lambda,s)(r)\,\volthmII(W;\ell_t(r)).
  \]
\item $\mathrm{Nef}_p(\lambda,s)$ means that $\sum_r\volthmcoef(\lambda,s)(r)\,(\deg\ell_1(r)\cdot W)\ge0$
  for every vertical curve $W\supseteq P_M$ of characteristic $p$, and $\mathrm{Nef}_{\mathrm f}(\lambda,s)$ means
  this for all primes $p$. $\mathrm{Nef}_{\mathrm g}(\lambda,s)$ means that
  $\sum_r\volthmcoef(\lambda,s)(r)\deg(W;\ell_1(r))\ge0$ for every arithmetic variety $W\supseteq P_M$ of
  dimension $2$.
\item $\mathrm{SP}(t,\lambda,s)$: the Levi form of $\sum_r\volthmcoef(\lambda,s)(r)\log h_{\ell_t(r)}$
  is $\ge0$ on the tangent cone at every complex point of $P_M$.
\item $\mathrm{Eff}(\nu)$: there are $q\ge1$, integers $n_r=q\,\volthmcoef(\nu,-1)(r)$ and a nonzero
  integral section (an element of $\Omega$ integral on every standard chart after dehomogenisation)
  of the line $\bigotimes_r\ell_1(r)^{n_r}$.
\end{enumerate}
\formal{GlobalTheory.CompanionTheory.Ar.Top.RootModel}
\formallabel{GlobalTheory.CompanionTheory.Ar.Top.RootModel.PP}
\end{definition}

\begin{remark}\label{pf:volthm:rem:limunder}
Formally $\volthmPP_M(x^m)$ is defined by a choice of limit along $t\to0^+$, which equals the limit whenever
the limit exists; all statements below about $\volthmPP$ are used under hypotheses (convergence of $M$, or
the conditions $\mathrm{I}_2$ and $\mathrm{Cv}$ below) that guarantee its existence.
\end{remark}

\begin{definition}[Root towers]\label{pf:volthm:def:tower}
A \emph{root tower} $T$ of type $(\Psi,N)$ consists of an integer $j$ with $1\le j<N$, a finite nonempty set
$\mathrm{Comp}$ (``components''), a finite set $\mathrm{bad}\subset\N$ containing $2$, a rational weight
$\nu$, and root models $M_{B,c}$ of type $(\Psi,N)$ for $B\in\N$ and $c\in\mathrm{Comp}$. We use the
derived quantities
($\lambda_s=\lambda-s\nu$, $\mathrm{scale}(s)=e^{-1/s^2}$):
\begin{align*}
&\vol(c)=\mathrm{geo}_{M_{2,c}}\bigl((\volthmlin\rho)^{n,+}\bigr)/\volthmfd(\rho),\qquad
  \vol_T=\textstyle\sum_c\vol(c),\\
&\volthmPP_{B,c}(F)=\volthmPP_{M_{B,c}}(F^+),\qquad \volthmPP_B=\textstyle\sum_c\volthmPP_{B,c},\qquad
  \Pi(B,\lambda)=\volthmPP_B(\volthmlin(\lambda)^{n+1}),\\
&\Pi^+(B,c,\lambda,s)=\volthmPP_{M_{B,c}}(\volthmlin^+(\lambda,s)^{n+1}),\qquad
  \mathrm{ar}(B,c,\lambda,s)=\mathrm{ar}_{\mathrm{scale}(s)}\bigl(\volthmlin^+(\lambda_s,s)^{n+1}\bigr),\\
&\mathrm{dg}(B,c,\lambda,s)=\mathrm{geo}_{M_{B,c}}\bigl(\volthmlin^+(\lambda_s,s)^n\bigr),\qquad
  \mathrm{ht}(B,c,\lambda,s;W)=\mathrm{ht}_{\mathrm{scale}(s)}(\lambda_s,s;W),\\
&\mathrm{slope}(B,c,u)=\frac{\volthmPP_{B,c}(\volthmlin(u)^{n+1})}{(n+1)\vol(c)\volthmfd(u)},
\end{align*}
the quantities $\mathrm{ar}$, $\mathrm{dg}$, $\mathrm{ht}$ being computed in $M_{B,c}$. We say $p$ is
\emph{in $S_B$} if $p$ is prime and $p\le B$ or $p\in\mathrm{bad}$.
\formal{GlobalTheory.CompanionTheory.Ar.Top.RootTower}
\formallabel{GlobalTheory.CompanionTheory.Ar.Top.RootTower.vol}
\formallabel{GlobalTheory.CompanionTheory.Ar.Top.RootTower.slope}
\formallabel{GlobalTheory.CompanionTheory.Ar.Top.RootTower.lamS}
\formallabel{GlobalTheory.CompanionTheory.Ar.Top.RootTower.volT}
\formallabel{GlobalTheory.CompanionTheory.Ar.Top.RootTower.PPc}
\formallabel{GlobalTheory.CompanionTheory.Ar.Top.RootTower.PPt}
\formallabel{GlobalTheory.CompanionTheory.Ar.Top.RootTower.PB}
\formallabel{GlobalTheory.CompanionTheory.Ar.Top.RootTower.mDeg}
\formallabel{GlobalTheory.CompanionTheory.Ar.Top.RootTower.mbarHt}
\end{definition}

\begin{definition}[The flag height functional]\label{pf:volthm:def:hh}
Let $P_{\mathrm{Fl}}$ be the multicone ideal of the generic point of $\mathrm{Fl}$ in Pl\"ucker coordinates:
the blocks are indexed by $(\tau,q)\in\Psi\times[N-1]$, the coordinates of block $(\tau,q)$ by increasing
$(q+1)$-tuples in $[N]$, and the point is given by the $(q+1)$-minors of the first $q+1$ rows of a generic
$N\times N$ matrix. For $(\tau,i)\in\Psi\times[N]$ let $\xi(\tau,i)$ be the Hermitian line of
multidegree $e_{\tau,i}-e_{\tau,i-1}$ (with $e_{\tau,-1}=e_{\tau,N-1}=0$) and Fubini--Study metric
$\prod_{(\tau,q)}\|z_{\tau,q}\|^{a_{\tau,q}}$. For an intersection system
$\volthmII$ the \emph{flag height functional} $\volthmhh=\volthmhh^{\volthmII}$ is the linear form on $\volthmR$
with $\volthmhh(x^m)=(-1)^{n+1}\volthmII(P_{\mathrm{Fl}};\xi^m)$ if $\deg m=n+1$ and $\volthmhh(x^m)=0$ otherwise.
\formal{GlobalTheory.CompanionTheory.Ar.Top.hhFl}
\end{definition}

% ---------------------------------------------------------------------------
\subsection{The statements}\label{pf:volthm:sec:statements}

\begin{definition}[Conclusions of the volume theorem]\label{pf:volthm:def:conclusions}
Let $T$ be a root tower of type $(\Psi,N)$, $\volthmII$ an intersection system, $h$ a linear form on
$\volthmR$, and $C,r,L\in\R$; write $\Lambda(B)=\log(Br)+|L|$ and
$\mathrm{br}(B,u)=\Lambda(B)\|u\|_1+\min(j,N-j)\,\mathrm{spread}(u)$.
\begin{enumerate}[label=(\alph*)]
\item \emph{Proportionality} $\mathrm{Pr}$: $\vol(c)>0$ for all $c$; $\mathrm{geo}_{M_{B,c}}(F^+)=\vol(c)\volthmdeg(F)$ for
  all $B,c,F$; and $\mathrm{geo}_{M_{B,c}}=\mathrm{geo}_{M_{B',c}}$ for all $B,B',c$.
\item \emph{Equal volumes} $\mathrm{EV}$: $\vol(c)=\vol(c')$ for all $c,c'$.
\item \emph{Limit formulas} $\mathrm{T}_1$: for all $u,g$ with $\mathrm{WH}_j(u,g)$, all $B$ with $B_0(u,g)<B+1$ and all $c$: for
  all sufficiently small $s>0$, $M_{B,c}$ satisfies $\mathrm{Nef}_{\mathrm f}(u_s,s)$ and
  $\mathrm{SP}(\mathrm{scale}(s),u_s,s)$; as $s\to0^+$,
  $\mathrm{ar}(B,c,u,s)\to\volthmPP_{B,c}(\volthmlin(u)^{n+1})$ and $\mathrm{dg}(B,c,u,s)\to\vol(c)\volthmfd(u)$;
  $\vol(c)\volthmfd(u)>0$; and $\mathrm{ar}(B,c,u,s)/((n+1)\mathrm{dg}(B,c,u,s))\to\mathrm{slope}(B,c,u)$.
\item \emph{Asymptotic formula} $\mathrm{AF}(h)$: there are $C'$ and $B_1$ such that for every $B\ge B_1$ there are
  $\alpha,\beta\in\R^\Psi$ with
  \[
  \Bigl|\volthmPP_B(x^m)-\vol_T\Bigl(h(x^m)-\sum_\tau(\alpha_\tau-\beta_\tau)\volthmSh_j^\tau(x^m)-
  \sum_\tau\beta_\tau\volthmTot^\tau(x^m)\Bigr)\Bigr|\le C'\frac{\log B}{B}
  \]
  for all monomials $x^m$ of degree $n+1$. The \emph{averaged asymptotic formula} $\overline{\mathrm{AF}}(h)$ is the
  same statement with $\alpha,\beta\in\R$
  scalars, i.e.\ with $\sum_\tau(\alpha-\beta)\volthmSh_j^\tau+\sum_\tau\beta\volthmTot^\tau$.
\item \emph{Total lower bound} $\mathrm{T}_2(C,r,L)$: for all $u,g$ with $\mathrm{WH}_j(u,g)$ and all $B\ge128$ with $B_0(u,g)<B+1$
  and $|\mathrm{bad}|\le B$:
  $\Pi(B,u)\ge-C(n+1)\vol_T\volthmfd(u)\,\mathrm{br}(B,u)$.
\item \emph{Mean slope bound} $\mathrm{MS}(C,r,L)$: under the same hypotheses,
  $\frac1{|\mathrm{Comp}|}\sum_c\mathrm{slope}(B,c,u)\ge-C\,\mathrm{br}(B,u)$.
\end{enumerate}
\formallabel{GlobalTheory.CompanionTheory.Ar.Top.RootTower.Proportional}
\formallabel{GlobalTheory.CompanionTheory.Ar.Top.RootTower.ThmB}
\end{definition}

The data of the volume theorem are those of \cref{def:covolume:tower}: a CM field $E$ which is abelian over
$\Q$; a \emph{host} $H$ (a prime $d\equiv7\bmod16$, a square root of $-d$ in $\OO_E$ and a compatible
element; $\Psi=\Psi_H$, the \emph{CM type} of $H$, is the set of embeddings $E\to\C$ sending this root to a fixed root in
$\C$); a hermitian \emph{datum} $D$ of signature $(p,q)$ over $\OO_E$; moduli spaces $X^o_l$ for all levels
$l$; a level $l$, a coefficient datum $c$ and a model $Y$ of $X^o_l$ over $\OO_{c.F}$; a root tower $T$ of
type $(\Psi,p+q)$ and a \emph{good origin} $o$ of $T$ on $Y$, which realises the root models of $T$ on flag
bundles over the components of $Y$ with smooth and flat analytic realisations and ample reference models. The
model $Y$ is \emph{global} if it has integral models at all primes, with good reduction and compatible with
$Y$; the good origin satisfies the \emph{extra conditions} $(\mathrm{X})$ if the point maps of its analytic
systems are local holomorphic immersions and all its Hodge components have rank $4$ at the primes of $S_o(B)$,
the set of primes attached to $o$ at level $B$. For a good origin one has $j=p$, $S_o(B)$ is the set of primes
in $S_B$, and $\mathrm{bad}$ is the set of primes dividing $2\,\mathrm{disc}(E)\,M'_l$, $M'_l$ the modulus of the
level. Put
\[
\mathrm{rd}_E=|\mathrm{disc}\,E|^{1/[E:\Q]},\qquad
L_H=\Re\frac{L'(1,\chi_d)}{L(1,\chi_d)},
\]
with $\chi_d$ the quadratic character modulo $d$.

\begin{definition}[$\sfname{VolumeTheoremX}$]\label{pf:volthm:def:VolX}
$\sfname{VolumeTheoremX}$ is the statement: there is $C\ge0$ such that for all data as above and every
intersection system $\volthmII$, if $Y$ is global, $o$ satisfies $(\mathrm{X})$, $\volthmII$ is full, $q\ge1$,
$3q+1\le p+q$ and $|\mathrm{bad}|\le4|\Psi|(p+q)^2$, then $T$ satisfies $\mathrm{Pr}$,
$\mathrm{EV}$, $\mathrm{T}_1$, $\mathrm{AF}(\volthmhh^{\volthmII})$, $\mathrm{T}_2(C,\mathrm{rd}_E,L_H)$ and
$\mathrm{MS}(C,\mathrm{rd}_E,L_H)$ (with respect to $\volthmII$).
\formal{GlobalTheory.CompanionTheory.Ar.Top.VolumeTheoremX}
\formal{GlobalTheory.CompanionTheory.Ar.Top.GoodOrigin}
\formal{GlobalTheory.CompanionTheory.Ar.Top.GoodOrigin.Extra}
\formal{GlobalTheory.CompanionTheory.Ar.Top.IsGlobal}
\end{definition}

\begin{definition}[Inputs of the volume theorem]\label{pf:volthm:def:Hinputs}
The following are propositions of the development; they are hypotheses of \cref{pf:volthm:main}.
\begin{enumerate}[label=(\roman*)]
\item[$(\mathrm{H}_0)$] The \emph{general inputs}: the conjunction of fourteen statements, namely the moduli
  inputs (existence of moduli spaces and uniformisations, Siegel families and integral models at all primes,
  compatibility of toroidal compactifications with change of level and with Hecke operators); a
  transformation formula for theta constants under change of theta section and Lie basis; extension of
  homomorphisms of semi-abelian schemes over normal bases; a numerical statement on Segre embeddings of chart
  bundles; resolution of singularities in projective coordinates; strong approximation for special unitary
  groups; the Hodge--Tate statement for $p$-divisible groups; GAGA for sections of line bundles; a bound for
  the maximal slope of Lie algebras of principal semi-abelian schemes; the \emph{analytic and
  intersection-theoretic inputs} $(\mathrm{AI})$ (convergence and regularisation of plurisubharmonic weights,
  pull-backs, existence of resolutions, the degree formula, nonnegativity of arithmetic intersection numbers,
  the recursion for rational sections, locality in the metrics, the projection formula, the line and
  fibre-dimension statements, and the Kleiman inequality over fields, denoted $(\mathrm{K})$ below); the
  \emph{Hilbert--Samuel package} (which contains the inputs $(\mathrm{Z}^\pm)$ of \cref{pf:volthm:def:inputs}); an
  analytic package; a \emph{height package} (which contains $(\mathrm{GH})$ of \cref{pf:volthm:def:inputs} and a
  relative form $(\mathrm{Z}^-_{\mathrm{rel}})$ of $(\mathrm{Z}^-)$); and $(\mathrm{AP})$ of
  \cref{pf:volthm:def:inputs}.
\item[$(\mathrm{H}_1)$] Integrability of the fibre sums of sheet densities of good linear projections, which makes
  the integrals $\int_P$ well defined.
\item[$(\mathrm{H}_2)$] Existence, uniqueness, bimultiplicativity and compatibility of the Weil pairing on torsion
  points of polarised abelian schemes.
\item[$(\mathrm{H}_3)_{C_h,C_u}$] For all admitted component data with a cusp weight $U$ and a tied
  regularisation, the regularised Hodge lines are semipositive, nef on the fibres and ample on the generic
  fibre, the heights of their points are bounded below by $-C_hfj'\Lambda(B)-\frac pqK_A$, and the mean of $U$ is
  at most $C_ufj'\Lambda(B)$ (here $f=|\Psi|$ and $j'=N-j$).
\item[$(\mathrm{H}_4)$] At all places over a prime $p\notin S_0$ which is odd, unramified in $E$ and Hodge--Tate
  free, the Chern numbers on the special fibre of the root lines of a flag with two levels are given by an
  explicit formula in terms of Frobenius labels.
\item[$(\mathrm{H}_5)$] At all places over a Hodge--Tate free prime the root lines restrict to a trivial flag on
  every component of the special fibre.
\item[$(\mathrm{H}_6)$] The arithmetic degree of a flat generator of the line bundle of a flag is
  $\frac{p-q}2\Re\log\Omega_E+q\,(2\log\sqrt d-\log\mathrm{rd}_E)$, where
  $\log\Omega_E=\gamma+\log(2\pi/\mathrm{rd}_E)-\frac{L'}{L}(1,\chi)$ for the Kronecker character $\chi$.
\item[$(\mathrm{H}_7)$] Twisting the universal object by automorphisms of the host is realised by automorphisms of
  the moduli space.
\item[$(\mathrm{H}_8)$] Let $M,M'$ be root models realised on $Y$ by two Hodge ties with the same set of primes
  $S$, and let $\Phi$ be an automorphism of the moduli space which twists the action of $\OO_E$ by a host
  automorphism, maps the compactification of the first realisation to that of the second and preserves the data
  at the primes of $S$. Then the normalised intersection numbers of the Chern-root lines of $M$ and of $M'$ (with
  labels permuted accordingly) have the same limit as $t\to0^+$.
\item[$(\mathrm{H}_9)$] Two realisations with the same singularities of block projections of one arithmetic
  variety have normalised intersection numbers (weighted by the degrees of the function field extensions) with
  the same limit as $t\to0^+$.
\end{enumerate}
\formal{GlobalTheory.CompanionTheory.Ar.Top.Given}
\formal{GlobalTheory.CompanionTheory.Ar.PB.Inputs}
\formal{ReciprocityTheory.ChartRestriction.G.Geo.PtsIntegrableStatement}
\formal{ReciprocityTheory.ChartRestriction.G.Out.WeilSt}
\formal{GlobalTheory.CompanionTheory.Ar.Cusp.NineGeoU}
\formal{GlobalTheory.CompanionTheory.Ar.HodgeLatticeFlagBound}
\formal{GlobalTheory.CompanionTheory.Ar.TrivFlagOnHodge}
\formal{GlobalTheory.CompanionTheory.Ar.Fl.FlatDegree}
\formal{GlobalTheory.CompanionTheory.Ar.TwistAut}
\formal{GlobalTheory.CompanionTheory.Ar.Cmp.CmpPair}
\formal{GlobalTheory.CompanionTheory.Ar.Cmp.CmpSt}
\end{definition}

These hypotheses are discharged where the theorem is applied:
\lref{GlobalTheory.CompanionTheory.Ar.Top.goodTower_and_volume} proves the existence of good towers together
with $\sfname{VolumeTheoremX}$ from the inputs of \cref{thm:tower}.

\begin{theorem}[Volume theorem]\label{pf:volthm:main}
$(\mathrm{H}_0)$, $(\mathrm{H}_1)$, $(\mathrm{H}_2)$, $(\mathrm{H}_3)_{C_h,C_u}$ for some reals $C_h,C_u$, $(\mathrm{H}_4)$, $(\mathrm{H}_5)$,
$(\mathrm{H}_6)$, $(\mathrm{H}_7)$, $(\mathrm{H}_8)$ and $(\mathrm{H}_9)$ imply
$\sfname{VolumeTheoremX}$. One can take $C=2\max(C_h+C_u/2,0)+17$.
\formallabel{GlobalTheory.T.volume_theorem_of_inputs}
\formallabel{GlobalTheory.CompanionTheory.Ar.Cusp.volumeTheoremX_of_inputs}
\end{theorem}

The proof is given in \cref{pf:volthm:sec:proof}. Its core is the following abstract statement.

\begin{definition}[Tower conditions]\label{pf:volthm:def:conditions}
Let $T$ be a root tower of type $(\Psi,N)$ with index $j$, $\volthmII$ an intersection system, $h$ a linear form
on $\volthmR$, $\mathrm{frob}\colon\N\times\Psi\times\N\to\Psi\sqcup\Psi$, and $C,r,L\in\R$. The
quantifiers ``for all $B\ge2$'' etc.\ range over integers, and $M=M_{B,c}$ in each statement.
\begin{enumerate}[label=(\arabic*)]
\item \emph{Convergence} $\mathrm{Cv}$: every $M_{B,c}$ converges.
\item \emph{Coefficient bound} $\mathrm{I}_1$: there is $C_2$ with $|\volthmPP_M(x^m)|\le C_2\log B$ for all $B\ge2$, $c$ and monomials
  $x^m\in\volthmR^+$.
\item \emph{Regularisation rate} $\mathrm{I}_2$: for every $a$ there are $s_0>0$ and $C_3$ such that
  $|\mathrm{ar}(B,c,\lambda,s)-\Pi^+(B,c,\lambda_s,s)|\le C_3s$ whenever $B\ge2$, $\|\lambda\|_1\le a$,
  $0<s\le s_0$ and $\lambda_s$ is dominant.
\item \emph{Semipositivity} $\mathrm{A}_1$: for every $a$ there is $s_0>0$ such that $M$ satisfies
  $\mathrm{SP}(\mathrm{scale}(s),\lambda_s,s)$ under the same hypotheses.
\item \emph{Nefness at small primes} $\mathrm{M}_1$: $M$ satisfies $\mathrm{Nef}_p(\mu,0)$ whenever $p$ is in $S_B$
  and $\mu$ is dominant.
\item \emph{Nefness at large primes} $\mathrm{M}_2$: $M$ satisfies $\mathrm{Nef}_p(\mu,0)$ whenever $p$ is a prime not in $S_B$, $\mu$ is
  dominant and $\mathcal N_j(p,\mu)$ holds.
\item \emph{Nefness of the reference line} $\mathrm{M}_0$: $M$ satisfies $\mathrm{Nef}_{\mathrm f}(0,1)$.
\item \emph{Positivity of shifted lines} $\mathrm{I}_3$: for every $a$ there is $s_0>0$ such that $M$ satisfies $\mathrm{Nef}_{\mathrm f}(\lambda_s,s)$ and
  $\mathrm{SP}(\mathrm{scale}(s),\lambda_s,s)$ whenever $B\ge2$, $\|\lambda\|_1\le a$, $0<s\le s_0$,
  $\lambda_s$ is dominant and $\mathcal N_j(B,\lambda_s)$ holds.
\item \emph{Lower height bound} $\mathrm{I}_5$: for every $c_0>0$ there are $\theta_0>0$ and $C$ such that for all $B\ge2$, $c$, rational
  $\lambda$ with $\|\lambda\|\le1$, $0<\theta\le\theta_0$ and rational $s=\theta^2$, if
  $\volthmgap_j(\lambda,\tau)\ge\theta$ for all $\tau$ and $\lambda$ has compact gaps $\ge c_0$, then
  $\mathrm{ht}(B,c,\lambda,s;W)\ge-C\log(B/s)$ for every point $W$ of $M$.
\item \emph{Upper height bound} $\mathrm{I}_6$: under the same quantifiers and hypotheses, there is a Zariski dense set of points $W$ of $M$
  with $\mathrm{ht}(B,c,\lambda,s;W)\le C\log B$.
\item \emph{Generic nefness and effectivity} $\mathrm{I}_7$: $M$ satisfies $\mathrm{Nef}_{\mathrm g}(\lambda,0)$ for all dominant $\lambda$,
  $\mathrm{Nef}_{\mathrm g}(0,1)$, and $\mathrm{Eff}(\nu)$.
\item \emph{Antisymmetry} $\mathrm{I}_8(h)$: for $B\ge2$, $c$, $\tau$ and $b=a+1\ne j$, and every weight $\lambda$,
  \begin{equation}\label{pf:volthm:eq:I8}
  \volthmPP_{B,c}(\volthmlin(s_{\tau,a,b}\lambda)^{n+1})-\vol(c)\,h(\volthmlin(s_{\tau,a,b}\lambda)^{n+1})=
  -\bigl(\volthmPP_{B,c}(\volthmlin(\lambda)^{n+1})-\vol(c)\,h(\volthmlin(\lambda)^{n+1})\bigr).
  \end{equation}
\item \emph{Flag height conditions} $\mathrm{I}_9(h)$ (on $h$ alone): $h(\volthmlin(\lambda)^{n+1})=0$ for dominant $\lambda$
  with $\lambda_{\tau,a}=\lambda_{\tau,a+1}$ for some $\tau,a$; $h(\volthmlin(\lambda)^{n+1})\ge0$ for dominant
  $\lambda$; and $h(p_1(\tau)Q)=0$ for $Q$ homogeneous of degree $n$.
\item \emph{Euler product identity} $\mathrm{EP}(\mathrm{frob})$: for $2\le B_0\le B$, $c$, $\tau$, $1\le k\le N$ and $Q$ homogeneous of
  degree $n+1-k$,
  \begin{equation}\label{pf:volthm:eq:prop6}
  \volthmPP_{B_0,c}(p_k(\tau)Q)=\volthmPP_{B,c}(p_k(\tau)Q)+\vol(c)\sum_{\substack{B_0<p\le B\\ p\notin\mathrm{bad}}}
  \delta_j\bigl(p,\tau,\mathrm{frob}(p,\tau,\cdot),k,Q\bigr),
  \end{equation}
  the sum over primes.
\item \emph{Galois invariance} $\mathrm{GI}$: there is a subgroup $\Gamma\subseteq\mathrm{Sym}(\Psi)$ acting transitively on $\Psi$
  with $\volthmPP_B(\gamma F)=\volthmPP_B(F)$ for $\gamma\in\Gamma$, $B\ge2$ and $F\in\volthmR$.
\item \emph{Flat sections} $\mathrm{FL}(C,r,L)$: there are reals $h(B,c,\tau)$ with
  $\volthmPP_{B,c}(p_1(\tau)Q)=h(B,c,\tau)\vol(c)\volthmdeg(Q)$ for $B\ge2$ and $Q$ homogeneous of degree $n$,
  and $|h(B,c,\tau)|\le CN\Lambda(B)$ for $B\ge128$ with $|\mathrm{bad}|\le B$.
\item \emph{Slope condition} $\mathrm{SL}(C,r,L)$: for $B\ge128$ with $4|\Psi|N^2\le B$ and $|\mathrm{bad}|\le B$,
  \[
  \volthmPP_B(\mathrm{tp}_j)\ge-C|\Psi|\min(j,N-j)\Lambda(B)\,(d_X+1)\vol_T\mathrm{td}_j.
  \]
\item \emph{Constant shape} $\mathrm{CS}(\mathrm{frob})$ (on $\mathrm{frob}$): whether $\mathrm{frob}(p,\tau,m)$ lies in the first copy of $\Psi$
  does not depend on $\tau$.
\item \emph{Label invariance} $\mathrm{LI}(h)$ (on $h$): $h(\gamma F)=h(F)$ for all $\gamma\in\mathrm{Sym}(\Psi)$.
\end{enumerate}
\formallabel{GlobalTheory.CompanionTheory.Ar.Top.RootTower.I1}
\end{definition}

\begin{definition}[Inputs of the tower theorem]\label{pf:volthm:def:inputs}
\begin{enumerate}[label=(\alph*)]
\item[$(\mathrm{Z}^-)$] For an intersection family $\volthmII$, an arithmetic variety $P$ of
  dimension $D+1$ and a Hermitian line $A$ which is semipositive on $P$, nef on the fibres and conjugation
  invariant: if $e\deg_\C(W)\le\volthmII(W;A)$ for all arithmetic curves $W\supseteq P$, then
  $(D+1)e\deg(P;A)\le\volthmII(P;A^{D+1})$.
\item[$(\mathrm{Z}^+)$] Under the same assumptions, if $F$ is a Zariski dense set of
  arithmetic curves $W\supseteq P$ with $\volthmII(W;A)\le e'\deg_\C(W)$, then
  $\volthmII(P;A^{D+1})\le(D+1)e'\deg(P;A)$.
\item[$(\mathrm{FH})$] Given $(\mathrm{AI})$, for every full $\volthmII$ the polynomial
  \[
  H_{\mathrm{Fl}}(\lambda)=\sum_w\prod_k\lambda_{w(k)}\,\volthmII\bigl(P_{\mathrm{Fl}};\xi(w(1)),\dots,\xi(w(n+1))\bigr)
  \]
  (sum over words $w$ of length $n+1$ in $\Psi\times[N]$) is homogeneous of degree $n+1$, vanishes when a gap
  $\lambda_{\tau,a}-\lambda_{\tau,a+1}$ vanishes, and is $\ge0$ on dominant weights.
\item[$(\mathrm{GH})$] For full $\volthmII$, $1\le j\le N-1$ and $e\in\R$: if $\frac12\log\det(1+ZZ^*)\le e$ for
  every complex $(N-j)\times j$ matrix $Z$ with entries of absolute value $1$, and the relevant flag integral is
  positive, then a lower bound by $-(d_X+1)\cdot(\text{flag integral})\cdot|\Psi|\,e$ holds for the arithmetic
  height of the corresponding test class of $\mathrm{Fl}$; by the identities in the proof of
  \cref{pf:volthm:prop:hh}(d) it amounts to the bound $\volthmhh(\mathrm{tp}_j)\le(d_X+1)\,\mathrm{td}_j\,|\Psi|\,e$.
\item[$(\mathrm{AP})$] Let $1\le j<N$, $h$ a polynomial on $\R^{\Psi\times[N]}$ and, for each $\tau$, $K_\tau$ a
  compact set in the hyperplane $H_\tau=\{x_{\tau,j-1}=x_{\tau,j}\}$ containing a nonempty relatively open subset of
  $H_\tau$, with $h=0$ on $K_\tau$. Then there is $C_0$ such that for every polynomial $P$ and all $v,\varepsilon$:
  if $P-vh$ is homogeneous of degree $n+1$ and changes sign under every simple reflection
  $x_{\tau,a}\leftrightarrow x_{\tau,a+1}$ with $a+1\ne j$, and $|P|\le\varepsilon$ on every $K_\tau$, then there is
  a linear form $l$ invariant under the group $W_j$ of \cref{pf:volthm:lem:invlin} with
  $|\mathrm{coeff}_m(P-vh-\Delta l)|\le C_0\varepsilon$ for all $m$.
\end{enumerate}
\formal{ReciprocityTheory.ChartRestriction.G.Geo.SuccessiveMinimaLowerBound}
\formal{ReciprocityTheory.ChartRestriction.G.Geo.SuccessiveMinimaUpperBound}
\formallabel{GlobalTheory.CompanionTheory.Ar.PB.FlagHeight}
\formal{GroundTheory.Heights₂.Heights.LowerBoundFree.GrassmannianHeightZTwo.Statement}
\formal{GroundTheory.Heights₂.Heights.AntisymPerturbation.Statement}
\end{definition}

The inputs $(\mathrm{Z}^-)$ and $(\mathrm{Z}^+)$ are the two halves of Zhang's theorem of successive minima
\cite[Thm.~5.2]{Zhang1995}, in the coordinate formulation of the development.

\begin{theorem}[Tower theorem]\label{pf:volthm:thm:tower}
Assume $(\mathrm{FH})$, $(\mathrm{AI})$, $(\mathrm{Z}^-)$,
$(\mathrm{Z}^+)$, $(\mathrm{AP})$ and $(\mathrm{GH})$. Let
$\volthmII$ be a full intersection system, $\Psi$ nonempty, $T$ a root tower of type $(\Psi,N)$,
$\mathrm{frob}$ a map as above, and $C,C'\ge0$, $r\ge1$, $L$ reals. If $\mathrm{CS}(\mathrm{frob})$ and
$T$ satisfies $\mathrm{Pr}$, $\mathrm{I}_1$, $\mathrm{I}_2$, $\mathrm{A}_1$, $\mathrm{M}_1$,
$\mathrm{M}_2$, $\mathrm{M}_0$, $\mathrm{I}_5$, $\mathrm{I}_6$, $\mathrm{I}_7$, $\mathrm{I}_8(\volthmhh)$,
$\mathrm{EP}(\mathrm{frob})$, $\mathrm{GI}$, $\mathrm{FL}(C',r,L)$, $\mathrm{EV}$ and
$\mathrm{SL}(C,r,L)$, then $T$ satisfies $\mathrm{T}_1$, $\mathrm{AF}(\volthmhh)$,
$\mathrm{T}_2(2C+3C'+5,r,L)$ and $\mathrm{MS}(2C+3C'+5,r,L)$.
\formallabel{GlobalTheory.CompanionTheory.Ar.Top.chain_tower}
\formallabel{GlobalTheory.T.tower_weil_inputs_combined}
\end{theorem}

The proof occupies \cref{pf:volthm:sec:models2,pf:volthm:sec:weil,pf:volthm:sec:asym,pf:volthm:sec:lower}
and is completed in \cref{pf:volthm:sec:lower}.

% ---------------------------------------------------------------------------
\subsection{Root models: rational weights, heights and degrees}\label{pf:volthm:sec:models2}

In this subsection $M$ is a root model of type $(\Psi,N)$ and $\volthmII$ a full intersection system.

\begin{lemma}\label{pf:volthm:lem:intcoef}
For a rational weight $\lambda$ and $s\in\Q$ there are an integer $q\ge1$ and integers $n_r$ with
$n_r=q\,\volthmcoef(\lambda,s)(r)$ for all $r\in(\Psi\times[N])\sqcup\{\infty\}$.
\formallabel{RootFieldNumerics.exists_integer_coefficients_rational}
\formallabel{GlobalTheory.CompanionTheory.Ar.Top.clear_rat_weight_coef}
\formallabel{GlobalTheory.T.rational_weight_integral_coeffs}
\end{lemma}

\begin{proof}
Take $q$ a common denominator of the finitely many rationals $-\lambda_r$ and $s/k_A$.
\end{proof}

\begin{lemma}[Combination lines]\label{pf:volthm:lem:comb}
Let $t,s\in\R$, $\lambda$ a weight, $q\ge1$ and $n_r=q\,\volthmcoef(\lambda,s)(r)$ integers, and
$L=\bigotimes_r\ell_t(r)^{n_r}$.
\begin{enumerate}[label=(\alph*)]
\item If $M$ satisfies $\mathrm{SP}(t,\lambda,s)$, then $L$ is semipositive on $P_M$ and conjugation
  invariant.
\item $(\deg L\cdot W)=q\sum_r\volthmcoef(\lambda,s)(r)(\deg\ell_1(r)\cdot W)$ for every
  $W$; hence $\mathrm{Nef}_{\mathrm f}(\lambda,s)$ implies that $L$ is nef on the fibres of $P_M$.
\item For every point $W$ of $M$: $\deg_\C(W)>0$ and $\volthmII(W;L)=q\,\mathrm{ht}_t(\lambda,s;W)\deg_\C(W)$.
\item $\int_{P_M}(\log h_L)^n=\delta q^n\,\mathrm{geo}_M(\volthmlin^+(\lambda,s)^n)$, and more generally
  $\int_{P_M}\prod_{k}\log h_{L_k}=\delta q^n\,\mathrm{geo}_M(\prod_k\volthmlin^+(\lambda_k,s_k))$ for lines
  $L_k$ attached in the same way to $(\lambda_k,s_k)$ and a common $q$.
\item $\volthmII(P_M;L,\dots,L)=\delta q^{n+1}\,\mathrm{ar}_t(\volthmlin^+(\lambda,s)^{n+1})$.
\end{enumerate}
\formallabel{GlobalTheory.T.root_model_semipositive_conj}
\formallabel{GlobalTheory.CompanionTheory.Ar.Top.semipositive_integral_multiple}
\formallabel{GlobalTheory.T.root_model_curve_degree_nef}
\formallabel{GlobalTheory.CompanionTheory.Ar.Top.ht_combLine_integral_multiple}
\formallabel{GlobalTheory.T.root_model_point_height_degree}
\formallabel{GlobalTheory.T.arith_intersection_multilinear}
\formallabel{GlobalTheory.CompanionTheory.Ar.Top.intersection_combLine}
\formallabel{GlobalTheory.T.root_model_log_integral_formula}
\formallabel{GlobalTheory.T.root_model_arith_intersection_formula}
\end{lemma}

\begin{proof}
(a) is \lref{RootFieldNumerics.root_model_line_semipositive}, which states exactly this: the weight of
$L$ is $q$ times $\sum_r\volthmcoef(\lambda,s)(r)\log h_{\ell_t(r)}$, and products of conjugation-invariant
metrics are conjugation invariant. (b) is \lref{GlobalTheory.CompanionTheory.Ar.Top.nefFib_integral_multiple}
(linearity of the degree on curves in the multidegree). (c) $\deg_\C(W)>0$ for an arithmetic curve is
\lref{GroundTheory.Heights₂.Heights.LowerBoundFree.degC_pos_of_isArithVar_one}; by multilinearity of a
full intersection system in its smooth entries
(\lref{RootFieldNumerics.intersection_multilinear_one_position}: $\volthmII(P;\dots,\bigotimes_j\ell_j^{n_j},
\dots)=\sum_jn_j\volthmII(P;\dots,\ell_j,\dots)$),
$\volthmII(W;L)=\sum_rq\volthmcoef(\lambda,s)(r)\volthmII(W;\ell_t(r))=q\,\mathrm{ht}_t(\lambda,s;W)\deg_\C(W)$.
(d) and (e) are \lref{GlobalTheory.CompanionTheory.Ar.Top.geoFun_comb_mixed} and
\lref{GlobalTheory.CompanionTheory.Ar.Top.top_comb_eq_arFun}, which state exactly these identities; they
follow by multilinearity and symmetry, the factor $\delta k_A^{m(\infty)}$ in the definitions of
$\mathrm{ar}_t$ and $\mathrm{geo}_M$ compensating the normalisation $s/k_A$ of the coefficient of $A$.
\end{proof}

\begin{proposition}[Zhang's bounds on a root model]\label{pf:volthm:prop:zhang}
Let $t\in\R$, $\lambda$ a rational weight, $s\in\Q$ and $e\in\R$, and assume that $M$ satisfies
$\mathrm{SP}(t,\lambda,s)$ and $\mathrm{Nef}_{\mathrm f}(\lambda,s)$.
\begin{enumerate}[label=(\alph*)]
\item If $(\mathrm{Z}^-)$ holds and $\mathrm{ht}_t(\lambda,s;W)\ge e$ for all points $W$ of $M$, then
  \[
  (n+1)\,e\,\mathrm{geo}_M(\volthmlin^+(\lambda,s)^n)\le\mathrm{ar}_t(\volthmlin^+(\lambda,s)^{n+1}).
  \]
\item If $(\mathrm{Z}^+)$ holds and $F$ is a Zariski dense set of points $W$ of $M$ with
  $\mathrm{ht}_t(\lambda,s;W)\le e$, then
  \[
  \mathrm{ar}_t(\volthmlin^+(\lambda,s)^{n+1})\le(n+1)\,e\,\mathrm{geo}_M(\volthmlin^+(\lambda,s)^n).
  \]
\end{enumerate}
\formallabel{GlobalTheory.CompanionTheory.Ar.Top.zhang_lower_bound_root_model}
\formallabel{GlobalTheory.T.zhang_lower_bound_rational}
\formallabel{GlobalTheory.CompanionTheory.Ar.Top.zhang_upper_bound_root_model}
\formallabel{GlobalTheory.T.zhang_upper_bound_dense_points}
\end{proposition}

\begin{proof}
Choose $q,n_r$ by \cref{pf:volthm:lem:intcoef} and let $L$ be as in \cref{pf:volthm:lem:comb}. By (a) and (b)
of that lemma $L$ satisfies the positivity hypotheses of the successive minima bounds on $P_M$, which has
dimension $n+1$. (a) By \cref{pf:volthm:lem:comb}(c), $qe\deg_\C(W)\le\volthmII(W;L)$ for every arithmetic
curve $W\supseteq P_M$, so $(\mathrm{Z}^-)$ gives
$(n+1)qe\deg(P_M;L)\le\volthmII(P_M;L^{n+1})$. By \cref{pf:volthm:lem:comb}(d),(e) this reads
$(n+1)qe\,\delta q^n\mathrm{geo}_M(\volthmlin^+(\lambda,s)^n)\le\delta q^{n+1}\mathrm{ar}_t(\volthmlin^+(\lambda,s)^{n+1})$;
divide by $\delta q^{n+1}>0$. (b) is identical, with $(\mathrm{Z}^+)$ applied to the dense
set $F$.
\end{proof}

\begin{lemma}[Nef cones]\label{pf:volthm:lem:nefcone}
The curve degree $(\lambda,s)\mapsto\sum_r\volthmcoef(\lambda,s)(r)(\deg\ell_1(r)\cdot W)$ is linear.
Hence for $s\ge0$: $\mathrm{Nef}_p(\lambda,0)$ and $\mathrm{Nef}_p(0,1)$ imply $\mathrm{Nef}_p(\lambda,s)$, and
$\mathrm{Nef}_{\mathrm g}(\lambda,0)$ and $\mathrm{Nef}_{\mathrm g}(0,1)$ imply $\mathrm{Nef}_{\mathrm g}(\lambda,s)$.
\formallabel{GlobalTheory.CompanionTheory.Ar.Top.curveDeg_linear_nef}
\formallabel{GlobalTheory.T.curve_degree_linear_nef_cone}
\end{lemma}

\begin{proof}
$\volthmcoef(\lambda,s)$ is linear in $(\lambda,s)$, and $\volthmcoef(\lambda,s)=\volthmcoef(\lambda,0)+s\,\volthmcoef(0,1)$.
\end{proof}

\begin{lemma}[Auxiliary inputs]\label{pf:volthm:lem:aux}
\begin{enumerate}[label=(\alph*)]
\item For all coordinate data $(K,\iota)$ there is an intersection family over $\Q$ satisfying the axioms of
  intersection theory over a field, and any two such families agree on projective varieties.
\item If $P$ is an arithmetic variety of dimension $n+1$ and $s$ is a section of a Hermitian line on $P$, every
  component $W$ of the divisor of $s$ is either horizontal and an arithmetic variety of dimension $n$, or
  vertical, of prime characteristic and of dimension $n+|K|$.
\end{enumerate}
\formallabel{GlobalTheory.T.exists_unique_rational_statement}
\formallabel{GlobalTheory.CompanionTheory.Ar.Link.divSuppArith_holds}
\end{lemma}

\begin{proof}
(a) is \lref{GlobalTheory.CompanionTheory.Ar.PB.Known.fieldNumbers} for the field $\Q$, and (b) is
\lref{ReciprocityTheory.ChartRestriction.G.Geo.divSuppArith}.
\end{proof}

\begin{proposition}[Degree bounds]\label{pf:volthm:prop:degbounds}
Assume $(\mathrm{AI})$. Let $T$ be a root tower satisfying $\mathrm{Pr}$ and $\mathrm{I}_7$
with respect to a full $\volthmII$. For every $B$, $c$, rational weight $\lambda$ and rational $s\ge0$ such that
$\lambda$ and $\lambda_s$ are dominant,
\[
0\le\mathrm{dg}(B,c,\lambda,s)\le\vol(c)\volthmfd(\lambda).
\]
\formallabel{RootFieldNumerics.tower_rational_weight_bound}
\formallabel{GlobalTheory.CompanionTheory.Ar.Top.deg_bounds}
\formallabel{GlobalTheory.T.tower_component_height_bound}
\end{proposition}

\begin{proof}
The two positivity results used below assume, besides the component $(\mathrm{K})$ of
$(\mathrm{AI})$ (nonnegativity of intersection numbers of nef classes over a field), the two statements
of \cref{pf:volthm:lem:aux}.

Write $M=M_{B,c}$ and $\ell(\mu,\sigma)=\volthmlin^+(\mu,\sigma)$. Since $\nu$ is rational, $\lambda_s$ is a
rational weight. By $\mathrm{I}_7$ and \cref{pf:volthm:lem:nefcone}, $M$ satisfies $\mathrm{Nef}_{\mathrm g}(\lambda,0)$
and $\mathrm{Nef}_{\mathrm g}(\lambda_s,s)$.

\emph{Nonnegativity.} Let $L$ be the line attached to $(\lambda_s,s)$ by \cref{pf:volthm:lem:comb} (with $t=1$).
By multilinearity of geometric degrees
(\lref{GlobalTheory.T.log_integral_multilinear}), $\deg(W;L)=q\sum_r\volthmcoef(\lambda_s,s)(r)
\deg(W;\ell_1(r))\ge0$ for every arithmetic surface $W\supseteq P_M$. By the arithmetic Kleiman criterion in the
form \lref{GlobalTheory.T.ptsIntegral_log_nonneg} (if each of $L_1,\dots,L_n$ has nonnegative degree on
all arithmetic surfaces containing $P$, then $\int_P\prod\log h_{L_k}\ge0$), $\int_{P_M}(\log h_L)^n\ge0$, and
\cref{pf:volthm:lem:comb}(d) gives $\mathrm{dg}(B,c,\lambda,s)=\mathrm{geo}_M(\ell(\lambda_s,s)^n)\ge0$.

\emph{Upper bound.} By $\mathrm{Pr}$,
$\vol(c)\volthmfd(\lambda)=\mathrm{geo}_M((\volthmlin\lambda)^{n,+})=\mathrm{geo}_M(\ell(\lambda,0)^n)$. Since
$\ell(\lambda,0)-\ell(\lambda_s,s)=s\,\ell(\nu,-1)$, the identity
$a^n-b^n=(a-b)\sum_{i<n}a^ib^{n-1-i}$ gives
\[
\vol(c)\volthmfd(\lambda)-\mathrm{dg}(B,c,\lambda,s)=s\sum_{i=0}^{n-1}
\mathrm{geo}_M\bigl(\ell(\lambda,0)^i\,\ell(\lambda_s,s)^{n-1-i}\,\ell(\nu,-1)\bigr).
\]
Each term is $\ge0$: by \cref{pf:volthm:lem:comb}(d) it is a positive multiple of
$\int_{P_M}\prod_{k<n-1}\log h_{L_k}\cdot\log h_{L'}$, where the $L_k$ are lines attached to $(\lambda,0)$ or
$(\lambda_s,s)$, hence have nonnegative degree on arithmetic surfaces, and $L'$ is the line attached to
$(\nu,-1)$, which has a nonzero integral section by $\mathrm{Eff}(\nu)$ (part of $\mathrm{I}_7$). The
nonnegativity of such integrals (``effective times nef'') is
\lref{GlobalTheory.T.arith_intersection_section_bound}, a result of the development whose proof we do not
reproduce: for an arithmetic variety $P_y$ of dimension $e+2$ with $y$ birational, smooth lines
$L_1,\dots,L_e$ of nonnegative degree on all arithmetic surfaces containing $P_y$, and a smooth line $A$ with a
nonzero integral section, $\int_{P_y}\log h_{L_1}\cdots\log h_{L_e}\log h_A\ge0$.
\end{proof}

% ---------------------------------------------------------------------------
\subsection{Limit formulas and the Weil bound on walls}\label{pf:volthm:sec:weil}

\begin{proposition}[$\mathrm{I}_3$]\label{pf:volthm:prop:I3}
If a root tower satisfies $\mathrm{A}_1$, $\mathrm{M}_1$, $\mathrm{M}_2$ and $\mathrm{M}_0$, it satisfies
$\mathrm{I}_3$. More precisely, $\mathrm{M}_1$ and $\mathrm{M}_2$ imply that $M_{B,c}$ satisfies
$\mathrm{Nef}_{\mathrm f}(\mu,0)$ for every dominant $\mu$ with $\mathcal N_j(B,\mu)$.
\formallabel{GlobalTheory.CompanionTheory.Ar.Top.RootTower.tower_I3_of_nef_inputs}
\formallabel{GlobalTheory.T.tower_I3_of_semipositive_nef}
\formallabel{GlobalTheory.CompanionTheory.Ar.Top.RootTower.nefFib_of_nefThreshold}
\formallabel{GlobalTheory.CompanionTheory.Ar.Top.NefThreshold.mono}
\end{proposition}

\begin{proof}
Let $\mu$ be dominant with $\mathcal N_j(B,\mu)$, and $p$ a prime. If $p$ is in $S_B$, $\mathrm{M}_1$ gives
$\mathrm{Nef}_p(\mu,0)$. Otherwise $p>B$; since $\volthmgap_j(\mu,\tau)\ge0$ for dominant $\mu$, the condition
$\mathcal N_j(b,\mu)$ is monotone in $b$, so $\mathcal N_j(p,\mu)$ holds and $\mathrm{M}_2$ gives
$\mathrm{Nef}_p(\mu,0)$. Given $a$, take $s_0$ from $\mathrm{A}_1$. For $\lambda,s$ as in $\mathrm{I}_3$, apply
this to $\mu=\lambda_s$; with $\mathrm{M}_0$ and \cref{pf:volthm:lem:nefcone} ($s\ge0$) we get
$\mathrm{Nef}_{\mathrm f}(\lambda_s,s)$, and $\mathrm{A}_1$ gives semipositivity.
\end{proof}

\begin{proposition}[Limit formulas, $\mathrm{T}_1$]\label{pf:volthm:prop:T1}
Let $\Psi$ be nonempty. A root tower satisfying $\mathrm{Pr}$, $\mathrm{I}_2$ and $\mathrm{I}_3$
satisfies $\mathrm{T}_1$.
\formallabel{GlobalTheory.CompanionTheory.Ar.Top.LimitVolumeRatio.tower_T1_limit_of_inputs}
\formallabel{GlobalTheory.T.tower_T1_limit}
\end{proposition}

\begin{proof}
Let $u,g$ satisfy $\mathrm{WH}_j(u,g)$, $B_0(u,g)<B+1$, and fix $c$; write $M=M_{B,c}$. Since
$|\Psi|\ge1$, $N\ge2$ and $\|u\|/g\ge0$, $B_0(u,g)\ge16$, so $B\ge16\ge2$.

\emph{Eventual dominance and threshold.} As $u$ is regular dominant and $u_s=u-s\nu$ depends continuously on
$s$, $u_s$ is dominant for $s$ near $0$. At $s=0$ the threshold holds strictly:
$4|\Psi|\,\|u\|<(B-2)\volthmgap_j(u,\tau)$ for all $\tau$. Indeed, with $F=|\Psi|$ and $x=\|u\|$,
$B-2>B_0(u,g)-3=4FN^2+4FN^2x/g-3$, so $(B-2)g>4FN^2g+4FN^2x-3g\ge4Fx+g(4FN^2-3)>4Fx$, and
$\volthmgap_j(u,\tau)\ge g$. Since both sides are continuous in $s$ ($\volthmgap_j(u_s,\tau)=\volthmgap_j(u,\tau)-
s\volthmgap_j(\nu,\tau)$), $\mathcal N_j(B,u_s)$ holds for $s$ near $0$. Now $\mathrm{I}_3$ (with $a=\|u\|_1$)
gives the first assertion of $\mathrm{T}_1$.

\emph{Limits.} By $\mathrm{I}_2$ (with $a=\|u\|_1$) there are $s_1,C_2$ with
$|\mathrm{ar}(B,c,u,s)-\Pi^+(B,c,u_s,s)|\le C_2s$ for small $s>0$. Moreover
$\volthmlin^+(u_s,s)=\volthmlin(u)^++s\,(x_\infty-\volthmlin(\nu)^+)$, so for any linear form $\Phi$ on
$\volthmR^+$ and $d\ge0$, $s\mapsto\Phi(\volthmlin^+(u_s,s)^d)$ is a polynomial in $s$ with value
$\Phi(\volthmlin(u)^{d,+})$ at $0$. Applied to $\Phi=\volthmPP_M$, $d=n+1$, this gives
$\Pi^+(B,c,u_s,s)\to\volthmPP_{B,c}(\volthmlin(u)^{n+1})$, hence the same limit for
$\mathrm{ar}(B,c,u,s)$. Applied to $\Phi=\mathrm{geo}_M$, $d=n$, with $\mathrm{Pr}$, it gives
$\mathrm{dg}(B,c,u,s)\to\mathrm{geo}_M(\volthmlin(u)^{n,+})=\vol(c)\volthmfd(u)$, which is $>0$ by
$\mathrm{Pr}$ and \cref{pf:volthm:lem:weyl} ($u$ is regular dominant). The limit of the quotient follows.
\end{proof}

\begin{lemma}[Elementary facts on weights]\label{pf:volthm:lem:weights}
Let $1\le j<N$.
\begin{enumerate}[label=(\alph*)]
\item For every $d\ge0$ there is $C_p\ge0$ such that for every linear form $\Phi$ on $\volthmR^+$ with
  $|\Phi(x^m)|\le b$ for all monomials, and all weights $\lambda,\mu$ with $\|\lambda\|,\|\mu\|\le1$ and
  $|s|,|s'|\le1$: $|\Phi(\volthmlin^+(\lambda,s)^d)-\Phi(\volthmlin^+(\mu,s')^d)|\le C_pb(\|\lambda-\mu\|+|s-s'|)$;
  moreover $(\lambda,s)\mapsto\Phi(\volthmlin^+(\lambda,s)^d)$ is continuous.
\item There is $C_8\ge0$ with $\volthmfd(\lambda)\le C_8\volthmgap_j(\lambda,\tau_0)$ for all dominant $\lambda$ with
  $\|\lambda\|\le1$ and all $\tau_0$.
\item Let $c_0>0$, $\tau_0\in\Psi$ and $\Phi$ a continuous function on weights. Call $\lambda$ \emph{in the
  wall} $\mathcal W_j(\tau_0,c)$ if $\|\lambda\|\le\frac12$, $\volthmgap_j(\lambda,\tau_0)=0$,
  $\volthmgap_j(\lambda,\tau)\ge c$ for $\tau\ne\tau_0$, and $\lambda$ has compact gaps $\ge c$.
  If $|\Phi|\le b$ on all rational weights in $\mathcal W_j(\tau_0,c_0/2)$, then $|\Phi|\le b$ on
  $\mathcal W_j(\tau_0,c_0)$.
\end{enumerate}
\formallabel{GlobalTheory.CompanionTheory.Ar.Top.linA_pow_continuous_lipschitz}
\formallabel{GlobalTheory.T.arith_functional_lipschitz_weights}
\formallabel{RootFieldNumerics.flagDeg_le_gap}
\formallabel{GlobalTheory.CompanionTheory.Ar.Top.flagDeg_le_mul_gap}
\formallabel{GlobalTheory.T.flag_degree_le_gap}
\formallabel{RootFieldNumerics.continuous_bound_wall_extension}
\formallabel{GlobalTheory.CompanionTheory.Ar.Top.rat_points_suffice_on_walls}
\formallabel{GlobalTheory.T.continuous_wall_bound_extend}
\end{lemma}

\begin{proof}
(a) $\volthmlin^+(\lambda,s)^d=\sum_m c_m(\lambda,s)x^m$ where each $c_m$ is a polynomial in the finitely many
coordinates of $(\lambda,s)$; on the unit ball each $c_m$ is Lipschitz with a constant depending only on
$d,\Psi,N$, and $\Phi(\volthmlin^+(\lambda,s)^d)=\sum_mc_m(\lambda,s)\Phi(x^m)$.

(b) By \cref{pf:volthm:lem:weyl}, $\volthmfd(\lambda)=n!\prod_{\tau}\prod_{i<k}f_{\tau,i,k}$ with
$f_{\tau,i,k}=(\lambda_{\tau,i}-\lambda_{\tau,k})/(k-i)\in[0,2]$ for $\lambda$ dominant with $\|\lambda\|\le1$,
and the factor $f_{\tau_0,j-1,j}$ equals $\volthmgap_j(\lambda,\tau_0)$. Bounding the other factors by $2$ gives
$C_8=n!\,2^{|\Psi|N^2}$.

(c) Let $\lambda\in\mathcal W_j(\tau_0,c_0)$. For $k\in\N$ choose rationals $q_k(\tau,i)$ with
$|q_k(\tau,i)|\le|\lambda_{\tau,i}|$, $|q_k(\tau,i)-\lambda_{\tau,i}|\le\frac{c_0}{4(k+1)}$, and
$q_k(\tau,j-1)=q_k(\tau,j)$ whenever $\lambda_{\tau,j-1}=\lambda_{\tau,j}$ (choose the rational approximation as
a function of the real number $\lambda_{\tau,i}$). Then $q_k$ is rational, $\|q_k\|\le\frac12$,
$\volthmgap_j(q_k,\tau_0)=0$, and all other gaps change by at most $c_0/2$, so
$q_k\in\mathcal W_j(\tau_0,c_0/2)$; and $q_k\to\lambda$. By continuity $|\Phi(\lambda)|=\lim|\Phi(q_k)|\le b$.
\end{proof}

\begin{lemma}[Perturbation off the wall]\label{pf:volthm:lem:perturb}
Let $T$ be a root tower and $c_0,\theta_0>0$. There is $B_1$ such that for all $B\ge B_1$, $\tau_0$ and
$\lambda^0\in\mathcal W_j(\tau_0,c_0)$, with $\theta=5|\Psi|/B$ and
$\lambda=\lambda^0+2\theta\,\mathbf 1_{\{\tau_0\}\times\{i<j\}}$: $B\ge3$, $0<\theta\le\min(\theta_0,1)$,
$\|\lambda\|\le1$, $\|\lambda_{\theta^2}\|\le1$, $\volthmgap_j(\lambda,\tau)\ge\theta$ for all $\tau$,
$\volthmgap_j(\lambda,\tau_0)=2\theta$, $\lambda$ has compact gaps $\ge c_0$, $\lambda$ and $\lambda_{\theta^2}$ are
dominant, $\mathcal N_j(B,\lambda_{\theta^2})$ holds, and
$\|\lambda_{\theta^2}-\lambda^0\|\le2\theta+\|\nu\|\theta^2$.
\formallabel{RootFieldNumerics.perturbed_weight_wall_bounds}
\formallabel{GlobalTheory.CompanionTheory.Ar.Top.hypotheses_hold_for_large_bound}
\formallabel{GlobalTheory.T.perturbed_weight_gap_bounds}
\end{lemma}

\begin{proof}
Elementary; $B_1$ is chosen so that $\theta=5|\Psi|/B$ is small compared with $\theta_0$, $1/4$, $c_0$ and
$1/(1+\|\nu\|)$. Adding $2\theta$ to the entries $(\tau_0,i)$, $i<j$, does not change the compact gaps (both
indices of a compact gap lie on the same side of $j$), raises $\volthmgap_j(\cdot,\tau_0)$ from $0$ to $2\theta$,
leaves the other gaps $\ge c_0\ge\theta$, and changes the norm by $2\theta$; as $\lambda^0$ is dominant
(its gaps are $\ge0$), so is $\lambda$. Subtracting $\theta^2\nu$ changes every gap by at most
$2\|\nu\|\theta^2$, which keeps all gaps positive and the norm $\le1$ for $\theta$ small. Finally
$4|\Psi|\|\lambda_{\theta^2}\|\le4|\Psi|$, while
$(B-2)\volthmgap_j(\lambda_{\theta^2},\tau)\ge(B-2)(\theta-2\|\nu\|\theta^2)=5|\Psi|(1-\frac2B)(1-2\|\nu\|\theta)$,
which is $\ge4|\Psi|$ for $B$ large.
\end{proof}

\begin{proposition}[Squeeze]\label{pf:volthm:prop:squeeze}
Assume $(\mathrm{Z}^-)$ and $(\mathrm{Z}^+)$, and let $T$ be a root
tower satisfying $\mathrm{I}_3$, $\mathrm{I}_5$ and $\mathrm{I}_6$ with respect to a full $\volthmII$. For every
$c_0>0$ there are $\theta_0>0$ and $C_s\ge0$ such that
\[
|\mathrm{ar}(B,c,\lambda,s)|\le C_s\log(B/s)\,\mathrm{dg}(B,c,\lambda,s)
\]
for all $B\ge2$, $c$, rational $\lambda$ with $\|\lambda\|\le1$, $0<\theta\le\theta_0$ and rational
$s=\theta^2$, such that $\volthmgap_j(\lambda,\tau)\ge\theta$ for all $\tau$, $\lambda$ has compact gaps $\ge c_0$,
$\lambda_s$ is dominant, $\mathcal N_j(B,\lambda_s)$ holds and $\mathrm{dg}(B,c,\lambda,s)\ge0$.
\formallabel{GlobalTheory.CompanionTheory.Ar.Top.mbarTop_squeeze}
\formallabel{GlobalTheory.T.tower_gap_height_uniform_bound}
\end{proposition}

\begin{proof}
Take $s_0$ from $\mathrm{I}_3$ with $a=|\Psi|N$, $(\theta_5,C_5)$ from $\mathrm{I}_5$ and $(\theta_6,C_6)$ from
$\mathrm{I}_6$ for $c_0$, and put $\theta_0=\min(\theta_5,\theta_6,1,s_0)$ and
$C_s=(n+1)(\max(C_5,0)+\max(C_6,0))$. Then $s=\theta^2\le\theta\le\min(1,s_0)$ and
$\|\lambda\|_1\le|\Psi|N$, so $\mathrm{I}_3$ gives $\mathrm{Nef}_{\mathrm f}(\lambda_s,s)$ and
$\mathrm{SP}(\mathrm{scale}(s),\lambda_s,s)$ for $M=M_{B,c}$; $\lambda_s$ is rational. By $\mathrm{I}_5$ every
point $W$ has $\mathrm{ht}_{\mathrm{scale}(s)}(\lambda_s,s;W)\ge-C_5\log(B/s)$, and by $\mathrm{I}_6$ a Zariski dense
set of points has $\mathrm{ht}_{\mathrm{scale}(s)}(\lambda_s,s;W)\le C_6\log B$. \Cref{pf:volthm:prop:zhang} with
$t=\mathrm{scale}(s)$ and the weight $\lambda_s$ gives
\[
-(n+1)C_5\log(B/s)\,\mathrm{dg}\le\mathrm{ar}\le(n+1)C_6\log B\,\mathrm{dg}
\]
($\mathrm{ar}=\mathrm{ar}(B,c,\lambda,s)$, $\mathrm{dg}=\mathrm{dg}(B,c,\lambda,s)$). As $0\le\log B\le\log(B/s)$
(since $s\le1$) and $\mathrm{dg}\ge0$, the claim follows.
\end{proof}

\begin{proposition}[Weil bound on walls]\label{pf:volthm:prop:ThmW}
Assume $(\mathrm{AI})$, $(\mathrm{Z}^-)$ and $(\mathrm{Z}^+)$, and let
$\volthmII$ be full. Let $T$ be a root tower satisfying $\mathrm{Pr}$, $\mathrm{I}_1$, $\mathrm{I}_2$,
$\mathrm{I}_3$, $\mathrm{I}_5$, $\mathrm{I}_6$ and $\mathrm{I}_7$. Then $T$ satisfies $\mathrm{WB}$: for every
$c_0>0$ there are $C$ and $B_1$ such that
\[
|\Pi(B,\lambda)|\le C\,\frac{\log B}{B}\qquad\text{for all }B\ge B_1,\ \tau_0\in\Psi,\
\lambda\in\mathcal W_j(\tau_0,c_0).
\]
\formallabel{GlobalTheory.CompanionTheory.Ar.Top.thmW_tower}
\formallabel{GlobalTheory.T.tower_weil_theorem}
\end{proposition}

\begin{proof}
Fix $c_0>0$. Take $(\theta_0,C_s)$ from \cref{pf:volthm:prop:squeeze} for $c_0/2$; $(s_0,C_3)$ from $\mathrm{I}_2$
with $a=|\Psi|N$; $C_2$ from $\mathrm{I}_1$; $C_8$ from \cref{pf:volthm:lem:weights}(b); $C_p$ from
\cref{pf:volthm:lem:weights}(a) with $d=n+1$; and $B_1$ from \cref{pf:volthm:lem:perturb} for $c_0/2$ and
$\min(\theta_0,s_0)$. Put
\[
C=5|\Psi|\sum_c\bigl(|C_3|+6C_sC_8\vol(c)+C_p|C_2|(3+\|\nu\|)\bigr).
\]
Let $B\ge B_1$ and $\tau_0\in\Psi$. The function $\lambda\mapsto\Pi(B,\lambda)=\sum_c\volthmPP_{M_{B,c}}
(\volthmlin^+(\lambda,0)^{n+1})$ is continuous (\cref{pf:volthm:lem:weights}(a)), so by
\cref{pf:volthm:lem:weights}(c) it suffices to show $|\Pi(B,l)|\le C\log B/B$ for rational
$l\in\mathcal W_j(\tau_0,c_0/2)$. For such $l$ let $\theta=5|\Psi|/B$, $\lambda=l+2\theta\mathbf 1_{\{\tau_0\}
\times\{i<j\}}$ and $s=\theta^2$ (rational); \cref{pf:volthm:lem:perturb} provides all the properties listed
there. Fix $c$ and put
\[
a=\mathrm{ar}(B,c,\lambda,s),\quad p=\Pi^+(B,c,\lambda_s,s),\quad
q=\volthmPP_{B,c}(\volthmlin(l)^{n+1})=\volthmPP_{M_{B,c}}(\volthmlin^+(l,0)^{n+1}).
\]
\begin{enumerate}[label=(\roman*)]
\item By \cref{pf:volthm:prop:degbounds}, $0\le\mathrm{dg}(B,c,\lambda,s)\le\vol(c)\volthmfd(\lambda)\le
  \vol(c)C_8\volthmgap_j(\lambda,\tau_0)=2\theta C_8\vol(c)$, and by \cref{pf:volthm:prop:squeeze},
  $|a|\le C_s\log(B/s)\,\mathrm{dg}$. Since $5|\Psi|\ge1$, $B/s=B^3/(5|\Psi|)^2\le B^3$, so
  $0\le\log(B/s)\le3\log B$ and $|a|\le6C_sC_8\vol(c)\,\theta\log B$.
\item By $\mathrm{I}_2$ ($\|\lambda\|_1\le|\Psi|N$, $s\le s_0$, $\lambda_s$ dominant), $|a-p|\le C_3s\le|C_3|\theta\log B$
  (using $\theta\le1\le\log B$, as $B\ge3$).
\item By $\mathrm{I}_1$, $|\volthmPP_{M_{B,c}}(x^m)|\le|C_2|\log B$ for all monomials, so
  \cref{pf:volthm:lem:weights}(a) gives $|p-q|\le C_p|C_2|\log B\,(\|\lambda_s-l\|+s)\le
  C_p|C_2|\log B\,(2\theta+\|\nu\|\theta^2+\theta^2)\le C_p|C_2|(3+\|\nu\|)\,\theta\log B$.
\end{enumerate}
Hence $|q|\le|a|+|a-p|+|p-q|\le(|C_3|+6C_sC_8\vol(c)+C_p|C_2|(3+\|\nu\|))\theta\log B$, and summing over $c$,
$|\Pi(B,l)|\le C\log B/B$.
\end{proof}

% ---------------------------------------------------------------------------
\subsection{The asymptotic formula}\label{pf:volthm:sec:asym}

\begin{lemma}[The test polynomial]\label{pf:volthm:lem:testpoly}
Let $1\le j<N$. Then $\mathrm{tp}_j$ is homogeneous of degree $n+1$ and $\mathrm{td}_j>0$. If $N-j\le j$, then
$\volthmSh_j^\tau(\mathrm{tp}_j)=0$ and $\volthmTot^\tau(\mathrm{tp}_j)=(d_X+1)(N-j)\,\mathrm{td}_j$ for all $\tau$; if
$N-j>j$, then $\volthmSh_j^\tau(\mathrm{tp}_j)=\volthmTot^\tau(\mathrm{tp}_j)=-(d_X+1)\,j\,\mathrm{td}_j$ for all $\tau$.
\formallabel{GlobalTheory.CompanionTheory.Ar.Top.testPoly_facts}
\formallabel{GlobalTheory.T.test_polynomial_shift_total}
\end{lemma}

\begin{proof}
$\mu_j$ is linear and $\phi_j$ is homogeneous of degree $|\Psi|(\binom j2+\binom{N-j}2)$; since
$\binom N2=\binom j2+\binom{N-j}2+j(N-j)$, $\mathrm{tp}_j$ has degree $d_X+1+\deg\phi_j=n+1$. The positivity
$\mathrm{td}_j>0$ is \lref{RootFieldNumerics.TF.testDeg_pos} ($\mathrm{td}_j$ is, up to the normalisation of
$\volthmdeg$, the degree of the product over $\Psi$ of the Grassmannians $\mathrm{Gr}(j,N)$ in their Pl\"ucker
embeddings, $\phi_j$ being the class of a fibre of $\mathrm{Fl}\to\prod_\Psi\mathrm{Gr}(j,N)$); we do not
reproduce its proof. For the values, note that $\phi_j$ is, for each $\tau$, a product of two Vandermonde
polynomials in the variables $x_{\tau,i}$, $i<j$, and $x_{\tau,i}$, $i\ge j$, which are invariant under
translation of either group of variables; hence $\sum_{i<j}\partial_{\tau,i}\phi_j=\sum_i\partial_{\tau,i}\phi_j=0$.
If $N-j\le j$, $\partial_{\tau,i}\mu_j=\mathbf 1_{i\ge j}$, so by the Leibniz rule
$\sum_{i<j}\partial_{\tau,i}\mathrm{tp}_j=0$ and $\sum_i\partial_{\tau,i}\mathrm{tp}_j=(d_X+1)(N-j)\mu_j^{d_X}\phi_j$.
If $N-j>j$, $\partial_{\tau,i}\mu_j=-\mathbf 1_{i<j}$, and both sums equal $-(d_X+1)j\mu_j^{d_X}\phi_j$. Apply
$\volthmdeg$.
\end{proof}

\begin{proposition}[The flag height functional]\label{pf:volthm:prop:hh}
Assume $(\mathrm{FH})$ and $(\mathrm{AI})$, let $\volthmII$ be full, $\Psi$ nonempty and $N\ge2$, and write
$h=\volthmhh^{\volthmII}$.
\begin{enumerate}[label=(\alph*)]
\item $h(\volthmlin(\lambda)^{n+1})=H_{\mathrm{Fl}}(\lambda)$ for every weight $\lambda$.
\item $h(p_1(\tau)Q)=0$ for $Q$ homogeneous of degree $n$.
\item $h$ satisfies $\mathrm{I}_9(h)$ and $\mathrm{LI}(h)$.
\item If $1\le j<N$ and $(\mathrm{GH})$ holds, then
  $h(\mathrm{tp}_j)\le(d_X+1)\,\mathrm{td}_j\,|\Psi|\,\frac{\min(j,N-j)}2\log(N+1)$.
\end{enumerate}
\formallabel{GlobalTheory.CompanionTheory.Ar.Top.hhFl_linRoot_pow}
\formallabel{GlobalTheory.T.flag_height_polynomial_eval}
\formallabel{RootFieldNumerics.flag_height_power_sum_zero}
\formallabel{GlobalTheory.CompanionTheory.Ar.Top.hhFl_pk_one_mul}
\formallabel{GlobalTheory.T.flag_height_power_sum_one_zero}
\formallabel{GlobalTheory.CompanionTheory.Ar.Top.flag_height_I9}
\formallabel{GlobalTheory.T.flag_height_I9_properties}
\formallabel{GlobalTheory.CompanionTheory.Ar.Top.hhFl_labelInv}
\formallabel{GlobalTheory.T.flag_height_label_invariant}
\formallabel{GlobalTheory.CompanionTheory.Ar.Top.HhTestProof.hhTestBound_hhFl}
\formallabel{GlobalTheory.T.flag_height_test_bound}
\end{proposition}

\begin{proof}
By \lref{GlobalTheory.T.flag_variety_boundary_dimension_drop}, $P_{\mathrm{Fl}}$ is an arithmetic variety of
dimension $n+1$. All Fubini--Study lines are smooth, and $\volthmII(P_{\mathrm{Fl}};\cdot)$ is symmetric in its
entries.

(a) Expanding, $\volthmlin(\lambda)^{n+1}=(-1)^{n+1}\sum_w\bigl(\prod_k\lambda_{w(k)}\bigr)x_{w(1)}\cdots x_{w(n+1)}$
over words $w$ of length $n+1$ in $\Psi\times[N]$, and by symmetry $h(x_{w(1)}\cdots x_{w(n+1)})=(-1)^{n+1}
\volthmII(P_{\mathrm{Fl}};\xi(w(1)),\dots,\xi(w(n+1)))$. Hence $h(\volthmlin(\lambda)^{n+1})=\sum_w\prod_k\lambda_{w(k)}
\volthmII(P_{\mathrm{Fl}};\xi(w))$, which is the polarisation $H_{\mathrm{Fl}}(\lambda)$.

(b) Since the Fubini--Study line of multidegree $a+b$ is the product of those of multidegrees $a$ and $b$,
multilinearity of $\volthmII$ makes $(a_1,\dots,a_{n+1})\mapsto\volthmII(P_{\mathrm{Fl}};\mathrm{FS}(a_1),\dots,
\mathrm{FS}(a_{n+1}))$ multilinear in the multidegrees. For a word $w$ of length $n$,
$h(p_1(\tau)x_w)=(-1)^{n+1}\sum_i\volthmII(P_{\mathrm{Fl}};\xi(w),\xi(\tau,i))$, and
$\sum_i(e_{\tau,i}-e_{\tau,i-1})=0$, so this vanishes. Expand $Q$ in monomials.

(c) The first two parts of $\mathrm{I}_9$ follow from (a) and $(\mathrm{FH})$ (for dominant $\lambda$ with
$\lambda_{\tau,a}=\lambda_{\tau,a+1}$ the gap at $(\tau,a)$ vanishes); the third is (b). For $\gamma\in\mathrm{Sym}(\Psi)$
the permutation of blocks $(\tau,q)\mapsto(\gamma\tau,q)$ maps $P_{\mathrm{Fl}}$ onto itself and $\xi(\tau,i)$ to
$\xi(\gamma\tau,i)$; the projection formula of a full intersection system for this block map (relative degree
$1$) gives $\volthmII(P_{\mathrm{Fl}};\xi(\gamma w))=\volthmII(P_{\mathrm{Fl}};\xi(w))$ for every word $w$, hence
$h(\gamma x^m)=h(x^m)$.

(d) Put $\mu'_j=\sum_\tau\sum_{i\ge j}x_{\tau,i}$. If $N-j\le j$ then $\mu_j=\mu'_j$. Otherwise
$\mu'_j-\mu_j=\sum_\tau p_1(\tau)$, so $\mu_j'^{k}\phi_j-\mu_j^{k}\phi_j$ is $\sum_\tau p_1(\tau)$ times a homogeneous
polynomial; by (b) (for $k=d_X+1$; the cofactor has degree $n$ by \cref{pf:volthm:lem:testpoly}) and by
\cref{pf:volthm:lem:degfl}(b) (for $k=d_X$) we may replace $\mu_j$ by $\mu'_j$ in $h(\mathrm{tp}_j)$ and in
$\mathrm{td}_j$. After the substitution $x\mapsto-x$, which multiplies monomials of degree $d$ by $(-1)^d$,
$h(\mu_j'^{d_X+1}\phi_j)$ and $\volthmdeg(\mu_j'^{d_X}\phi_j)$ become $-h^{\mathrm{Fl}}(F_A)$ and $J$, where
$h^{\mathrm{Fl}}(F_A)$ is the arithmetic height of the test class $F_A$ and $J$ the flag integral appearing in
$(\mathrm{GH})$ for the coefficient $(-1)^{d_X}$ (both identities are direct computations). Here $J=\mathrm{td}_j>0$. The hypothesis on $Z$ holds with
$e=\frac{\min(j,N-j)}2\log(N+1)$: for a complex $(N-j)\times j$ matrix $Z$ with entries of absolute value $1$,
$1+ZZ^*$ is positive definite with diagonal entries $1+j$, so Hadamard's inequality gives
$0<\det(1+ZZ^*)\le(1+j)^{N-j}$, and $\det(1+ZZ^*)=\det(1+Z^*Z)\le(1+N-j)^j$; both are $\le(N+1)^{\min(j,N-j)}$.
Hence $(\mathrm{GH})$ gives $h(\mathrm{tp}_j)=-h^{\mathrm{Fl}}(F_A)\le(d_X+1)\mathrm{td}_j|\Psi|e$.
\end{proof}

\begin{lemma}[From antisymmetric perturbations to the functional]\label{pf:volthm:lem:corb}
Let $1\le j$, $R$ a linear form on $\volthmR$, $v>0$, $\varepsilon\ge0$, $P_l$ a homogeneous polynomial of degree
$n+1$ on weights with $P_l(\lambda)=R(\volthmlin(\lambda)^{n+1})$ for all $\lambda$, and $l$ a $W_j$-invariant
linear form with $|\mathrm{coeff}_m(P_l-\Delta l)|\le\varepsilon$ for all $m$. Then there are
$\alpha,\beta\in\R^\Psi$ with
\[
\Bigl|R(x^m)+v\Bigl(\sum_\tau(\alpha_\tau-\beta_\tau)\volthmSh_j^\tau(x^m)+\sum_\tau\beta_\tau\volthmTot^\tau(x^m)
\Bigr)\Bigr|\le\varepsilon
\]
for every monomial $x^m$ of degree $n+1$.
\formallabel{GlobalTheory.CompanionTheory.Ar.Top.cor_b_poly_functional}
\formallabel{GlobalTheory.T.antisymmetric_perturbation_bound}
\end{lemma}

\begin{proof}
By \cref{pf:volthm:lem:invlin} write $l(x)=\sum_\tau(a_\tau-b_\tau)s_\tau(x)+\sum_\tau b_\tau t_\tau(x)$ with
$s_\tau=\sum_{i<j}x_{\tau,i}$, $t_\tau=\sum_ix_{\tau,i}$, and by \cref{pf:volthm:lem:weyl} let
$\volthmfd=c\Delta$ with $c>0$. Put $K=v(n+1)c$, $\alpha=a/K$, $\beta=b/K$, and
$\Phi=R+v(\sum_\tau(\alpha_\tau-\beta_\tau)\volthmSh_j^\tau+\sum_\tau\beta_\tau\volthmTot^\tau)$. By
\cref{pf:volthm:lem:degfl}(d), for every weight $x$,
\[
\Phi(\volthmlin(x)^{n+1})=P_l(x)-v(n+1)\volthmfd(x)\sum_\tau\bigl((\alpha_\tau-\beta_\tau)s_\tau(x)+\beta_\tau t_\tau(x)\bigr)
=P_l(x)-\Delta(x)l(x).
\]
So the polarisation of $\Phi$ (\cref{pf:volthm:lem:polar}) is $P_l-\Delta l$, and
$|\Phi(x^m)|\le|\mathrm{coeff}_m(P_l-\Delta l)|\le\varepsilon$.
\end{proof}

\begin{proposition}[Asymptotic formula, $\mathrm{AF}$]\label{pf:volthm:prop:ThmB}
Assume $(\mathrm{AP})$. Let $T$ be a root tower and $h$ a linear form on $\volthmR$ such that $T$
satisfies $\mathrm{Pr}$, $\mathrm{WB}$ and $\mathrm{I}_8(h)$, and $h$ satisfies $\mathrm{I}_9(h)$. Then
$T$ satisfies $\mathrm{AF}(h)$.
\formallabel{RootFieldNumerics.antisymmetry_tower_theorem}
\formallabel{GlobalTheory.CompanionTheory.Ar.Top.thmB_functional}
\formallabel{GlobalTheory.T.tower_theorem_B}
\end{proposition}

\begin{proof}
$\vol_T>0$ by $\mathrm{Pr}$. Let $h^\sharp$ be the polarisation of $h$ in degree $n+1$
(\cref{pf:volthm:lem:polar}): a homogeneous polynomial with $h^\sharp(x)=h(\volthmlin(x)^{n+1})$.

\emph{The sets $K_\tau$.} For $\tau\in\Psi$ let $x^0_\tau(\sigma,i)=-\frac1{4(N+1)}\bigl(i-\mathbf 1_{\sigma=\tau,\,
i\ge j}\bigr)$ and let $K_\tau$ be the intersection of the hyperplane $H_\tau=\{x_{\tau,j-1}=x_{\tau,j}\}$ with the
closed ball of radius $\frac1{16(N+1)}$ (sup norm) around $x^0_\tau$. Then $x^0_\tau\in H_\tau$, $K_\tau$ is compact
and contains the relatively open set $H_\tau\cap(\text{open ball})$. Every $x$ in the closed ball with
$x_{\tau,j-1}=x_{\tau,j}$ lies in $\mathcal W_j(\tau,\frac1{8(N+1)})$ and is dominant: with
$\delta=\frac1{4(N+1)}$, each entry of $x$ is within $\delta/4$ of $-\delta(i-\mathbf 1_{\sigma=\tau,i\ge j})$, so
consecutive differences $x_{\sigma,a}-x_{\sigma,a+1}$ are $\ge\delta/2$ except at $(\tau,j-1)$, where the
difference is $0$, and $|x_{\sigma,i}|\le\delta N+\delta/4\le\frac12$. In particular $h^\sharp=0$ on $K_\tau$ by
$\mathrm{I}_9(h)$. Let $C_0$ be the constant of $(\mathrm{AP})$ for $h^\sharp$ and the $K_\tau$, and
$(C_W,B_1)$ the constants of $\mathrm{WB}$ for $c_0=\frac1{8(N+1)}$.

\emph{Antisymmetry.} Let $B\ge\max(B_1,2)$ and let $P_B$ be the polarisation of $\volthmPP_B$ in degree $n+1$. Then
$P_B-\vol_Th^\sharp$ is homogeneous of degree $n+1=|\Psi|\binom N2+1$, and for every simple reflection
$s=s_{\tau,a,a+1}$ with $a+1\ne j$,
\[
(P_B-\vol_Th^\sharp)(sx)=\sum_c\Bigl(\volthmPP_{B,c}(\volthmlin(sx)^{n+1})-\vol(c)h(\volthmlin(sx)^{n+1})\Bigr)
=-(P_B-\vol_Th^\sharp)(x)
\]
by $\mathrm{I}_8(h)$; so $P_B-\vol_Th^\sharp$ lies in the antisymmetric space of $(\mathrm{AP})$.
On $K_\tau$ we have $|P_B(x)|=|\Pi(B,x)|\le C_W\log B/B$ by $\mathrm{WB}$. Hence
$(\mathrm{AP})$ (with $v=\vol_T$, $\varepsilon=C_W\log B/B$) gives a $W_j$-invariant linear form
$l$ with $|\mathrm{coeff}_m(P_B-\vol_Th^\sharp-\Delta l)|\le C_0C_W\log B/B$.

\emph{Conclusion.} Apply \cref{pf:volthm:lem:corb} to $R=\volthmPP_B-\vol_T h$, whose polarisation is
$P_B-\vol_Th^\sharp$, and $v=\vol_T$: there are $\alpha,\beta$ with
$|\volthmPP_B(x^m)-\vol_T(h(x^m)-\sum_\tau(\alpha_\tau-\beta_\tau)\volthmSh_j^\tau(x^m)-\sum_\tau\beta_\tau
\volthmTot^\tau(x^m))|\le C_0C_W\log B/B$ for all monomials of degree $n+1$. This is $\mathrm{AF}(h)$ with
constant $C_0C_W$ and threshold $\max(B_1,2)$.
\end{proof}

\begin{proposition}[Galois averaging, $\overline{\mathrm{AF}}$]\label{pf:volthm:prop:ThmBA}
If $T$ satisfies $\mathrm{AF}(h)$ and $\mathrm{GI}$, and $h$ satisfies $\mathrm{LI}(h)$, then $T$
satisfies $\overline{\mathrm{AF}}(h)$ with the same constant.
\formallabel{GlobalTheory.CompanionTheory.Ar.Top.thmBA_of_thmB_propA}
\formallabel{GlobalTheory.T.tower_theorem_B_averaged}
\end{proposition}

\begin{proof}
Let $(C',B_1)$ be the constants of $\mathrm{AF}(h)$, $\Gamma$ the group of $\mathrm{GI}$, and $B\ge\max(B_1,2)$
with parameters $\alpha,\beta$. For a monomial $x^m$ of degree $n+1$ and $\gamma\in\Gamma$, apply $\mathrm{AF}$ to
$\gamma x^m$. By $\mathrm{GI}$, $\volthmPP_B(\gamma x^m)=\volthmPP_B(x^m)$; by $\mathrm{LI}$,
$h(\gamma x^m)=h(x^m)$; and by \cref{pf:volthm:lem:degfl}(e),
$\sum_\tau\alpha_\tau\volthmSh_j^\tau(\gamma x^m)=\sum_\tau\alpha_{\gamma\tau}\volthmSh_j^\tau(x^m)$, and similarly for
$\beta$ and $\volthmTot$. So the bound of $\mathrm{AF}$ holds for $x^m$ with $(\alpha\circ\gamma,\beta\circ\gamma)$.
Since $\Gamma$ is transitive, $A=\frac1{|\Gamma|}\sum_{\gamma\in\Gamma}\alpha_{\gamma\tau}$ and
$B_\beta=\frac1{|\Gamma|}\sum_\gamma\beta_{\gamma\tau}$ do not depend on $\tau$. Averaging the expressions inside
the absolute values over $\gamma\in\Gamma$ (they are affine in $(\alpha,\beta)$) and using the triangle inequality
gives
$|\volthmPP_B(x^m)-\vol_T(h(x^m)-(A-B_\beta)\sum_\tau\volthmSh_j^\tau(x^m)-B_\beta\sum_\tau\volthmTot^\tau(x^m))|\le
C'\log B/B$.
\end{proof}

% ---------------------------------------------------------------------------
\subsection{The identity at a cutoff and the lower bounds}\label{pf:volthm:sec:lower}

Throughout this subsection $T$ is a root tower with index $j$ and $h$ a linear form on $\volthmR$. For $B_0\in\N$
and $A_0,b_0\in\R$, $\mathcal J(B_0;A_0,b_0)$ is the statement: for all $\tau$, $1\le k\le N$ and $Q$
homogeneous of degree $n+1-k$,
\begin{multline}\label{pf:volthm:eq:ident}
\volthmPP_{B_0}(p_k(\tau)Q)=\vol_T\Bigl(h(p_k(\tau)Q)-A_0\sum_\sigma\volthmSh_j^\sigma(p_k(\tau)Q)\\
-b_0\sum_\sigma\volthmTot^\sigma(p_k(\tau)Q)+R_j(B_0,\tau,k,Q)\Bigr)
\end{multline}
(here $\mathrm{frob}$ is fixed and $R_j$ is formed with the set $\mathrm{bad}$ of $T$).

\begin{lemma}[Euler factors]\label{pf:volthm:lem:deltadecomp}
Let $1\le j\le N$ and assume $\mathrm{CS}(\mathrm{frob})$. There are reals $\kappa(p)$ such that for all
primes $p$, $\tau$, $k\ge1$ and $Q$,
\begin{multline*}
\delta_j(p,\tau,\mathrm{frob}(p,\tau,\cdot),k,Q)=-\frac{\log p}{p-1}\sum_\sigma\volthmSh_j^\sigma(p_k(\tau)Q)\\
+\kappa(p)\sum_\sigma\volthmTot^\sigma(p_k(\tau)Q)+r_j(p,\mathrm{frob}(p,\tau,\cdot),k,Q).
\end{multline*}
\formallabel{GlobalTheory.CompanionTheory.Ar.Top.deltaTerm_eq}
\formallabel{GlobalTheory.T.power_sum_delta_decomposition}
\end{lemma}

\begin{proof}
If $\Psi$ is empty there is nothing to prove; fix $\tau_0\in\Psi$. Let $\mathrm{rk}(\sigma^+)=j$ and
$\mathrm{rk}(\sigma^-)=N-j$, so that $\pi_{j,0}(\sigma)=\mathrm{rk}(\sigma)$ is constant, and put
$\kappa(p)=\frac{\log p}N\sum_{m\ge0}p^{-(m+1)}\mathrm{rk}(\mathrm{frob}(p,\tau_0,m))$. By
\cref{pf:volthm:lem:degfl}(c), $\sum_\sigma\volthmSh_j^\sigma(p_k(\tau)Q)=k\volthmdeg(p^{<j}_{k-1}(\tau)Q)$ and
$\sum_\sigma\volthmTot^\sigma(p_k(\tau)Q)=\mathbf 1_{k=1}N\volthmdeg(Q)$. For $k\ge2$ the claim is the definition
\eqref{pf:volthm:eq:delta}, \eqref{pf:volthm:eq:rest}. For $k=1$, $r_j=0$, $p^{<j}_0(\tau)=j$, and
$\volthmdeg(\pi_{j,0}(\mathrm{frob}(p,\tau,m))Q)=\mathrm{rk}(\mathrm{frob}(p,\tau_0,m))\volthmdeg(Q)$ by
$\mathrm{CS}$; so
\begin{align*}
\delta_j&=\log p\sum_mp^{-(m+1)}\mathrm{rk}(\mathrm{frob}(p,\tau_0,m))\volthmdeg(Q)-\frac{\log p}{p-1}j\volthmdeg(Q)\\
&=\kappa(p)N\volthmdeg(Q)-\frac{\log p}{p-1}j\volthmdeg(Q).
\end{align*}
\end{proof}

\begin{lemma}[Summability of the remainders]\label{pf:volthm:lem:rest-summable}
For all $\tau,k,Q$ there is $c$ with $|r_j(p,\mathrm{frob}(p,\tau,\cdot),k,Q)|\le c\log p/p^2$ for all $p\ge2$, and
for every $B_0$ the partial sums $\sum_{B_0<p\le B,\,p\notin\mathrm{bad}}r_j(p,\mathrm{frob}(p,\tau,\cdot),k,Q)$ over
primes converge to $R_j(B_0,\tau,k,Q)$ as $B\to\infty$.
\formallabel{RootFieldNumerics.rest_term_bound_convergent}
\formallabel{GlobalTheory.CompanionTheory.Ar.Top.restTerm_bound_and_tendsto}
\formallabel{GlobalTheory.T.remainder_term_summable}
\formallabel{GroundTheory.Heights₂.Heights.RestTerms.RestInnerTsumBound.proof}
\end{lemma}

\begin{proof}
For $k\le1$, $r_j=0$. For $k\ge2$ let $M=\max_{\sigma\in\Psi\sqcup\Psi}|k\volthmdeg(\pi_{j,k-1}(\sigma)Q)|$ (a finite
maximum). For $p\ge2$, $\sum_{m\ge0}p^{-k(m+1)}=p^{-k}/(1-p^{-k})\le2p^{-k}$, so the sum over $m$ converges absolutely
and $|r_j|\le\log p\cdot M\cdot2p^{-k}\le2M\log p/p^2$. Since $\sum_p\log p/p^2<\infty$, the family over primes is
absolutely summable, and its partial sums converge to its sum $R_j(B_0,\tau,k,Q)$.
\end{proof}

\begin{proposition}[Identity at a cutoff]\label{pf:volthm:prop:identat}
Assume $\mathrm{CS}(\mathrm{frob})$, and that $T$ satisfies $\mathrm{Pr}$,
$\mathrm{EP}(\mathrm{frob})$ and $\overline{\mathrm{AF}}(h)$. For every $B_0\ge2$ there are $A_0,b_0\in\R$ with
$\mathcal J(B_0;A_0,b_0)$.
\formallabel{GlobalTheory.CompanionTheory.Ar.Top.exists_identAt}
\formallabel{GlobalTheory.T.tower_identity_at_cutoff}
\end{proposition}

\begin{proof}
If $\Psi$ is empty the statement is vacuous; fix $\tau_0\in\Psi$. We have $N\ge2$ and $\vol_T>0$. Let
$(C,B_1)$ and $\alpha_B,\beta_B$ ($B\ge B_1$) be given by $\overline{\mathrm{AF}}(h)$, let $\kappa$ be as in
\cref{pf:volthm:lem:deltadecomp}, and for $B\ge B_0$ let $S(B)$ be the set of primes $p\notin\mathrm{bad}$ with
$B_0<p\le B$. Put
\[
a_B=\alpha_B-\beta_B+\sum_{p\in S(B)}\frac{\log p}{p-1},\qquad b_B=\beta_B-\sum_{p\in S(B)}\kappa(p).
\]

\emph{Step 1: a finite identity.} Let $\tau$, $1\le k\le N$ and $Q$ homogeneous of degree $n+1-k$, and $F=p_k(\tau)Q$,
which is homogeneous of degree $n+1$ (note $k\le N\le n+1$, as $\Psi\ne\emptyset$). Summing \eqref{pf:volthm:eq:prop6} over $c$ gives
$\volthmPP_{B_0}(F)=\volthmPP_B(F)+\vol_T\sum_{p\in S(B)}\delta_j(p,\tau,\mathrm{frob}(p,\tau,\cdot),k,Q)$. Let
$G_B(F)=\volthmPP_B(F)-\vol_T(h(F)-(\alpha_B-\beta_B)\sum_\sigma\volthmSh_j^\sigma(F)-\beta_B\sum_\sigma\volthmTot^\sigma(F))$;
by $\overline{\mathrm{AF}}$ and linearity, $|G_B(F)|\le(\sum_m|\mathrm{coeff}_m F|)\,C\log B/B\to0$. Inserting
\cref{pf:volthm:lem:deltadecomp},
\begin{multline}\label{pf:volthm:eq:finite}
\volthmPP_{B_0}(F)=\vol_T\Bigl(h(F)-a_B\sum_\sigma\volthmSh_j^\sigma(F)-b_B\sum_\sigma\volthmTot^\sigma(F)\\
+\sum_{p\in S(B)}r_j(p,\mathrm{frob}(p,\tau,\cdot),k,Q)\Bigr)+G_B(F).
\end{multline}

\emph{Step 2: convergence of $a_B,b_B$.} By \cref{pf:volthm:lem:degfl}(f) there are $Q_1$ of degree $n$ with
$\volthmdeg(Q_1)=1$ and $Q_2$ of degree $n-1$ with $\volthmdeg(p^{<j}_1(\tau_0)Q_2)=1$. By
\cref{pf:volthm:lem:degfl}(c), $F_1=p_1(\tau_0)Q_1$ has $\sum_\sigma\volthmTot^\sigma(F_1)=N\ne0$, and $F_2=p_2(\tau_0)Q_2$
has $\sum_\sigma\volthmSh_j^\sigma(F_2)=2\ne0$ and $\sum_\sigma\volthmTot^\sigma(F_2)=0$. In \eqref{pf:volthm:eq:finite} for
$F_2$ everything except $a_B$ converges as $B\to\infty$ (the prime sum by \cref{pf:volthm:lem:rest-summable},
$G_B\to0$), so $a_B$ converges, to $A_0$ say; then \eqref{pf:volthm:eq:finite} for $F_1$ shows that $b_B$
converges, to $b_0$ say.

\emph{Step 3.} Letting $B\to\infty$ in \eqref{pf:volthm:eq:finite} for arbitrary $(\tau,k,Q)$, using
\cref{pf:volthm:lem:rest-summable}, gives \eqref{pf:volthm:eq:ident}.
\end{proof}

\begin{lemma}[Expansion of the identity]\label{pf:volthm:lem:identexp}
If $\mathcal J(B_0;A_0,b_0)$ holds and $(Q_{\tau,k})$ is a power-sum decomposition of $F$, then
\[
\volthmPP_{B_0}(F)=\vol_T\Bigl(h(F)-A_0\sum_\sigma\volthmSh_j^\sigma(F)-b_0\sum_\sigma\volthmTot^\sigma(F)+
\sum_\tau\sum_{k=1}^NR_j(B_0,\tau,k,Q_{\tau,k})\Bigr).
\]
\formallabel{RootFieldNumerics.tower_point_count_identity}
\formallabel{GlobalTheory.CompanionTheory.Ar.Top.identOnDecomposedForm}
\formallabel{GlobalTheory.T.tower_height_identity_expansion}
\end{lemma}

\begin{proof}
Sum \eqref{pf:volthm:eq:ident} over $(\tau,k)$; all functionals involved are linear.
\end{proof}

\begin{lemma}[The flat coefficient bound]\label{pf:volthm:lem:flat}
Let $\Psi$ be nonempty. Assume $T$ satisfies $\mathrm{Pr}$ and $\mathrm{FL}(C',r,L)$, $h$ satisfies
$\mathrm{I}_9(h)$, $B_0\ge128$, $|\mathrm{bad}|\le B_0$ and $\mathcal J(B_0;A_0,b_0)$. Then
$|jA_0+Nb_0|\le C'N\Lambda(B_0)$.
\formallabel{GlobalTheory.CompanionTheory.Ar.Top.flat_85_bound}
\formallabel{GlobalTheory.T.cutoff_coefficients_bound}
\end{lemma}

\begin{proof}
Fix $\tau$ and $Q_1$ with $\volthmdeg(Q_1)=1$ (\cref{pf:volthm:lem:degfl}(f)), and put $F=p_1(\tau)Q_1$. By
\cref{pf:volthm:lem:degfl}(c), $\volthmSh_j^\sigma(F)=\mathbf 1_{\sigma=\tau}j$ (as $p^{<j}_0(\tau)=j$) and
$\volthmTot^\sigma(F)=\mathbf 1_{\sigma=\tau}N$; $h(F)=0$ by $\mathrm{I}_9(h)$; $R_j(B_0,\tau,1,Q_1)=0$ since $r_j=0$
for $k=1$; and $\volthmPP_{B_0}(F)=\sum_ch(B_0,c,\tau)\vol(c)$ by $\mathrm{FL}$. Hence \eqref{pf:volthm:eq:ident}
reads $\sum_ch(B_0,c,\tau)\vol(c)=-\vol_T(jA_0+Nb_0)$, and $|jA_0+Nb_0|\le\max_c|h(B_0,c,\tau)|\le C'N\Lambda(B_0)$.
\end{proof}

\begin{lemma}[Prime sums]\label{pf:volthm:lem:primesums}
\begin{enumerate}[label=(\alph*)]
\item For real $b>0$ and every $K$, $\sum_{p<K,\ p>b}\log p/p^2\le2\log4/b$ (sum over primes).
\item Let $k\ge2$, $\tau$, $Q$, and $M\ge0$ with $|k\volthmdeg(\pi_{j,k-1}(\sigma)Q)|\le M$ for all $\sigma\in\Psi\sqcup
  \Psi$. Then $|r_j(p,\mathrm{fr},k,Q)|\le2M\log p/p^k$ for every prime $p$ and every $\mathrm{fr}$, and for
  $0<b<B+1$: $|R_j(B,\tau,k,Q)|\le5.56\,M/b^{k-1}$.
\end{enumerate}
\formallabel{GroundTheory.Heights₂.Heights.RestTerms.Elem.prime_partial_le}
\formallabel{RootFieldNumerics.rest_term_prime_bound}
\formallabel{RootFieldNumerics.rest_sum_primes_bound}
\end{lemma}

\begin{proof}
(a) By Chebyshev's bound $\prod_{p\le x}p\le4^x$ (e.g.\ \cite[Thm.~415]{HardyWright2008}),
$\theta(x)=\sum_{p\le x,\,p>b}\log p\le x\log4$. Let $M_0=\lfloor b\rfloor+1>b$; terms with $p<M_0$ vanish. Abel
summation over $M_0\le p<K$ gives
$\sum\log p/p^2\le2\log4(\frac1{M_0}-\frac1K)+\frac{\log4}K\le\frac{2\log4}{M_0}\le\frac{2\log4}b$.
(b) The first bound is the estimate in the proof of \cref{pf:volthm:lem:rest-summable}. For $p>b$, $p\ge2$:
$2M\log p/p^k\le\frac{2M}{b^{k-2}}\cdot\frac{\log p}{p^2}$. The primes in $R_j(B,\cdot)$ satisfy $p\ge B+1>b$, so by (a)
$|R_j|\le\frac{2M}{b^{k-2}}\cdot\frac{2\log4}b\le\frac{5.56M}{b^{k-1}}$, as $4\log4\le5.56$.
\end{proof}

\begin{proposition}[Remainder for linear forms]\label{pf:volthm:prop:remlin}
Let $1\le j<N$, $\mathrm{WH}_j(u,g)$, $B_0(u,g)<B+1$, and $(Q_{\tau,k})$ a power-sum decomposition of
$\volthmlin(u)^{n+1}$. Then
\[
\Bigl|\sum_\tau\sum_{k=1}^NR_j(B,\tau,k,Q_{\tau,k})\Bigr|\le2(n+1)\volthmfd(u)\min(j,N-j)\,\mathrm{spread}(u).
\]
\formallabel{RootFieldNumerics.rest_sum_total_bound}
\formallabel{GlobalTheory.CompanionTheory.Ar.Top.restSum_total_bound}
\formallabel{GlobalTheory.T.lin_root_remainder_bound}
\end{proposition}

\begin{proof}
Put $E=(n+1)\volthmfd(u)>0$, $m_j=\min(j,N-j)$, $U_\tau=\max_iu_{\tau,i}-\min_iu_{\tau,i}$ and
$A=N(1+2\|u\|/g)$. The key input is the single-term bound
\lref{GlobalTheory.T.lie_term_flag_bound} (a result of the development, whose proof reduces
$\pi_{j,k-1}(\sigma)Q_{\tau,k}$ modulo the coinvariant relations to at most $m_j$ terms and bounds each by a divided
difference estimate for $u$; we do not reproduce it): for $2\le k\le N$ and all $\sigma$,
\[
|k\volthmdeg(\pi_{j,k-1}(\sigma)Q_{\tau,k})|\le E\,m_j\,\frac{2^{k-1}}{k!}\,U_\tau\,A^{k-1}.
\]
By \cref{pf:volthm:lem:primesums}(b) with $b=B_0(u,g)$ (note $0<b<B+1$),
$|R_j(B,\tau,k,Q_{\tau,k})|\le5.56\,Em_jU_\tau\,(2A/b)^{k-1}/k!$. Since $|\Psi|\ge1$ and $N\ge2$,
\[
2A=2N+4N\|u\|/g\le\tfrac12\cdot4|\Psi|N^2(1+\|u\|/g)=b/2,
\]
so $(2A/b)^{k-1}\le(1/2)^{k-1}$. For $k=1$ the
remainder is $0$. Using $2\cdot3^m\le(m+2)!$ one gets $\sum_{k=2}^N(1/2)^{k-1}/k!\le0.3$ by induction on $N$. Hence
$|\sum_kR_j(B,\tau,k,Q_{\tau,k})|\le5.56\cdot0.3\,Em_jU_\tau\le2Em_jU_\tau$; sum over $\tau$.
\end{proof}

The corresponding bound for the test polynomial is proved in \cref{pf:volthm:sec:remainder}; it rests on a
combinatorial bound for the individual terms (\cref{pf:volthm:prop:factor}).

\begin{proposition}[Remainder for the test polynomial]\label{pf:volthm:prop:remtest}
Let $1\le j<N$, $4|\Psi|N^2\le B_0$, and $(Q_{\tau,k})$ a power-sum decomposition of $\mathrm{tp}_j$. Then, for every
$\mathrm{bad}$ and $\mathrm{frob}$,
$\bigl|\sum_\tau\sum_{k=1}^NR_j(B_0,\tau,k,Q_{\tau,k})\bigr|\le1.6\,(d_X+1)\,\mathrm{td}_j$.
\formallabel{GlobalTheory.T.test_polynomial_remainder_bound}
\formallabel{GlobalTheory.CompanionTheory.Ar.Top.testRest_bound}
\formallabel{GlobalTheory.CompanionTheory.Ar.Top.TestRest.finite_bound}
\end{proposition}

\begin{proposition}[Lower bound at the cutoff]\label{pf:volthm:prop:lowerbase}
Let $\Psi$ be nonempty, $C,C'\ge0$ and $r\ge1$. Assume $T$ satisfies $\mathrm{Pr}$,
$\mathrm{FL}(C',r,L)$ and $\mathrm{SL}(C,r,L)$, and $h$ satisfies $\mathrm{I}_9(h)$ and the bound of
\cref{pf:volthm:prop:hh}(d). If $B_0\ge128$, $4|\Psi|N^2\le B_0$, $|\mathrm{bad}|\le B_0$ and
$\mathcal J(B_0;A_0,b_0)$ holds, then $A_0\ge-(2(C+C')+3)\Lambda(B_0)$.
\formallabel{GlobalTheory.CompanionTheory.Ar.Top.test_height_bound_slope_flat}
\formallabel{GlobalTheory.T.test_height_bound}
\end{proposition}

\begin{proof}
Write $f=|\Psi|$, $D=d_X$, $\mathrm{td}=\mathrm{td}_j>0$, $K=(D+1)\mathrm{td}>0$, $m=\min(j,N-j)\ge1$,
$\Lambda=\Lambda(B_0)$ and $\ell=\log(N+1)$. Decompose $\mathrm{tp}_j$ (\cref{pf:volthm:lem:testpoly,pf:volthm:lem:decomp})
and let $R$ be the corresponding remainder sum, $|R|\le1.6K$ by \cref{pf:volthm:prop:remtest}. By
\cref{pf:volthm:lem:identexp}, $\mathrm{SL}$ and $\vol_T>0$,
\begin{equation}\label{pf:volthm:eq:test}
-Cfm\Lambda K\le h(\mathrm{tp}_j)-A_0\sum_\sigma\volthmSh_j^\sigma(\mathrm{tp}_j)-b_0\sum_\sigma\volthmTot^\sigma(\mathrm{tp}_j)+R,
\end{equation}
and $h(\mathrm{tp}_j)\le Kf\frac m2\ell$. We use the elementary fact: if $K>0$, $f,m\ge1$,
$-Cfm\Lambda K\le h-xfKm+R$, $h\le Kf\frac m2\ell$ and $|R|\le1.6K$, then $x\le C\Lambda+\ell/2+1.6$ (indeed
$xfKm\le h+R+Cfm\Lambda K\le(C\Lambda+\ell/2+1.6)fKm$, as $K\le fmK$).

\emph{Case $N-j\le j$} ($m=N-j$). By \cref{pf:volthm:lem:testpoly}, $\sum_\sigma\volthmSh_j^\sigma(\mathrm{tp}_j)=0$ and
$\sum_\sigma\volthmTot^\sigma(\mathrm{tp}_j)=fKm$, so \eqref{pf:volthm:eq:test} gives $b_0\le C\Lambda+\ell/2+1.6$. With
\cref{pf:volthm:lem:flat}, $jA_0\ge-C'N\Lambda-Nb_0\ge-N(C\Lambda+\ell/2+1.6+C'\Lambda)$, and $N\le2j$, so
$A_0\ge-2(C\Lambda+\ell/2+1.6+C'\Lambda)$.

\emph{Case $N-j>j$} ($m=j$). Now both sums equal $-fKm$, so \eqref{pf:volthm:eq:test} gives
$-(A_0+b_0)\le C\Lambda+\ell/2+1.6$. Then $(N-j)A_0=N(A_0+b_0)-(jA_0+Nb_0)\ge-N(C\Lambda+\ell/2+1.6)-C'N\Lambda$ and
$N\le2(N-j)$, giving the same bound.

Finally $N+1\le B_0$ and $r\ge1$ give $\ell\le\log B_0\le\Lambda$, and $1.6\le\log128\le\Lambda$; hence
$2(C\Lambda+\ell/2+1.6+C'\Lambda)\le(2(C+C')+3)\Lambda$.
\end{proof}

\begin{proposition}[$\mathrm{T}_2$]\label{pf:volthm:prop:T2}
Let $K,C'\ge0$ and $r\ge1$. Assume $\mathrm{CS}(\mathrm{frob})$, that $T$ satisfies $\mathrm{Pr}$,
$\mathrm{EP}(\mathrm{frob})$, $\overline{\mathrm{AF}}(h)$ and $\mathrm{FL}(C',r,L)$, that $h$ satisfies $\mathrm{I}_9(h)$,
and that $A_0\ge-K\Lambda(B_0)$ whenever $B_0\ge128$, $4|\Psi|N^2\le B_0$, $|\mathrm{bad}|\le B_0$ and
$\mathcal J(B_0;A_0,b_0)$. Then $T$ satisfies $\mathrm{T}_2(K+C'+2,r,L)$.
\formallabel{GlobalTheory.CompanionTheory.Ar.Top.tower_T2_of_inputs}
\formallabel{GlobalTheory.T.tower_T2}
\end{proposition}

\begin{proof}
Let $\mathrm{WH}_j(u,g)$, $B\ge128$, $B_0(u,g)<B+1$ and $|\mathrm{bad}|\le B$. If $\Psi$ is empty, $\volthmlin(u)=0$
and both sides of $\mathrm{T}_2$ vanish; assume $\Psi\ne\emptyset$. Then $4|\Psi|N^2\le B_0(u,g)<B+1$, so
$4|\Psi|N^2\le B$. By \cref{pf:volthm:prop:identat} there are $A_0,b_0$ with $\mathcal J(B;A_0,b_0)$; by
hypothesis $A_0\ge-K\Lambda$ ($\Lambda=\Lambda(B)\ge0$), and by \cref{pf:volthm:lem:flat}
$|jA_0+Nb_0|\le C'N\Lambda$. Let $E=(n+1)\volthmfd(u)>0$ (\cref{pf:volthm:lem:weyl}). Decomposing
$\volthmlin(u)^{n+1}$ (\cref{pf:volthm:lem:decomp}) and applying \cref{pf:volthm:lem:identexp},
\[
\Pi(B,u)=\vol_T\Bigl(h(\volthmlin(u)^{n+1})+E\sum_\sigma\bigl(A_0s_\sigma+b_0t_\sigma\bigr)+R\Bigr),\qquad
s_\sigma=\sum_{i<j}u_{\sigma,i},\ t_\sigma=\sum_iu_{\sigma,i},
\]
by \cref{pf:volthm:lem:degfl}(d), where $|R|\le2E\,m_j\,\mathrm{spread}(u)$ by \cref{pf:volthm:prop:remlin}
($m_j=\min(j,N-j)$). By $\mathrm{I}_9(h)$, $h(\volthmlin(u)^{n+1})\ge0$. For each $\sigma$ write
$A_0s_\sigma+b_0t_\sigma=A_0(s_\sigma-\frac jNt_\sigma)+(jA_0+Nb_0)\frac{t_\sigma}N$. By \cref{pf:volthm:lem:partial}
(as $u_\sigma$ is decreasing), $0\le s_\sigma-\frac jNt_\sigma\le\|u_\sigma\|_1$ and $|t_\sigma|\le\|u_\sigma\|_1$, so
$A_0s_\sigma+b_0t_\sigma\ge-(K+C')\Lambda\|u_\sigma\|_1$, and summing,
$\sum_\sigma(A_0s_\sigma+b_0t_\sigma)\ge-(K+C')\Lambda\|u\|_1$. Therefore
\begin{multline*}
\Pi(B,u)\ge-\vol_TE\bigl((K+C')\Lambda\|u\|_1+2m_j\mathrm{spread}(u)\bigr)\\
\ge-(K+C'+2)(n+1)\vol_T\volthmfd(u)\,\mathrm{br}(B,u).\qedhere
\end{multline*}
\end{proof}

\begin{proposition}[Mean slopes]\label{pf:volthm:prop:meanslope}
If $T$ satisfies $\mathrm{Pr}$, $\mathrm{EV}$ and $\mathrm{T}_2(C,r,L)$, it satisfies
$\mathrm{MS}(C,r,L)$.
\formallabel{RootFieldNumerics.mean_slope_of_proportional_volume}
\formallabel{GlobalTheory.CompanionTheory.Ar.Top.mean_slope_of_volume_eq}
\formallabel{GlobalTheory.T.mean_slope_proportional_volume}
\end{proposition}

\begin{proof}
By $\mathrm{EV}$ all $\vol(c)$ equal some $v>0$, so $\vol_T=|\mathrm{Comp}|v$ and
$\frac1{|\mathrm{Comp}|}\sum_c\mathrm{slope}(B,c,u)=\Pi(B,u)/((n+1)\vol_T\volthmfd(u))$. Under the hypotheses of
$\mathrm{MS}$, $u$ is regular dominant, so the denominator is positive, and $\mathrm{T}_2$ gives the claim.
\end{proof}

\begin{proof}[Proof of \cref{pf:volthm:thm:tower}]
Since $1\le j<N$, $N\ge2$. Write $h=\volthmhh^{\volthmII}$.
\begin{enumerate}[label=(\arabic*)]
\item $\mathrm{I}_3$ holds by \cref{pf:volthm:prop:I3}, and $\mathrm{T}_1$ by \cref{pf:volthm:prop:T1}.
\item $\mathrm{WB}$ holds by \cref{pf:volthm:prop:ThmW}, using the component $(\mathrm{K})$ of
  $(\mathrm{AI})$ and \cref{pf:volthm:lem:aux}.
\item By \cref{pf:volthm:prop:hh}, $h$ satisfies $\mathrm{I}_9(h)$ and $\mathrm{LI}(h)$, and the bound of
  \cref{pf:volthm:prop:hh}(d) holds for $j$.
\item $\mathrm{AF}(h)$ holds by \cref{pf:volthm:prop:ThmB}, and $\overline{\mathrm{AF}}(h)$ by \cref{pf:volthm:prop:ThmBA}.
\item By \cref{pf:volthm:prop:lowerbase} the hypothesis of \cref{pf:volthm:prop:T2} holds with
  $K=2(C+C')+3\ge0$, so $\mathrm{T}_2(K+C'+2,r,L)$ holds, and $K+C'+2=2C+3C'+5$.
\item $\mathrm{MS}(2C+3C'+5,r,L)$ follows by \cref{pf:volthm:prop:meanslope}.\qedhere
\end{enumerate}
\end{proof}

% ---------------------------------------------------------------------------
\subsection{The remainder for the test polynomial}\label{pf:volthm:sec:remainder}

In this subsection we prove \cref{pf:volthm:prop:remtest}. Throughout, $f=|\Psi|\ge1$, $1\le j<N$, and the
variables are $x_{\tau,i}$, $(\tau,i)\in\Psi\times[N]$. The \emph{small block} is $B_s=\{i\ge j\}$ with sign
$\varsigma=1$ if $N-j\le j$, and $B_s=\{i<j\}$ with $\varsigma=-1$ otherwise; it has $r=\min(j,N-j)$ elements, and
$d_X=fr(N-r)$. The \emph{small power sums} are $\pi^s_m(\tau)=\sum_{i\in B_s}(\varsigma x_{\tau,i})^m$; thus
$\mu_j=\sum_\tau\pi^s_1(\tau)$. In this subsection a \emph{decomposition} of a polynomial $F$ is a family
$(Q_{k,\tau})_{1\le k\le N,\,\tau\in\Psi}$ of polynomials, not necessarily homogeneous, with
$F=\sum_k\sum_\tau p_k(\tau)Q_{k,\tau}$; a power-sum decomposition in the sense of \cref{pf:volthm:sec:poly} is a
decomposition by homogeneous polynomials of the appropriate degrees.

\emph{The rational flag integral.} Let $V=\prod_\tau\prod_{i<i'}(x_{\tau,i}-x_{\tau,i'})$ and, for a polynomial $\Phi$ with
real coefficients,
\[
\textstyle\int\Phi=\sum_\alpha\mathrm{coeff}_\alpha(\Phi)\,\mathrm{coeff}_\alpha(V)\,\alpha!,\qquad
\alpha!=\prod_{\tau,i}\alpha(\tau,i)!,
\]
the sum over the exponents $\alpha$ of the monomials of $V$; thus $\int\Phi$ is the constant term of $V(\partial)\Phi$.
It is rational on rational polynomials. Let $\phi'$ be the product over
$\tau$ of the Vandermonde products $\prod_{i<i'}(y_i-y_{i'})$ of the variables $\varsigma x_{\tau,i}$, $i\in B_s$, and
of the other variables of $\tau$ (each in increasing order of $i$). Then $\phi'=\pm\phi_j$. Put
\[
F_s=\mu_j^{d_X+1}\phi',\qquad J=\textstyle\int\mu_j^{d_X}\phi'.
\]
For a decomposition $(Q_{k,\tau})$ of $F_s$ with rational coefficients, maps $s\colon\N\times\N\times
\Psi\to\Psi$ and signs $e(p,n,\tau,k)\in\{\pm1\}$, the \emph{prime model} at an integer $p$ is
\begin{equation}\label{pf:volthm:eq:primemodel}
R_p=\log p\sum_\tau\sum_{k=2}^N\sum_{n\ge0}p^{-(n+1)k}\,e(p,n+1,\tau,k)\;k\!\int\pi^s_{k-1}\bigl(s(p,n+1,\tau)\bigr)\,
Q_{k,\tau}.
\end{equation}

The factors of the prime model are controlled by the following bound, which is the combinatorial core of the
remainder estimate. Its proof in the development, which we do not reproduce, expresses these terms through Schur
polynomials of the small block (by divided differences), evaluates the coefficients with the Frobenius character
formula and counts of standard tableaux, and bounds the characters of symmetric groups that occur.

\begin{proposition}[Factor bound; \lref{GroundTheory.Heights₂.Heights.RestTerms.SameFactorBound.full},
\lref{GroundTheory.Heights₂.Heights.RestTerms.OtherFactorBound.full}]\label{pf:volthm:prop:factor}
$J\ne0$; $F_s$ has a decomposition with rational coefficients; and for every such decomposition, all
$2\le k\le N$ and all $\tau,\tau'\in\Psi$,
\[
\Bigl|k\int\pi^s_{k-1}(\tau')\,Q_{k,\tau}\Bigr|\le(d_X+1)\,|J|\,\frac{X^{k-1}}{k!},\qquad X=2\bigl(r(N-r)+N\bigr).
\]
\end{proposition}

\begin{lemma}[The prime model]\label{pf:volthm:lem:primemodel}
Let $B_0\ge4fN^2$ and let $(Q_{k,\tau})$, $s$, $e$ be as above.
\begin{enumerate}[label=(\alph*)]
\item For every prime $p>B_0$, $|R_p|\le1.5\,fX\,\frac{\log p}{p^2}\,(d_X+1)|J|$.
\item For every set $P$ of primes $>B_0$, the family $(R_p)_{p\in P}$ is absolutely summable and
$\bigl|\sum_{p\in P}R_p\bigr|\le1.6\,(d_X+1)|J|$.
\end{enumerate}
\formallabel{GroundTheory.Heights₂.Heights.RestTerms.RestPrimeBound.full}
\formallabel{GroundTheory.Heights₂.Heights.RestTerms.RestPrimeBound.proof}
\formallabel{GroundTheory.Heights₂.Heights.RestTerms.RestPrimeBoundOfFactor.full}
\formallabel{GroundTheory.Heights₂.Heights.RestTerms.RestPrimeBoundOfFactor.proof}
\formallabel{GroundTheory.Heights₂.Heights.RestTerms.RestInnerTsumBound.full}
\formallabel{GroundTheory.Heights₂.Heights.RestTerms.RestKSumBound.full}
\formallabel{GroundTheory.Heights₂.Heights.RestTerms.GeometricTail.intersectionNumber_baseChange}
\formallabel{GroundTheory.Heights₂.Heights.RestTerms.ExponentialTail.intersectionNumber_baseChange}
\formallabel{GroundTheory.Heights₂.Heights.RestTerms.XOverPrime.intersectionNumber_baseChange}
\formallabel{GroundTheory.Heights₂.Heights.RestTerms.RestBound.full}
\formallabel{GroundTheory.Heights₂.Heights.RestTerms.RestBound.proof}
\formallabel{GroundTheory.Heights₂.Heights.RestTerms.RestBoundOfPrimeBound.full}
\formallabel{GroundTheory.Heights₂.Heights.RestTerms.RestBoundOfPrimeBound.proof}
\end{lemma}

\begin{proof}
(a) Put $Y=(d_X+1)|J|$. As $r(N-r)\le N^2/4$ and $N\ge2$, $X\le N^2/2+2N$, and $p>B_0\ge4fN^2\ge4N^2$ gives
$X/p\le\frac18+\frac1{2N}\le\frac38$. For fixed $\tau,k$ the coefficients $c_n=e(p,n+1,\tau,k)\,k\int\pi^s_{k-1}(s(p,n+1,\tau))
Q_{k,\tau}$ satisfy $|c_n|\le YX^{k-1}/k!$ by \cref{pf:volthm:prop:factor}, and $\sum_{n\ge0}p^{-(n+1)k}=
p^{-k}/(1-p^{-k})\le2p^{-k}$; so the series over $n$ converges absolutely and is bounded by $2p^{-k}YX^{k-1}/k!$. Summing
over $k\ge2$, with $u=X/p\le\frac38$ and $1/(k+2)!\le1/(2\,k!)$,
\[
\sum_{k\ge2}\frac{2X^{k-1}}{p^kk!}=\frac{2X}{p^2}\sum_{k\ge0}\frac{u^k}{(k+2)!}\le\frac{X}{p^2}e^{3/8}\le
\frac{1.5\,X}{p^2},
\]
and summing over the $f$ elements $\tau$ gives (a).

(b) By (a) and \cref{pf:volthm:lem:primesums}(a), the partial sums of $|R_p|$ over $p\in P$ are bounded by
$1.5fXY\cdot2\log4/B_0$, so the family is absolutely summable with the same bound for its sum. Finally
$1.5fX\cdot2\log4/B_0\le3\log4\,X/(4N^2)\le3\log4\,(\frac18+\frac1{2N})\le\frac98\log4\le1.6$.
\end{proof}

\emph{Flag integrals.} A linear form $I$ on $\volthmR$ is a \emph{flag integral} if $I(p_m(\tau)G)=0$ whenever
$1\le m\le n$ and $G$ is homogeneous of degree $n-m$, and $I(G)\ne0$ for some $G$ homogeneous of degree $n$. By
\cref{pf:volthm:lem:degfl}(a),(b), $\volthmdeg$ is a flag integral (it is nonzero on the monomial
$x^{m_{\mathrm{id}}}$). We use two facts about them, which are results of the development that we do not reproduce.
\begin{itemize}
\item \emph{Proportionality}: for every flag integral $I$ there is $\kappa\ne0$ with $I(G)=\kappa\int G$ for all $G$
homogeneous of degree $n$. This combines
\lref{GroundTheory.Heights₂.Heights.RestTerms.FlagIntegralProportional.intersectionNumber_baseChange}, which
states that every flag integral agrees in degree $n$ with a nonzero multiple of a fixed linear form (the coefficient of
$x^{m_{\mathrm{id}}}$ after antisymmetrisation in every block), and
\lref{GroundTheory.Heights₂.Heights.RestTerms.flagIntegralReal_eq_mul_ddConst}, which states that $\int$ is a
positive multiple of that form.
\item \emph{Transfer} (\lref{GroundTheory.Heights₂.Heights.RestTerms.RestBoundEq.flagIntegralReal_transfer_of_syzygy},
which states this): let $(Q^0_{k,\tau})$ be a decomposition of a rational polynomial $F$ with rational coefficients and
$(Q_{k,\tau})$ one with real coefficients, and let $G$ be a polynomial such that $\int G\,S_{m,\tau}=0$ for every decomposition $(S_{k,\tau})$ of $0$
with rational coefficients; then $\int G\,Q_{m,\tau}=\int G\,Q^0_{m,\tau}$. (The real decompositions of $0$ are real
combinations of rational ones.)
\end{itemize}

For $\sigma\in\Psi\sqcup\Psi$ let $\bar\sigma\in\Psi$ be its label, and put $\epsilon(\sigma,m)=(-1)^m$ if $\sigma$ lies in
the copy of $\Psi$ attached to the small block (the first copy if $B_s=\{i<j\}$, the second if $B_s=\{i\ge j\}$), and
$\epsilon(\sigma,m)=-1$ otherwise.

\begin{lemma}[Lie classes]\label{pf:volthm:lem:lieclass}
Let $I$ be a flag integral, $1\le m<N$, $\sigma\in\Psi\sqcup\Psi$ and $G$ homogeneous of degree $n-m$. Then
$I(\pi_{j,m}(\sigma)G)=\epsilon(\sigma,m)\,I(\pi^s_m(\bar\sigma)G)$.
\formallabel{GroundTheory.Heights₂.Heights.RestTerms.RestBoundEq.lieClass_reduce}
\end{lemma}

\begin{proof}
By definition $\pi_{j,m}(\tau^+)=\sum_{i<j}x_{\tau,i}^m$ and $\pi_{j,m}(\tau^-)=(-1)^m\sum_{i\ge j}x_{\tau,i}^m$, while
$\sum_{i<j}x_{\tau,i}^m+\sum_{i\ge j}x_{\tau,i}^m=p_m(\tau)$ and $I(p_m(\tau)G)=0$ (note $m\le N-1\le n$). If $\sigma$ lies in
the copy of the small block, $\pi_{j,m}(\sigma)=(-1)^m\pi^s_m(\bar\sigma)$ (for $\varsigma=-1$ because
$(-x)^m=(-1)^mx^m$); otherwise $\pi_{j,m}(\sigma)\equiv-\pi^s_m(\bar\sigma)$ modulo $p_m(\bar\sigma)$, in both cases of
$\varsigma$.
\end{proof}

\begin{lemma}[Decompositions of zero]\label{pf:volthm:lem:syzygy}
Let $(S_{k,\tau})$ be a decomposition of $0$ with rational coefficients, $2\le m\le N$ and $\tau,\tau'\in\Psi$. Then
$\int\pi^s_{m-1}(\tau')\,S_{m,\tau}=0$.
\formallabel{GroundTheory.Heights₂.Heights.RestTerms.RestBoundEq.syzygy_term_zero}
\end{lemma}

\begin{proof}
Let $B_0=4fN^2$, let $p>B_0$ be a prime, $(Q^0_{k,\tau})$ a rational decomposition of $F_s$
(\cref{pf:volthm:prop:factor}), and for $t\in\Q$ let $Q^t=Q^0+tS$, again a decomposition of $F_s$. Take $s$ constant equal to
$\tau'$, $e_1\equiv1$, and $e_2$ equal to $1$ except $e_2(p,1,\tau,m)=-1$. In the difference of the prime models
\eqref{pf:volthm:eq:primemodel} for $(Q^t,s,e_1)$ and $(Q^t,s,e_2)$ only the term $n=0$, $(\tau,k)=(\tau,m)$ survives, so
it equals $2\,p^{-m}\log p\cdot m\int\pi^s_{m-1}(\tau')Q^t_{m,\tau}$, an affine function of $t$ with slope
$2p^{-m}\log p\cdot m\int\pi^s_{m-1}(\tau')S_{m,\tau}$. By \cref{pf:volthm:lem:primemodel}(b) applied to $P=\{p\}$, both
prime models are bounded by $1.6(d_X+1)|J|$ for all $t$, so the slope vanishes.
\end{proof}

\begin{proposition}[The remainder bound for flag integrals]\label{pf:volthm:prop:restbound}
Let $I$ be a flag integral, $v>0$, $B_0\ge4fN^2$, $(Q_{\tau,k})$ a power-sum decomposition of $F_s$ with real
coefficients, $\sigma\colon\N\times\N\times\Psi\to\Psi\sqcup\Psi$, and for an integer $p$ put
\[
R^I_p=v\log p\sum_\tau\sum_{k=2}^N\sum_{n\ge0}p^{-(n+1)k}\,k\,I\bigl(\pi_{j,k-1}(\sigma(p,n+1,\tau))\,Q_{\tau,k}\bigr).
\]
Then for every set $P$ of primes $>B_0$ the family $(R^I_p)_{p\in P}$ is summable and
$\bigl|\sum_{p\in P}R^I_p\bigr|\le1.6\,(d_X+1)\,v\,\bigl|I(\mu_j^{d_X}\phi')\bigr|$.
\formallabel{GroundTheory.Heights₂.Heights.RestTerms.RestBound.imp}
\formallabel{GroundTheory.Heights₂.Heights.RestTerms.RestBoundEq.core}
\end{proposition}

\begin{proof}
Let $\kappa\ne0$ be the constant of proportionality for $I$, $(Q^0_{k,\tau})$ a rational decomposition of $F_s$, and
$(Q^h_{k,\tau})$ its homogeneous parts of the right degrees (still a decomposition, since $F_s$ is homogeneous of
degree $n+1$). For $2\le k\le N$ the polynomial $\pi_{j,k-1}(\sigma)Q_{\tau,k}$ is homogeneous of degree $n$, so by
\cref{pf:volthm:lem:lieclass}, proportionality, the transfer and \cref{pf:volthm:lem:syzygy},
\[
k\,I\bigl(\pi_{j,k-1}(\sigma)Q_{\tau,k}\bigr)=\kappa\,\epsilon(\sigma,k-1)\,k\int\pi^s_{k-1}(\bar\sigma)\,Q^h_{k,\tau}.
\]
Hence $R^I_p=v\kappa R_p$ for the prime model of $(Q^h,s,e)$ with $s(p,n,\tau)=\overline{\sigma(p,n,\tau)}$ and
$e(p,n,\tau,k)=\epsilon(\sigma(p,n,\tau),k-1)$, and $I(\mu_j^{d_X}\phi')=\kappa J$. The claim follows from
\cref{pf:volthm:lem:primemodel}(b).
\end{proof}

\begin{proof}[Proof of \cref{pf:volthm:prop:remtest}]
If $\Psi$ is empty, the left side is $0$ and the right side is positive (\cref{pf:volthm:lem:testpoly}). Otherwise
$R_j(B_0,\tau,k,Q_{\tau,k})$ is the sum of the absolutely summable family $r_j(p,\mathrm{frob}(p,\tau,\cdot),k,Q_{\tau,k})$
over the primes $p>B_0$, $p\notin\mathrm{bad}$ (\cref{pf:volthm:lem:rest-summable}); it is the limit of the sums over
the finite sets $S$ of such primes with $p\le B$. It therefore suffices to show
\[
\Bigl|\sum_{p\in S}\sum_\tau\sum_{k=1}^Nr_j\bigl(p,\mathrm{frob}(p,\tau,\cdot),k,Q_{\tau,k}\bigr)\Bigr|\le1.6\,(d_X+1)\,\mathrm{td}_j
\]
for every finite set $S$ of primes $>B_0$. Comparing the Vandermonde products, $\mathrm{tp}_j=\epsilon F_s$ and
$\mu_j^{d_X}\phi_j=\epsilon\,\mu_j^{d_X}\phi'$ with the same sign $\epsilon=\pm1$, so $(\epsilon Q_{\tau,k})$ is a
decomposition of $F_s$. The terms with $k=1$ vanish, and for $k\ge2$, by \eqref{pf:volthm:eq:rest},
$\sum_\tau\sum_kr_j(p,\mathrm{frob}(p,\tau,\cdot),k,Q_{\tau,k})=\epsilon R^I_p$ with $I=\volthmdeg$, $v=1$, the
decomposition $(\epsilon Q_{\tau,k})$ and $\sigma(p,n+1,\tau)=\mathrm{frob}(p,\tau,n)$. By
\cref{pf:volthm:prop:restbound} with $P=S$,
the left side is at most $1.6(d_X+1)|\volthmdeg(\mu_j^{d_X}\phi')|=1.6(d_X+1)\,\mathrm{td}_j$, as $\mathrm{td}_j>0$.
\end{proof}

% ---------------------------------------------------------------------------
\subsection{Linking of good towers}\label{pf:volthm:sec:linking}

This subsection proves the two results behind the Euler product identity (\cref{pf:volthm:prop:prop6}): a
formula for the change of the arithmetic intersection numbers when the Chern-root lines are modified at finitely many
primes, and the construction of such modifications for the root models of a good origin at two cutoffs.

\begin{definition}[Second roots]\label{pf:volthm:def:secondroots}
Let $M$ be a root model of type $(\Psi,N)$ with coordinates $(K,\iota)$, ideal $P=P_M$ and Chern-root lines
$x_t(r)$ of multidegrees $a(r)$.
\begin{enumerate}[label=(\alph*)]
\item A system of \emph{second roots} $R$ of $M$ consists of multidegrees $a'(r)$, smooth conjugation-invariant
Hermitian lines $x'_t(r)$ of multidegree $a'(r)$ ($t\in\R$), and multihomogeneous integral polynomials $f_r,g_r$, for
$r\in\Psi\times[N]$. The root model $M^R$ is $M$ with $x_t(r)$ replaced by $x'_t(r)$ (same point, reference line,
$k_A$ and $\delta$).
\item Let $S$ be a finite set of primes. $R$ is a \emph{change at $S$} if for every $r$: $f_r,g_r\notin P$ and
$\deg f_r-\deg g_r=a(r)-a'(r)$; at every relevant prime $W\supseteq P$ which contains none of the primes of $S$ there are
$u,v\notin W$ with $uf_r-vg_r\in P$ (the divisor of $f_r/g_r$ on $P$ is vertical and lies over $S$); and for every
$t>0$ and every complex point $z$ of $P$, $|f_r(z)|\,h_{x'_t(r)}(z)=|g_r(z)|\,h_{x_t(r)}(z)$.
\item Let $j\in\N$ and $\phi_p$ ($p\in S$) be permutations of $\Psi\sqcup\Psi$. $R$ is \emph{local at the places of $S$
with labels $\phi$} if there is a Dedekind domain $A\subseteq\Omega_M$, finite over $\Z$ and containing the
constants of the point of $M$, such that for every $p\in S$ and every localisation $O$ of $A$ at a prime above $p$
(with $P_O$ the ideal of relations of the point over $O$) the following hold.
  \begin{enumerate}[label=(\roman*)]
  \item \emph{Flag Chern comparison}: for every $\tau$, $k\ge1$ and every list $L$ of $n+1-k$ multidegrees, the
  vertical intersection number over $O$ of the change data $(a,a',f,g)$ of the roots of $\tau$ in the direction of the
  $k$-th power sum, against $L$, equals the intersection number on the divisor of $p$ in $P_O$ of the polynomial
  $\mathrm{br}_j(\phi_p,p,k,\tau)$ of \cref{pf:volthm:lem:bracket}, evaluated at the multidegrees $a(r)$, against $L$.
  \item \emph{Vanishing of power sums}: for every $\tau$, every component $W$ of the special fibre of $P_O$, every
  $k\ge1$ and every list $L$ of $n-k$ multidegrees, the intersection number on $W$ of the $k$-th power sum of the
  multidegrees $a'(\tau,i)$, $i\in[N]$, against $L$, vanishes.
  \end{enumerate}
\item A root tower $T$ is \emph{linked} by $\phi\colon\N\to\mathrm{Sym}(\Psi\sqcup\Psi)$ with respect to $\volthmII$ if
for $2\le B_0\le B$ and every $c$ there are a root model $M$ and second roots $R$ of $M$ such that $M$ and $M^R$
converge, $\volthmPP_M(F^+)=\volthmPP_{B_0,c}(F)$, $\volthmPP_{M^R}(F^+)=\volthmPP_{B,c}(F)$ and
$\mathrm{geo}_M(F^+)=\mathrm{geo}_{M_{B_0,c}}(F^+)$ for all $F$, and $R$ is a change at, and local at the places of, the
set of primes $B_0<p\le B$ with $p\notin\mathrm{bad}$, with labels $\phi$.
\end{enumerate}
\formallabel{GlobalTheory.CompanionTheory.Ar.Num.SecondRoots}
\formallabel{GlobalTheory.CompanionTheory.Ar.Num.PowerSumsVanish}
\formallabel{GlobalTheory.CompanionTheory.Ar.Num.LinkedTower}
\formal{GlobalTheory.CompanionTheory.Ar.Num.SecondRoots.IsChange}
\formal{GlobalTheory.CompanionTheory.Ar.Num.SecondRoots.LocalAtPlaces}
\formal{GlobalTheory.CompanionTheory.Ar.Num.FlagChernComparison}
\end{definition}

\begin{proposition}[Height change formula]\label{pf:volthm:prop:heightchange}
Let $M$ be a root model, $R$ second roots of $M$, $\volthmII$ an intersection family, $S$ a finite set of primes,
$j\in\N$ and $\phi_p$ ($p\in S$) permutations of $\Psi\sqcup\Psi$. If $R$ is a change at $S$ and local at the places of
$S$ with labels $\phi$, then for every $t>0$, $\tau\in\Psi$, $1\le k\le N$ and $Q$ homogeneous of degree $n+1-k$,
\[
\mathrm{ar}^M_t\bigl((p_k(\tau)Q)^+\bigr)=\mathrm{ar}^{M^R}_t\bigl((p_k(\tau)Q)^+\bigr)+\sum_{p\in S}\log p\cdot
\mathrm{geo}_M\bigl((\mathrm{br}_j(\phi_p,p,k,\tau)Q)^+\bigr).
\]
\formallabel{GlobalTheory.T.power_sum_height_change_model}
\formallabel{GlobalTheory.CompanionTheory.Ar.Num.power_sum_height_model_change}
\end{proposition}

\begin{proof}
Write $F=p_k(\tau)Q$, homogeneous of degree $n+1$, and, for a family $\ell$ of lines indexed by the roots,
$\volthmII(F;\ell)=\sum_m\mathrm{coeff}_m(F)\,\volthmII(P;\ell^m)$. By the definition of $\mathrm{ar}_t$,
$\delta\,\mathrm{ar}^M_t(F^+)=\volthmII(F;x_t)$ and $\delta\,\mathrm{ar}^{M^R}_t(F^+)=\volthmII(F;x'_t)$. The proof
has three steps, each a result of the development which we do not reproduce.
\begin{enumerate}[label=(\arabic*)]
\item \emph{Mixed expansion}
(\lref{GlobalTheory.T.root_power_sum_mixed_expansion}): in two sets of variables $x$, $x'$ indexed by the roots,
$F(x)-F(x')=\sum_r(x_r-x'_r)H_r(x,x')$ with explicit polynomials $H_r$ homogeneous of degree $n$, built from the
expansion of $p_k(\tau)(x)-p_k(\tau)(x')$ and from $Q$.
\item \emph{Isometric families}
(\lref{GlobalTheory.T.isometric_family_intersection_eq}, which states that for families of lines which are
isometric through pairs $(f_r,g_r)$ that are vertical off $S$, the difference of the intersection numbers of a mixed
expansion is $\sum_{p\in S}\log p$ times the vertical intersection number over $\Z_{(p)}$ of the divisors of
$f_r/g_r$ against the $H_r$): since $R$ is a change at $S$, $\volthmII(F;x_t)-\volthmII(F;x'_t)=\sum_{p\in S}\log p\cdot
V_p$, where $V_p$ is that vertical number. (In terms of the inductive formula for arithmetic intersection numbers:
replacing $x_t(r)$ by $x'_t(r)$ changes the numbers by the intersection with the divisor of $f_r/g_r$, which is
vertical over $S$, while the archimedean term vanishes because $f_r/g_r$, as a rational section of
$x_t(r)\otimes x'_t(r)^{-1}$, has norm $1$.)
\item \emph{The local term}: $V_p=\delta\,\mathrm{geo}_M((\mathrm{br}_j(\phi_p,p,k,\tau)Q)^+)$ for $p\in S$. Over the
discrete valuation ring $\Z_{(p)}$, the flag Chern comparison and the vanishing of power sums at all places above $p$
imply the same statements over $\Z_{(p)}$ (\lref{GlobalTheory.T.flag_m3_padic_global} and
\lref{GlobalTheory.T.power_sums_vanish_at_places}); with them,
\lref{GlobalTheory.T.power_sum_term_discrete_valuation} computes $V_p$ as the intersection number, on the
divisor of $p$, of $\mathrm{br}_j(\phi_p,p,k,\tau)Q$ evaluated at the multidegrees $a(r)$. As $p$ is a uniformiser of
$\Z_{(p)}$, this equals the intersection number of the same polynomial on the generic fibre
(\lref{GlobalTheory.T.uniformizer_divisor_cycle_number}), which is $\delta\,\mathrm{geo}_M$ of it
(\lref{GlobalTheory.T.root_model_geometric_generic_number}; the geometric degrees of the Chern-root lines are
the intersection numbers of their multidegrees on the generic fibre).
\end{enumerate}
Dividing by $\delta$ gives the formula.
\end{proof}

\begin{proposition}[Link step]\label{pf:volthm:prop:linkstep}
Assume the numerical Segre statement of $(\mathrm{H}_0)$, $(\mathrm{H}_4)$ and $(\mathrm{H}_5)$, and let $\volthmII$ be
full. Let $S_0\subseteq S$ be finite sets of primes, $F_x$ fixed data on $Y$, $o_0$ and $o_1$ root origins
(\lref{GlobalTheory.CompanionTheory.Ar.RootOrigin}) over $F_x$ of root models $M_0$ at $S_0$ and $M_1$ at $S$, and
$o_R$ a root origin over $F_x$ of a root model $R_f$ at the empty set whose flag has degree $1$; assume that the
realisations of the three origins have no boundary (their boundary ideal is the unit ideal), that the Hodge components of $o_0$ and $o_1$ carry vector
bundles and have equal ranks at the primes of $S_0$, that every prime dividing $\mathrm{disc}(E)$ lies in $S_0$, that
$M_0$ and $M_1$ converge, that $2\notin S\setminus S_0$, and that for $p\in S\setminus S_0$, $\phi_p$ is a Frobenius
label at $p$ (\lref{GlobalTheory.CompanionTheory.Ar.IsFrobLabels}: $\phi_p$ describes the action of Frobenius at
the primes of $\OO_{c.F}$ above $p$ on the embeddings of $\OO_E$). Then there are a root model $M$ and second roots $R$
of $M$ such that $M$ and $M^R$ converge, $\volthmPP_M(F^+)=\volthmPP_{M_0}(F^+)$,
$\mathrm{geo}_M(F^+)=\mathrm{geo}_{M_0}(F^+)$ and $\volthmPP_{M^R}(F^+)=\volthmPP_{M_1}(F^+)$ for all $F\in\volthmR$,
and $R$ is a change at, and local at the places of, $S\setminus S_0$ with labels $\phi$.
\formallabel{GlobalTheory.CompanionTheory.Ar.Top.Link.link_step}
\formallabel{GlobalTheory.CompanionTheory.Ar.Top.Link.pullModel}
\end{proposition}

\begin{proof}
The construction realises the three origins on a common model and modifies the Chern-root lines of $M_0$ by
rational functions with vertical divisors. We use the following results of the development, which we do not
reproduce.

\emph{Common realisation.} \lref{GlobalTheory.T.root_origin_comparison_subset} provides a field $\Omega$ of
characteristic $0$, coordinates $(K,J)$, a point $y$ over $\Omega$, block maps $v_0,v_1,v_R$ from the coordinates of
$M_0,M_1,R_f$ to those of $y$ and field embeddings under which $y$ restricts to the points of the three models, a Hodge
component $\mathfrak M$ of $Y$ at $S$ at the point $y$ with a flag, and the comparison of the two levels $S_0\subseteq
S$ of this flag. By \lref{GlobalTheory.T.hodge_cpt_birational_arith_var}, $y$ is birational and $P_y$ is an
arithmetic variety of dimension $n+1$. Let $M$ be the \emph{pull-back} of $M_0$ to $y$: the root model with point $y$,
lines pulled back along $v_0$, and $\delta_M=[\Q(y):\Q(y\circ v_0)]\,\delta_{M_0}$, the degree of the extension of
function fields times $\delta_{M_0}$.

\emph{The change data.} \lref{GlobalTheory.T.flag_height_bound_of_arith_var} provides, for each root $r$,
multihomogeneous integral polynomials $f_r,g_r$ such that the pair formed by the pulled-back multidegrees of $M_0$ and
$M_1$ and $(f_r,g_r)$ is vertical off the primes of $S\setminus S_0$. Let $t_0>0$ be the regularisation threshold of
the analytic system of $F_x$. Define $x'_t(r)$ as the pull-back along $v_1$ of the Chern-root line $x_t(r)$ of $M_1$
if $0<t\le t_0$, and otherwise as a smooth conjugation-invariant line which is isometric to $x_t(r)$ of $M$ through
$(f_r,g_r)$ (\lref{GlobalTheory.T.exists_isometric_herm_line} states that such a line exists for every pair
vertical off a finite set of primes). This defines second roots $R$ of $M$.

\emph{Change and locality.} By \lref{GlobalTheory.T.hodge_comp_root_model_primes}, using the numerical Segre
statement, $(\mathrm{H}_4)$, $(\mathrm{H}_5)$, the Frobenius labels and the facts that the primes of $S\setminus S_0$
are odd, unramified in $E$ and Hodge--Tate free for $\mathfrak M$, the pair data of $R$ are vertical off
$S\setminus S_0$ and $R$ is local at the places of $S\setminus S_0$ with labels $\phi$. For $0<t\le t_0$ the lines
$x'_t(r)$ and $x_t(r)$ are isometric through $(f_r,g_r)$ by \lref{GlobalTheory.T.root_origin_height_comparison_three}
(the metrics of the root lines of the two origins differ by the absolute values of $f_r/g_r$), and for $t>t_0$ by
construction. So $R$ is a change at $S\setminus S_0$.

\emph{The limit functionals.} By \lref{GlobalTheory.T.root_model_block_pullback}, a root model and its
pull-back along a block map (with the matching $\delta$) have the same convergence behaviour, the same limit
functional and the same geometric functional on polynomials in the Chern roots; applied to $M_0$ and $M$ this gives
the convergence of $M$ and the identities for $\volthmPP_M$ and $\mathrm{geo}_M$. Let $M_1'$ be the pull-back of $M_1$
to $y$, with the pulled-back reference line of $M_1$; it converges with $\volthmPP_{M_1'}=\volthmPP_{M_1}$ on
polynomials in the Chern roots. The reference lines of $M$ and $M_1'$ are isometric through a pair vertical off the
empty set (\lref{GlobalTheory.T.root_origin_height_comparison_two}), and replacing the reference line by such a
line does not change $\mathrm{ar}_t$ (\lref{GlobalTheory.T.arMono_eq_relined_isometric_vert}). Hence for
$0<t\le t_0$, $\mathrm{ar}^{M^R}_t=\mathrm{ar}^{M_1'}_t$ on all monomials, so $M^R$ converges and
$\volthmPP_{M^R}(F^+)=\volthmPP_{M_1'}(F^+)=\volthmPP_{M_1}(F^+)$.
\end{proof}

% ---------------------------------------------------------------------------
\subsection{The conditions for good towers}\label{pf:volthm:sec:conditions}

In this subsection we fix data as in \cref{pf:volthm:def:VolX} with $p,q\ge1$, $Y$ global and $\volthmII$ full; $o$ is
a good origin of $T$ on $Y$, and $\Psi=\Psi_H$, $N=p+q$, $j=p$. The good origin provides, for every
component $c$, a compact analytic system $S_c$ of dimension $n$ (a compact Hausdorff space with a finite
holomorphic atlas, potentials indexed by $(\Psi\times[N])\sqcup\{\infty\}$ and their regularisations), and for
every $B$ a realisation of $P_{M_{B,c}}$ by $S_c$ (a resolution with a map to $S_c$) such that
$\delta(M_{B,c})=[c.F:\Q]\cdot\deg$ and the lines of the realisation are the fundamental combinations
$-\sum_{k\le i}x_t(\tau,k)$ of the Chern-root lines. We call this datum the \emph{real tower} $R_o$ of $o$. We use
the following properties of $R_o$, which are results of the development that we do not reproduce, or parts of the data of $o$:
\begin{itemize}
\item $R_o$ is \emph{uniformised} for $j$ (every point of the open part of each $S_c$ has a period chart for the
  potentials of the Chern-root lines and the index $j$, and at some point the period map is an immersion into the
  flag variety): \lref{GlobalTheory.T.good_tower_real_agreement};
\item $R_o$ is \emph{reference-linked}: for $B\ge2$ the polynomial of arithmetic intersection numbers of the
  realisation of $M_{B,c}$ differs from a fixed polynomial $P^{\mathrm{fix}}_c$ by a vertical perturbation at a
  finite set of primes $S$ with weights $\alpha_p\ge0$, $\sum_{p\in S}\alpha_p\log p\le C_l\log B$:
  \lref{GlobalTheory.T.good_tower_real_I6};
\item $R_o$ is \emph{smooth} and \emph{flat}, and \emph{positive} (the reference potential is positive); these are
  part of the data of $o$.
\end{itemize}
The analytic inputs used below are three results of the development, which we do not reproduce:
\emph{agreement} $(\mathrm{AG})$, \lref{GlobalTheory.T.agreement_positive}, which states that the arithmetic
intersection numbers of a realisation converge as $t\to0^+$ and that at $t=e^{-1/s^2}$ the polarised forms differ
from their limit by at most $C_3s$ on coefficient vectors of bounded $\ell^1$-norm; \emph{coefficient stability}
$(\mathrm{CO})$, \lref{GlobalTheory.T.coefficients'}, which states that the coefficients of the limit polynomials
of two realisations of the same system differing by a vertical perturbation at $S$ with weights $\alpha$ differ by
at most $C_2\sum_{p\in S}\alpha_p\log p$; and \emph{Hirzebruch--Mumford proportionality} $(\mathrm{HM})$,
\lref{GlobalTheory.CompanionTheory.Ar.PB.proportional}, which states that the degree polynomial of a
uniformised system of dimension $n$ is $\mathrm{vol}(S)>0$ times the degree polynomial of the flag variety.

\begin{proposition}[Proportionality]\label{pf:volthm:prop:prop}
Assume $(\mathrm{HM})$ and $(\mathrm{AI})$, and let $T$ be a root tower with a real tower $R$ which
is uniformised (any $T$, $R$; $\Psi$ nonempty). Then $T$ satisfies $\mathrm{Pr}$.
\formallabel{GlobalTheory.CompanionTheory.Ar.Top.rootTower_proportional_of_unif}
\formallabel{GlobalTheory.T.real_tower_proportional}
\end{proposition}

\begin{proof}
By the geometric agreement \lref{GlobalTheory.T.real_tower_geometric_agreement} (an identity of the development, not
reproduced, expressing $[c.F:\Q]\,\mathrm{geo}_{M_{B,c}}(\volthmlin^+(\lambda,s)^n)$ through the degree polynomial of
$S_c$ evaluated at $(\lambda,-s)$, with $k_A$ independent of $B$) and $(\mathrm{HM})$, for each $c$ there is
$\kappa_c>0$ with $\mathrm{geo}_{M_{B,c}}(\volthmlin(\lambda)^{n,+})=\kappa_c\volthmfd(\lambda)$ for all $B,\lambda$.
Since $\volthmfd(\rho)=n!>0$, $\vol(c)=\kappa_c>0$. The linear forms
$F\mapsto\mathrm{geo}_{M_{B,c}}(F^+)$ and $\kappa_c\volthmdeg$ vanish on homogeneous polynomials of degree $\ne n$ and
agree on all $n$-th powers of linear forms; as these span the homogeneous polynomials of degree $n$, they agree.
Finally, for $B,B'$ the geometric agreement expresses $\mathrm{geo}_{M_{B,c}}$ and $\mathrm{geo}_{M_{B',c}}$ on all
$\volthmlin^+(\lambda,s)^n$ by the same quantity (same $S_c$ and same $k_A$), and the same spanning argument in
$\volthmR^+$ gives $\mathrm{geo}_{M_{B,c}}=\mathrm{geo}_{M_{B',c}}$.
\end{proof}

\begin{proposition}[$\mathrm{Cv}$, $\mathrm{I}_2$ and $\mathrm{I}_1$]\label{pf:volthm:prop:I12}
Assume $(\mathrm{AG})$ and $(\mathrm{AI})$, and let $R$ be a real tower of $T$. Then $T$ satisfies
$\mathrm{Cv}$ and $\mathrm{I}_2$. If moreover $(\mathrm{CO})$ holds and $R$ is reference-linked, $T$
satisfies $\mathrm{I}_1$.
\formallabel{GlobalTheory.CompanionTheory.Ar.Top.real_tower_convergence_I2}
\formallabel{GlobalTheory.T.real_tower_conv_I2}
\formallabel{GlobalTheory.CompanionTheory.Ar.Top.TowerOverReals.tower_I1_of_reference_linked}
\formallabel{GlobalTheory.T.real_tower_I1}
\end{proposition}

\begin{proof}
The arithmetic agreement \lref{GlobalTheory.T.real_tower_arith_agreement} (not reproduced here)
identifies, for each $(B,c)$, $[c.F:\Q]\,\mathrm{ar}_t$ and $[c.F:\Q]\,\volthmPP_{M_{B,c}}$ on
$\volthmlin^+(\lambda,s)^{n+1}$ with the polarised arithmetic intersection numbers of the realisation (at the
parameter $t$, resp.\ in the limit), evaluated at the coefficient vector of $(\lambda,s/k_A)$, and shows that the
convergence of the realisation numbers implies convergence of $M_{B,c}$.

$\mathrm{Cv}$: apply $(\mathrm{AG})$ to $S_c$ (its parameter $t_0$ is positive) to get convergence. $\mathrm{I}_2$:
given $a$, apply $(\mathrm{AG})$ to each $S_c$ with the bound $2(a+\|\nu\|_1)+1$, and let $s_0$ be the minimum
of $1$ and the resulting thresholds and $C_3=\sum_c|C_3(c)|$. For $0<s\le s_0$ and $\|\lambda\|_1\le a$, the
coefficient vector of $(\lambda_s,s/k_A)$ has $\ell^1$-norm $\le\|\lambda\|_1+\|\nu\|_1+1\le2(a+\|\nu\|_1)+1$, so
$[c.F:\Q]\,|\mathrm{ar}(B,c,\lambda,s)-\Pi^+(B,c,\lambda_s,s)|\le C_3(c)s$, and $[c.F:\Q]\ge1$.

$\mathrm{I}_1$: let $C_l,(P^{\mathrm{fix}}_c)$ be the data of reference-linkedness, $C_2(c)$ the constants of
$(\mathrm{CO})$ for $S_c$, and $A=\sum_c\sum_m|\mathrm{coeff}_m(P^{\mathrm{fix}}_c)|$. For $B\ge2$ and a monomial
$x^m$ of degree $n+1$ (other degrees give $0$), $|\volthmPP_{M_{B,c}}(x^m)|$ is bounded by the absolute value of the
corresponding coefficient of the polarised limit polynomial (\cref{pf:volthm:lem:polar}, with the arithmetic
agreement and $[c.F:\Q]\ge1$); by $(\mathrm{CO})$ this coefficient differs from that of
$P^{\mathrm{fix}}_c$ by at most $|C_2(c)||C_l|\log B$. As $\log B\ge\log2$, $\mathrm{I}_1$ holds with constant
$A/\log2+\sum_c|C_2(c)||C_l|$.
\end{proof}

\begin{proposition}[Nefness and effectivity]\label{pf:volthm:prop:nef}
Assume $(\mathrm{H}_0)$.
\begin{enumerate}[label=(\alph*)]
\item $T$ satisfies $\mathrm{M}_1$.
\item $T$ satisfies $\mathrm{M}_2$.
\item Every root model whose reference line $A$ has nonnegative multidegree satisfies $\mathrm{Nef}_{\mathrm f}(0,1)$; hence
  $T$ satisfies $\mathrm{M}_0$.
\item If $o$ satisfies $(\mathrm{X})$, every $M_{B,c}$ satisfies $\mathrm{Eff}(\nu)$; with
  \lref{GlobalTheory.T.real_tower_nef_general} it follows that $T$ satisfies $\mathrm{I}_7$.
\end{enumerate}
\formallabel{GlobalTheory.CompanionTheory.Ar.Top.or_m1nef}
\formallabel{GlobalTheory.T.good_tower_M1_nef}
\formallabel{GlobalTheory.CompanionTheory.Ar.Top.m2_of_good_tower}
\formallabel{GlobalTheory.T.good_tower_M2}
\formallabel{RootFieldNumerics.root_model_nef_fibres_base}
\formallabel{GlobalTheory.CompanionTheory.Ar.Top.rootModel_nefFib_zero_one_of_nonneg}
\formallabel{GlobalTheory.T.root_model_nef_fibres}
\formallabel{RootFieldNumerics.global_model_tower_bounds}
\formallabel{GlobalTheory.CompanionTheory.Ar.Top.effDiff_of_goodOrigin}
\formallabel{GlobalTheory.T.good_tower_real_conditions}
\end{proposition}

\begin{proof}
(a) Let $p$ be in $S_B$; by construction of $o$, $p\in S_o(B)$, the Hodge component of $o$ at level $B$ is
Hodge--Tate free at $p$ (this is part of the data of $o$), and the root lines of the flag of $o$ satisfy the conclusion of
\lref{GlobalTheory.T.root_lines_weight_bound_of_free}, which, given the component
Hodge--Tate statement in $(\mathrm{H}_0)$, states that the root model satisfies $\mathrm{Nef}_p(\mu,0)$ for every
dominant $\mu$.

(b) Let $p$ be a prime not in $S_B$. Then $p\notin\mathrm{bad}$, so $p>2$ (as $2\in\mathrm{bad}$) and
$p\nmid\mathrm{disc}(E)$ (as $\mathrm{bad}$ contains the primes dividing $2\,\mathrm{disc}(E)M'_l$). Since $Y$ is global,
\lref{GlobalTheory.T.hodge_cpt_global_nef_bound} applies: for dominant $\mu$ with $\mathcal N_p(p,\mu)$ (recall
$j=p$) the root model of $o$ at level $B$ satisfies $\mathrm{Nef}_p(\mu,0)$.

(c) $\mathrm{Nef}_{\mathrm f}(0,1)$ asks that $\frac1{k_A}(\deg A\cdot W)\ge0$ for vertical curves $W$; the degree
of a nonnegative multidegree on a curve is $\ge0$ (\cref{pf:volthm:lem:curvedeg}). For the root
models of a good origin the multidegree of $A$ is $e_{k_0}$ (part of the data of the root origin).

(d) By \lref{GlobalTheory.T.root_origin_weight_identity} the root origin provides integers $n_r$ with
$n_r=k_Aq'\,\volthmcoef(\nu,-1)(r)$ (the fixed data of $o$ at the component carry integers $n_r=k_Aq'\nu_r$) and a nonzero
integral section of $\bigotimes_r\ell_1(r)^{n_r}$; $k_Aq'\ge1$. The two nefness parts of $\mathrm{I}_7$ are
\lref{GlobalTheory.T.real_tower_nef_general} (for real towers which are smooth and flat, given
$(\mathrm{NP})$, \lref{GlobalTheory.T.nef_polarization}).
\end{proof}

For a permutation $\phi$ of $\Psi\sqcup\Psi$, an integer $p\ge2$, $k\ge1$ and $\tau\in\Psi$ put
\begin{gather*}
c_{\phi,p,k}(\tau,\sigma)=\sum_{m\ge0:\ \phi^{m+1}(\tau^+)=\sigma}p^{-k(m+1)}-\frac{\mathbf 1_{\sigma=\tau^+}}{p-1},\\
\mathrm{br}_j(\phi,p,k,\tau)=k\sum_{\sigma\in\Psi\sqcup\Psi}c_{\phi,p,k}(\tau,\sigma)\,\pi_{j,k-1}(\sigma).
\end{gather*}

\begin{lemma}[Bracket polynomials]\label{pf:volthm:lem:bracket}
For $\phi,p,k,\tau$ as above and every $Q\in\volthmR$,
$\log p\cdot\volthmdeg(\mathrm{br}_j(\phi,p,k,\tau)Q)=\delta_j\bigl(p,\tau,(m\mapsto\phi^{m+1}(\tau^+)),k,Q\bigr)$.
\formallabel{GlobalTheory.CompanionTheory.Ar.Num.log_mul_degFl_bracket_mul}
\end{lemma}

\begin{proof}
For every function $g$ on $\Psi\sqcup\Psi$,
\[
\sum_\sigma c_{\phi,p,k}(\tau,\sigma)\,g(\sigma)=\sum_{m\ge0}p^{-k(m+1)}g(\phi^{m+1}(\tau^+))-\frac{g(\tau^+)}{p-1},
\]
by exchanging the finite sum over $\sigma$ with the series over $m$, which converges absolutely as $p^{-k}<1$; this
is the first part of \lref{GlobalTheory.CompanionTheory.Ar.Num.co_series_three_step}, which states exactly this
identity. Apply it to $g(\sigma)=\volthmdeg(\pi_{j,k-1}(\sigma)Q)$, use the linearity of $\volthmdeg$ and
$\pi_{j,k-1}(\tau^+)=p^{<j}_{k-1}(\tau)$, and compare with \eqref{pf:volthm:eq:delta}.
\end{proof}

\begin{proposition}[Euler product identity]\label{pf:volthm:prop:prop6}
Assume $(\mathrm{H}_9)$, $(\mathrm{AG})$, $(\mathrm{H}_4)$, $(\mathrm{H}_5)$ and
$(\mathrm{H}_0)$, and that $o$ satisfies $(\mathrm{X})$. Then there is $\phi\colon\N\to\mathrm{Sym}(\Psi\sqcup\Psi)$
such that $\mathrm{frob}(p,\tau,m)=\phi_p^{m+1}(\tau^+)$ satisfies $\mathrm{CS}$ and $T$ is \emph{linked} by $\phi$
(\cref{pf:volthm:def:secondroots}(d)); consequently $T$ satisfies $\mathrm{EP}(\mathrm{frob})$.
\formallabel{GlobalTheory.CompanionTheory.Ar.Top.Link.goodOrigin_linked}
\formallabel{GlobalTheory.T.good_tower_global_conditions}
\formallabel{GlobalTheory.CompanionTheory.Ar.Num.power_sum_term_linked_tower}
\formallabel{GlobalTheory.T.linked_tower_power_sum}
\formallabel{GlobalTheory.CompanionTheory.Ar.Num.PP_eq_add_of_arFun_eq_add}
\formallabel{GlobalTheory.T.limit_height_constant_shift}
\end{proposition}

\begin{proof}
\emph{Frobenius labels.} Since $E$ is abelian over $\Q$, \lref{GlobalTheory.T.frobenius_labels_shape} (not reproduced
here) gives permutations $\phi_p$ of $\Psi\sqcup\Psi$ which are Frobenius labels at $p$ (they describe the
action of Frobenius on the embeddings of $\OO_E$ into the residue fields of $c.F$) and such that the induced
$\mathrm{frob}$ satisfies $\mathrm{CS}$.

\emph{Linking.} For $2\le B_0\le B$ and a component $c$ apply the link step (\cref{pf:volthm:prop:linkstep}) to the
root origins of $o$ at the levels $B_0$ and $B$ over the fixed data of $c$ (so $S_0=S_o(B_0)$ and $S=S_o(B)$, and
$S\setminus S_0$ is the set of primes $B_0<p\le B$ outside $\mathrm{bad}$) and to the reference root origin of $o$
at $c$; its hypotheses are supplied by
the numerical Segre statement in $(\mathrm{H}_0)$, $(\mathrm{H}_4)$, $(\mathrm{H}_5)$,
$\mathrm{Cv}$ (\cref{pf:volthm:prop:I12}), the smoothness of $R_o$, and the following facts from $o$: the sets
$S_o(B_0)\subseteq S_o(B)$ are the primes in $S_{B_0}$, $S_B$; primes dividing $\mathrm{disc}(E)$ lie in
$S_o(B_0)$; the ranks $R_p$ of the Hodge components at $p\in S_o(B_0)$ agree at levels $B_0$ and $B$ (both equal
$4$ by $(\mathrm{X})$); and $2\in S_o(B_0)$.

\emph{From linking to $\mathrm{EP}$.} Fix $B_0\le B$, $c$, $\tau$, $1\le k\le N$, $Q$ homogeneous of degree
$n+1-k$, and $M,R$ as in the definition of linking, with prime set $S$. By the height change formula
(\cref{pf:volthm:prop:heightchange}), for every $t>0$,
$\mathrm{ar}^M_t(F)=\mathrm{ar}^{M^R}_t(F)+\sum_{p\in S}\log p\cdot\mathrm{geo}_M\bigl((\mathrm{br}_{j}(\phi_p,p,k,\tau)Q)^+\bigr)$
with $F=(p_k(\tau)Q)^+$ and $\mathrm{br}_j$ as in \cref{pf:volthm:lem:bracket}. If $\mathrm{ar}^M_t(F)=
\mathrm{ar}^{M^R}_t(F)+c'$ for all $t>0$ and both models converge, then $\volthmPP_M(F)=\volthmPP_{M^R}(F)+c'$ (uniqueness of
limits). Hence
$\volthmPP_{B_0,c}(p_k(\tau)Q)=\volthmPP_{B,c}(p_k(\tau)Q)+\sum_{p\in S}\log p\cdot\mathrm{geo}_{M_{B_0,c}}((\mathrm{br}_jQ)^+)$.
By $\mathrm{Pr}$ the last geometric degrees are $\vol(c)\volthmdeg(\mathrm{br}_jQ)$, and
$\log p\cdot\volthmdeg(\mathrm{br}_j(\phi_p,p,k,\tau)Q)=\delta_j(p,\tau,\mathrm{frob}(p,\tau,\cdot),k,Q)$ by
\cref{pf:volthm:lem:bracket}. This is \eqref{pf:volthm:eq:prop6}.
\end{proof}

\begin{proposition}[Galois symmetry]\label{pf:volthm:prop:propA}
Assume $(\mathrm{AG})$, $(\mathrm{H}_7)$, $(\mathrm{H}_8)$ and $(\mathrm{H}_0)$, and that $o$ satisfies
$(\mathrm{X})$. Then $T$ satisfies $\mathrm{GI}$.
\formallabel{GlobalTheory.CompanionTheory.Ar.Top.propA_of_goodOrigin}
\formallabel{GlobalTheory.T.good_tower_propA_of_global_extra}
\end{proposition}

\begin{proof}
By \cref{pf:volthm:prop:I12}, $T$ satisfies $\mathrm{Cv}$. The theorem
\lref{GlobalTheory.T.tower_origin_vector_bundle} (not reproduced) gives $\mathrm{GI}$ for the tower origin
underlying $o$, given $(\mathrm{H}_8)$, $(\mathrm{H}_7)$, $\mathrm{Cv}$, the Hodge--Tate freeness of the Hodge
components of $o$, the first part of $(\mathrm{X})$, and the equality of the ranks $R_p$ of the Hodge components of
any two components at $p\in S_o(B)$ (both are $4$ by the second part of $(\mathrm{X})$). The group $\Gamma$ is the
image of the automorphisms of the host acting on $\Psi_H$, which is transitive as $E$ is abelian.
\end{proof}

\begin{lemma}[The terms with $k=1$]\label{pf:volthm:lem:deltaone}
Let $j\le N$, $p\ge2$, $\tau\in\Psi$ and $\mathrm{fr}\colon\N\to\Psi\sqcup\Psi$. There is $e\in\R$ with
$|e|\le N/(p-1)$ and $\delta_j(p,\tau,\mathrm{fr},1,Q)=e\log p\,\volthmdeg(Q)$ for all $Q\in\volthmR$. Moreover
$e=0$ if $\mathrm{fr}$ takes values in the first copy of $\Psi$, and $e=(N-2j)p/(p^2-1)$ if $\mathrm{fr}(m)$ lies in
the first copy exactly for odd $m$.
\formallabel{GlobalTheory.T.delta_term_log_degree_bound}
\end{lemma}

\begin{proof}
The polynomials $\pi_{j,0}(\sigma^+)=j$ and $\pi_{j,0}(\sigma^-)=N-j$ are constants; write $\rho(\sigma)$ for
$\pi_{j,0}(\sigma)$, and note $p^{<j}_0(\tau)=j$. By \eqref{pf:volthm:eq:delta} with $k=1$ and linearity,
$\delta_j(p,\tau,\mathrm{fr},1,Q)=e\log p\,\volthmdeg(Q)$ with
\[
e=\sum_{m\ge0}p^{-(m+1)}\rho(\mathrm{fr}(m))-\frac j{p-1},
\]
the series converging absolutely. As $0\le\rho\le N$ and $\sum_{m\ge0}p^{-(m+1)}=\frac1{p-1}$, both terms lie in
$[0,\frac N{p-1}]$, so $|e|\le\frac N{p-1}$. If $\mathrm{fr}$ takes values in the first copy, then
$\rho(\mathrm{fr}(m))=j$ for all $m$ and $e=0$. In the alternating case $\rho(\mathrm{fr}(m))$ is $N-j$ for even $m$
and $j$ for odd $m$, so $\sum_mp^{-(m+1)}\rho(\mathrm{fr}(m))=\frac{(N-j)p+j}{p^2-1}$ and
$e=\frac{(N-j)p+j-j(p+1)}{p^2-1}=\frac{(N-2j)p}{p^2-1}$.
\end{proof}

\begin{proposition}[Flat sections, $\mathrm{FL}$]\label{pf:volthm:prop:propF}
Assume $(\mathrm{H}_1)$ and the locality of intersection numbers in the metrics and the degree formula from
$(\mathrm{AI})$, and let $T$ be a root tower and $\mathrm{frob}$ such that $T$ satisfies $\mathrm{Cv}$,
$\mathrm{Pr}$ and $\mathrm{EP}(\mathrm{frob})$, $r\ge1$, all elements of $\mathrm{bad}$ are prime, and for
every $c,\tau$ there is a flat datum $R_{c,\tau}$ of $M_{2,c}$ at $\tau$ (\cref{pf:volthm:def:flat}) whose height
$\mathrm{ht}(R_{c,\tau})$ is a \emph{near record}: $\mathrm{ht}(R_{c,\tau})=h_{\mathrm{rec}}+\sum_{p\in\mathrm{bad}}d_p$ with
$|h_{\mathrm{rec}}|\le N(\frac32\log r+\frac{|L|}2)$, $|d_p|\le(\frac N{p-1}+Nv_p)\log p$, $v_p\ge0$ and
$\sum_{p\in\mathrm{bad}}v_p\log p\le\log r$. Then $T$ satisfies $\mathrm{FL}(4,r,L)$.
\formallabel{GlobalTheory.CompanionTheory.Ar.Flat.propFTop}
\formallabel{GlobalTheory.T.PropFTop_holds}
\end{proposition}

\begin{proof}
For $p\ge2$ there is $E_p(\tau)$ with $|E_p(\tau)|\le N/(p-1)$ and
$\delta_j(p,\tau,\mathrm{frob}(p,\tau,\cdot),1,Q)=E_p(\tau)\log p\,\volthmdeg(Q)$ for all $Q$
(\cref{pf:volthm:lem:deltaone}). Let $G(B)$ be the set of primes $2<p\le B$, $p\notin\mathrm{bad}$, and put
$h(B,c,\tau)=\mathrm{ht}(R_{c,\tau})-\sum_{p\in G(B)}E_p(\tau)\log p$.

\emph{The identity.} For $B\ge2$ and $Q$ homogeneous of degree $n$, the push-forward formula for flat data
\lref{GlobalTheory.T.root_flat_pushforward_geoFun} (not reproduced here; it uses the three analytic inputs
and $\mathrm{Cv}$) and $\mathrm{Pr}$ give
$\volthmPP_{2,c}(p_1(\tau)Q)=\mathrm{ht}(R_{c,\tau})\vol(c)\volthmdeg(Q)$. By \eqref{pf:volthm:eq:prop6} with $B_0=2$, $k=1$,
$\volthmPP_{2,c}(p_1(\tau)Q)=\volthmPP_{B,c}(p_1(\tau)Q)+\vol(c)\sum_{p\in G(B)}E_p(\tau)\log p\,\volthmdeg(Q)$, so
$\volthmPP_{B,c}(p_1(\tau)Q)=h(B,c,\tau)\vol(c)\volthmdeg(Q)$.

\emph{The bound.} Let $B\ge128$ with $|\mathrm{bad}|\le B$. Then $h(B,c,\tau)=h_{\mathrm{rec}}+\sum_{p\in S}d'_p$ over
$S=\mathrm{bad}\cup G(B)$, with $d'_p=d_p$ on $\mathrm{bad}$ and $d'_p=-E_p(\tau)\log p$ on $G(B)$; so
$|d'_p|\le(\frac N{p-1}+Nv'_p)\log p$ with $v'_p=v_p$ on $\mathrm{bad}$ and $0$ on $G(B)$. Every $p\in S$ is prime with
$p\le B$ or $p\in\mathrm{bad}$, so by \lref{GlobalTheory.T.sum_log_prime_div_pred_bound} (not reproduced
here; it splits the primes $\le B$, bounded by $\log B+\frac{17}5$ via a Mertens-type estimate, from the at most
$B$ primes $>B$ in $\mathrm{bad}$, each contributing at most $\log(B+1)/B$),
$\sum_{p\in S}\frac{\log p}{p-1}\le3\log B+1$. Hence
\[
|h(B,c,\tau)|\le N(\tfrac32\log r+\tfrac{|L|}2)+N(3\log B+1)+N\log r\le4N(\log B+\log r+|L|)=4N\Lambda(B),
\]
using $\log B\ge1$ and $\log r\ge0$.
\end{proof}

\emph{Component settings.} The slope condition is verified on each component separately. A \emph{component
setting} $X$ consists of: integers $N\ge2$ and $1\le j\le N-1$ (put $f=|\Psi|$, $j'=N-j$ and $d_X=fjj'$); an
arithmetic variety $P_X$ of dimension $d_X+1$, with a boundary set of complex points, a constant $K_A$ and further
structure not needed here; Hodge data with a Hodge line $\mu$ of geometric degree $\deg\mu$; cusp data with a cusp
weight of mean $\bar U$; tied regularisations, which provide regularised Hermitian lines $\mathcal M_{a,b}(B,t)$ on
$P_X$ (for integers $1\le a\le b$, a cutoff $B$ and $t>0$) and an arithmetic limit $\mathrm{ar}_X(B)$; and reals
$r,L$, with $\Lambda_X(B)=\log(Br)+|L|$. The \emph{slope} of $X$ at $B$ is
$\mathrm{sl}_X(B)=\mathrm{ar}_X(B)/((d_X+1)\deg\mu)$. In the following conditions, ``for small $t$'' means: for every
$B$ in the stated range and all $1\le a\le b$ there is $t_0>0$ such that the property holds for
$\mathcal M=\mathcal M_{a,b}(B,t)$ whenever $0<t\le t_0$.
\begin{itemize}
\item \emph{Positivity}: for $B\ge4f+2$ and small $t$, $\mathcal M$ is semipositive on $P_X$, nef on the fibres and
  ample on the generic fibre.
\item \emph{Conjugation invariance}: for $B\ge4f+2$ and small $t$, $\mathcal M$ is conjugation invariant.
\item \emph{Height bounds} with constant $C_h$: for $B\ge4fN^2$ and small $t$, every arithmetic curve
  $W\supseteq P_X$, whether or not it meets the boundary, satisfies
  $\volthmII(W;\mathcal M)/(b\deg_\C(W))\ge-C_hfj'\Lambda_X(B)-\frac abK_A$.
\item \emph{Cusp mean bound} with constant $C_u$: if $j'\ge1$ and $3j'+1\le N$, then
  $\bar U\le C_ufj'\Lambda_X(B)$ for all $B\ge4fN^2$.
\item \emph{Stokes estimate} at $B$: for every $e\in\R$ such that for all $1\le a\le b$ and small $t$
  \begin{equation}\label{pf:volthm:eq:stokeshyp}
  (d_X+1)\bigl(e-\tfrac abK_A\bigr)\,\frac{\deg(P_X;\mathcal M)}{b^{d_X}}\le
  \frac{\volthmII(P_X;\mathcal M^{d_X+1})}{b^{d_X+1}},
  \end{equation}
  one has $(d_X+1)(e-\bar U/2)\deg\mu\le\mathrm{ar}_X(B)$.
\end{itemize}

\begin{lemma}[Slopes from the Stokes estimate]\label{pf:volthm:lem:slX}
Assume $(\mathrm{Z}^-_{\mathrm{rel}})$, let $\volthmII$ be an intersection family, and let $X$ be a component setting
with positivity, conjugation invariance, the height bounds with constant $C_h$, the cusp mean bound with constant
$C_u$, $\deg\mu>0$ and $3j'+1\le N$. If $B\ge4f+2$, $B\ge4fN^2$ and the Stokes estimate holds at $B$, then
$\mathrm{sl}_X(B)\ge-(C_h+C_u/2)fj'\Lambda_X(B)$.
\formallabel{GlobalTheory.CompanionTheory.Ar.Top.slX_lower_of_stokesAt}
\end{lemma}

\begin{proof}
Put $e=-C_hfj'\Lambda_X(B)$; we verify \eqref{pf:volthm:eq:stokeshyp}. Given $1\le a\le b$, let $t_0$ be the
smallest of the four thresholds of positivity, conjugation invariance and the two height bounds, let $0<t\le t_0$
and $\mathcal M=\mathcal M_{a,b}(B,t)$. Then $\mathcal M$ is semipositive on $P_X$, nef on the fibres and
conjugation invariant, and every arithmetic curve $W\supseteq P_X$ satisfies
$\volthmII(W;\mathcal M)/\deg_\C(W)\ge e'$ with $e'=b(e-\frac abK_A)$ (recall $\deg_\C(W)>0$,
\lref{GroundTheory.Heights₂.Heights.LowerBoundFree.degC_pos_of_isArithVar_one}). The relative lower bound
$(\mathrm{Z}^-_{\mathrm{rel}})$ is the statement $(\mathrm{Z}^-)$ with semipositivity required along the tangent
cones of an auxiliary ideal $P_0\subseteq P$; with $P_0=P=P_X$ it gives
$(d_X+1)e'\deg(P_X;\mathcal M)\le\volthmII(P_X;\mathcal M^{d_X+1})$, and dividing by $b^{d_X+1}>0$ gives
\eqref{pf:volthm:eq:stokeshyp}. By the Stokes estimate, $(d_X+1)(e-\bar U/2)\deg\mu\le\mathrm{ar}_X(B)$. Since
$j\le N-1$ we have $j'\ge1$, so the cusp mean bound gives $\bar U\le C_ufj'\Lambda_X(B)$. As $\deg\mu>0$, we get
$\mathrm{ar}_X(B)\ge-(d_X+1)\deg\mu\,(C_h+C_u/2)fj'\Lambda_X(B)$; divide by $(d_X+1)\deg\mu$.
\end{proof}

\begin{proposition}[Slope condition]\label{pf:volthm:prop:slope}
\begin{enumerate}[label=(\alph*)]
\item Assume $(\mathrm{Z}^-_{\mathrm{rel}})$, let $\volthmII$ be an intersection family, $\Psi$ nonempty, $T$ a root tower satisfying
  $\mathrm{Pr}$ with $3(N-j)+1\le N$, and suppose that for all $B\ge128$ with $4|\Psi|N^2\le B$ and
  $|\mathrm{bad}|\le B$ every component has a \emph{strict base degree setting} with constants $(C_h,C_u)$ and
  parameters $(r,L)$. Then $T$ satisfies $\mathrm{SL}(C_h+C_u/2,r,L)$.
\item If $C\le C'$, $r\ge1$ and $T$ satisfies $\mathrm{Pr}$, then $\mathrm{SL}(C,r,L)$ implies
  $\mathrm{SL}(C',r,L)$.
\item Assume $(\mathrm{H}_3)_{C_h,C_u}$, $(\mathrm{H}_2)$, $(\mathrm{AG})$, $(\mathrm{HM})$,
  $(\mathrm{H}_9)$, $(\mathrm{H}_0)$ and $(\mathrm{H}_1)$, that $o$ satisfies $(\mathrm{X})$,
  $q\le p$ and $|\mathrm{bad}|\le4|\Psi|N^2$. Then for every $B\ge4|\Psi|N^2$ every component has a strict base degree
  setting with constants $(C_h,C_u)$ and parameters $(\mathrm{rd}_E,L_H)$.
\end{enumerate}
\formallabel{GlobalTheory.CompanionTheory.Ar.Top.slope_lower_of_base_settings}
\formallabel{GlobalTheory.T.slope_lower_base_settings}
\formallabel{GlobalTheory.CompanionTheory.Ar.Top.ppc_lower_of_setting}
\formallabel{GlobalTheory.CompanionTheory.Ar.Cusp.slopeLower_mono_top}
\formallabel{GlobalTheory.CompanionTheory.Ar.Cusp.TopBase.base_setting_of_good_origin}
\formallabel{GlobalTheory.T.base_setting_of_good_tower}
\end{proposition}

\begin{proof}
A strict base degree setting of $(B,c)$ is a component setting $X$ with $N(X)=N$, $j(X)=j$, $r(X)=r$ and
$L(X)=L$ which satisfies positivity, conjugation invariance, the height bounds with constant $C_h$, the cusp mean
bound with constant $C_u$, $\deg\mu>0$, the existence of the arithmetic limits and the Stokes estimate at $B$, and
for which
\begin{equation}\label{pf:volthm:eq:slX}
\volthmPP_{B,c}(\mathrm{tp}_j)=\mathrm{sl}_X(B)\,(d_X+1)\vol(c)\,\mathrm{td}_j .
\end{equation}

(a) Here $\min(j,N-j)=N-j$, as $3(N-j)+1\le N$. Let $B$ be as in $\mathrm{SL}$; then
$B\ge4|\Psi|N^2\ge16|\Psi|\ge4|\Psi|+2$. For each component, \cref{pf:volthm:lem:slX}, applied to its setting $X$
(for which $f=|\Psi|$, $j'=N-j$ and $\Lambda_X=\Lambda$), gives $\mathrm{sl}_X(B)\ge-(C_h+C_u/2)|\Psi|(N-j)\Lambda(B)$.
Multiply by $(d_X+1)\vol(c)\mathrm{td}_j>0$
(\cref{pf:volthm:lem:testpoly}), use \eqref{pf:volthm:eq:slX}, and sum over $c$.

(b) The factor $|\Psi|\min(j,N-j)\Lambda(B)(d_X+1)\vol_T\mathrm{td}_j$ is $\ge0$ (as $\Lambda(B)\ge\log128>0$ for
$r\ge1$).

(c) Fix $B$ and a component $a$; note $B\ge4|\Psi|N^2\ge4|\Psi|+2$. \Cref{pf:volthm:prop:generic}, applied at the
generic point of the Hodge component of $o$ at level $B$ and $a$, with $S_0=\mathrm{bad}$, $S_1$ the set of primes
$\le B$ outside $\mathrm{bad}$ and $j=p$, gives a component datum with Hodge data and a component origin whose big side
is the first copy of $\Psi$; the origin is good, and together with the hypotheses on $\mathrm{bad}$, $B$ and $q\le p$
it is admitted at $B$. \Cref{pf:volthm:prop:cuspweight} provides a cusp weight $U$, a compact model and tied
regularisation settings; $(\mathrm{H}_3)_{C_h,C_u}$ then gives positivity, the height bounds and the cusp mean bound;
\lref{GlobalTheory.T.corrected_weight_regularization} (a result of the development which we do not reproduce)
turns these into a component setting $X$ with $\mathrm{rd}(X)=\mathrm{rd}_E$, $L(X)=L_H$ satisfying positivity,
conjugation invariance, the height bounds and the cusp mean bound. \Cref{pf:volthm:prop:slopeidentity} gives
$\deg\mu>0$, the existence of the arithmetic limit and the identity \eqref{pf:volthm:eq:slX}, and
\cref{pf:volthm:prop:stokes} gives the Stokes estimate at $B$.
\end{proof}

% ---------------------------------------------------------------------------
\subsection{Base degree settings for good origins}\label{pf:volthm:sec:basesettings}

This subsection proves the four results assembled in \cref{pf:volthm:prop:slope}(c): the construction of
component data at the generic points of the Hodge components, of cusp weights with regularisations, the slope
identity \eqref{pf:volthm:eq:slX}, and the Stokes estimate.

\emph{Component data.} We use the following notions of the development, defined elsewhere.
\begin{itemize}
\item A \emph{component datum} $D_c$ (\lref{GroundTheory.Heights₂.Heights.LowerBoundFree.ComponentData}) consists
of integers $N\ge2$ and $1\le j\le N-1$; a number field with an embedding into a field $\Omega$; a birational point $y$
over $\Omega$ whose ideal $P_X$ is an arithmetic variety of dimension $d_X+1$ ($d_X=|\Psi|j(N-j)$) with smooth generic
fibre; a closed, torus-stable and conjugation-stable \emph{boundary} set of complex points of $P_X$, covered by finitely
many simple normal crossing charts; vector bundles $W_\tau$ ($\tau\in\Psi$) on the open part with Hermitian metrics;
a finite set $S$ of primes with vertical divisors (given by multidegrees and rational functions) at its elements; and
a positive conjugation-invariant line $A'$ of positive multidegree whose heights on all points are bounded below by
$-K_A$. \emph{Hodge data} (\lref{GroundTheory.Heights₂.Heights.LowerBoundFree.HodgeData}) add Hodge subbundles
and the \emph{Hodge line} $\mu$ with its singular metric.
\item A \emph{component origin} (\lref{GlobalTheory.CompanionTheory.Ar.CptOrigin}) of $D_c$ over $Y$ at a set
$S_0$ of primes realises $D_c$ on a Hodge component of $Y$ at $S_0$, with a choice of the \emph{big} and the
\emph{small} side among the two copies of $\Psi$; it is \emph{good} if the Hodge component is Hodge--Tate free at the
primes of $S_0$ and of $S$.
\item The origin is \emph{admitted} at $B_1$ (\lref{GlobalTheory.CompanionTheory.Ar.Cusp.admitted}) if $Y$ is
global, the origin is good, $S_0$ is the set of primes dividing $2\,\mathrm{disc}(E)M'_l$ and has at most $4|\Psi|N^2$
elements, $B_1\ge4|\Psi|+2$, $S$ is the set of primes $\le B_1$ outside $S_0$, and $N-j\le j$.
\item A function $U\ge0$ on the complex points is a \emph{cusp weight} for the small side if at every interior complex
point its value is the cusp weight of the polarised lattice of the fibre, computed with the labels of the small side.
\item A \emph{regularisation setting} for a weight $w$ on the open part, relative to a compact complex model of
$P_X$, consists of a smooth reference line of the multidegree of $w$, chart data representing $w$ minus the reference
weight by plurisubharmonic functions (of log-log type for the Hodge line, bounded for the corrected weight), a
mollifier and a cutoff family; it is \emph{tied} if its charts are compatible with the coordinates of $P_X$. A
\emph{component setting} (\cref{pf:volthm:lem:slX}) combines Hodge data, cusp data with a cusp weight, a compact
model and tied regularisation settings of both kinds.
\end{itemize}

\begin{proposition}[Component data at generic points]\label{pf:volthm:prop:generic}
Assume $(\mathrm{H}_0)$ and $(\mathrm{H}_2)$, let $\volthmII$ be full, $Y$ global and $p,q\ge1$. Let $\xi_0$ be the generic
point of a component of $Y$, $S_0,S_1$ disjoint finite sets of primes with every prime dividing $\mathrm{disc}(E)$ in
$S_0$, and choose the index $j\in\{p,q\}$. Then there are a component datum $D_c$ of type $(\Psi,p+q)$ with index $j$
and $S=S_1$, Hodge data on it, and a component origin of $D_c$ over $Y$ at $S_0$ with generic point $\xi_0$, whose big
side is the first copy of $\Psi$ if $j=p$ and the second if $j=q$, whose Hodge component carries vector bundles, and
which is Hodge--Tate free at all primes of $S_0\cup S_1$.
\formallabel{GlobalTheory.T.full_intersection_generic_bound}
\formallabel{GlobalTheory.CompanionTheory.Ar.exists_componentData_origin}
\end{proposition}

\begin{proof}
This is a construction from the following results of the development, which we do not reproduce. By
\lref{GlobalTheory.T.global_snc_chart_system} there are coordinates, a point $y$ and a Hodge component of $Y$ at
$S_0$ with generic point $\xi_0$, with boundary equations, vertical divisor data at the primes of $S_1$, an
admissible Hodge line class, torsion and rank conditions and Hodge--Tate freeness at $S_0\cup S_1$, and simple normal
crossing data for the boundary. The transcendence degree of the coordinate field is
$|\Psi|j(p+q-j)=\dim X^o_l$, so $y$ is birational and $P_y$ is an arithmetic variety of dimension $d_X+1$
(\lref{GlobalTheory.T.birational_of_transcendence_degree}); the boundary points form a closed torus- and
conjugation-stable set (\lref{GlobalTheory.T.boundary_points_closed}) which is covered by torus-stable simple
normal crossing charts (\lref{GlobalTheory.T.exists_torus_stable_snc_charts}). A positive conjugation-invariant
line $A'$ with nonnegative heights exists by \lref{GlobalTheory.T.exists_positive_hermitian_line} (so $K_A=0$).
For each $\tau$ the sum of the Hodge lattices of the big and small labels is a vector bundle of rank $p+q$
(\lref{GlobalTheory.T.pair_lattice_sum_comap}), which gives chart bundles $W_\tau$
(\lref{GlobalTheory.T.exists_chart_bundle}) with Hermitian metrics satisfying the degree bound of
\lref{GlobalTheory.T.pair_lattice_degree_bound}. These data form the component datum, and the Hodge data with
the required intersection identity are provided by \lref{GlobalTheory.T.component_data_intersection_sum}.
\end{proof}

\begin{proposition}[Cusp weights and tied regularisations]\label{pf:volthm:prop:cuspweight}
Assume $(\mathrm{H}_0)$, let $E$ be abelian over $\Q$, and let $D_c$ be a component datum with Hodge data and a
component origin over $Y$ at $S_0$. Then there are a cusp weight $U$ for the small side, with $U\ge0$ and $U$ invariant
under the torus action on the open part, a compact complex model of $P_X$, a tied regularisation setting of log-log type
for the Hodge line, and a tied, conjugation-invariant regularisation setting of bounded type for the weight
$\phi_\mu+U/2$ (the weight of the Hodge line corrected by half the cusp weight).
\formallabel{GlobalTheory.T.cusp_weight_regularisations_tied}
\formallabel{GlobalTheory.CompanionTheory.Relay.cusp_weight_regularisations_tied}
\end{proposition}

\begin{proof}
We use the following results of the development, which we do not reproduce.
\lref{GlobalTheory.T.intersection_system_weight_left} provides a cusp weight $U\ge0$ for the small side, invariant
under the torus and under conjugation. By \lref{GlobalTheory.T.cusp_intersection_positive}, the Hodge weight is
plurisubharmonic in the charts and the corrected weight is plurisubharmonic and conjugation invariant; by
\lref{GlobalTheory.T.cusp_chart_log_growth_bounds} the corrected weight has at most logarithmic growth near the
boundary in the simple normal crossing charts. \lref{GlobalTheory.CompanionTheory.Relay.hodge_metric_frame_properties}
relates the Hodge weight to the metrics of frames, and a compact complex model exists by
\lref{GlobalTheory.CompanionTheory.Relay.exists_compactified_model}. Given these,
\lref{GlobalTheory.T.exists_regularization_of_hodge_metric} constructs the tied log-log regularisation setting of
the Hodge line, and \lref{GlobalTheory.T.exists_regularization_setting} the tied, conjugation-invariant bounded
regularisation setting of the corrected weight.
\end{proof}

\begin{proposition}[Arithmetic limits and the slope identity]\label{pf:volthm:prop:slopeidentity}
Assume $(\mathrm{AG})$, $(\mathrm{HM})$, $(\mathrm{H}_9)$, $(\mathrm{H}_0)$ and $(\mathrm{H}_1)$; let $E$ be abelian over
$\Q$, $Y$ global, $o$ a good origin of $T$ on $Y$ satisfying $(\mathrm{X})$, $\volthmII$ full, $1\le q\le p$, $2\le B\le
B_1$ and $a$ a component. Let $X$ be a component setting on a component datum $D_c$ with a good component origin over
$Y$ at $\mathrm{bad}$, with big side the first copy and small side the second copy of $\Psi$, with the same generic
point as the fixed data of $o$ at $a$, and with $S$ the set of primes $\le B_1$ outside $\mathrm{bad}$. Then
$\deg\mu>0$, the arithmetic limit $\mathrm{ar}_X(B)$ exists, and
$\volthmPP_{B,a}(\mathrm{tp}_j)=\mathrm{sl}_X(B)\,(d_X+1)\vol(a)\,\mathrm{td}_j$.
\formallabel{GlobalTheory.T.good_origin_intersection_tower}
\formallabel{GlobalTheory.CompanionTheory.Ar.Cusp.testPoly_eq_slX}
\end{proposition}

\begin{proof}
By \cref{pf:volthm:prop:prop} and the uniformisation of the real tower of $o$
(\lref{GlobalTheory.T.good_tower_real_agreement}), $T$ satisfies $\mathrm{Pr}$; and $\mathrm{td}_j>0$
(\cref{pf:volthm:lem:testpoly}). The following are results of the development which we do not reproduce. The
arithmetic numbers $\mathrm{ar}_{X,t}(B)$ of the regularised Hodge line $\mu_{B,t}$ converge as $t\to0^+$ to a limit
$L=\mathrm{ar}_X(B)$ (\lref{GlobalTheory.T.height_comparison_between_models}), its geometric degrees converge to
$\deg\mu$ (\lref{GlobalTheory.T.measure_limit_statement_holds}), and $\mu_{B,t}$ is smooth for small $t$. The key
input is the comparison \lref{GlobalTheory.T.good_origin_intersection_relay}: it realises the root model $M_{B,a}$ on
an arithmetic variety $P_3$ with a block map onto $P_X$: for small $t$ the line $\mu_{B,t}$ pulls back to a line
$M_t$, and there are lists of lines $N_{t,m}$ along the fibres, whose multidegrees do not depend on $t$, real
coefficients $c_m$ ($1\le m\le r$) and an integer $\delta>0$ such that
\[
\sum_mc_m\,\volthmII\bigl(P_3;M_t^{d_X+1},N_{t,m}\bigr)\longrightarrow\delta\,\volthmPP_{B,a}(\mathrm{tp}_j)\qquad(t\to0^+),
\]
and such that, for small $t$, $\delta\,\mathrm{geo}_{M_{B,a}}((\mu_j^{d_X}\phi_j)^+)$ is $\sum_mc_m$ times the integral
over $P_3$ of the product of the forms of $d_X$ copies of $M_t$ and of the lines $N_{1,m}$. By the projection
formula for the block map (part of $(\mathrm{AI})$) and its analogue for integrals
(\lref{GlobalTheory.T.ptsIntegral_fibre_append_eq_mul_image}), both sums equal $\varphi$ times the corresponding
quantity on $P_X$, where $\varphi=\sum_mc_m\varphi_m$ and $\varphi_m$ is the fibre degree of the multidegrees of
$N_{1,m}$. Passing to the limit $t\to0^+$:
\[
\delta\,\volthmPP_{B,a}(\mathrm{tp}_j)=\varphi L,\qquad
\delta\,\mathrm{geo}_{M_{B,a}}\bigl((\mu_j^{d_X}\phi_j)^+\bigr)=\varphi\deg\mu .
\]
By $\mathrm{Pr}$ the second left side is $\delta\vol(a)\mathrm{td}_j>0$, so $\varphi\deg\mu>0$; since
$\deg\mu\ge0$ (\lref{GlobalTheory.CompanionTheory.StokesCompare.degMu_nonneg}), $\deg\mu>0$. Finally $j=p$ and
$d_X=|\Psi|p\,q$ is the dimension in $D_c$, so
$\volthmPP_{B,a}(\mathrm{tp}_j)=\varphi L/\delta=\frac{L}{(d_X+1)\deg\mu}(d_X+1)\vol(a)\mathrm{td}_j=
\mathrm{sl}_X(B)(d_X+1)\vol(a)\mathrm{td}_j$.
\end{proof}

\begin{proposition}[The Stokes estimate]\label{pf:volthm:prop:stokes}
Assume $(\mathrm{AI})$, the analytic package, resolution in projective coordinates, GAGA for sections and the extension
of homomorphisms of semi-abelian schemes (all from $(\mathrm{H}_0)$); let $E$ be abelian over $\Q$, $\volthmII$ full, and
let $D_c$ be a component datum with Hodge data and a component origin admitted at some $B_1$, $U$ a cusp weight for the
small side, and $X$ a component setting with these Hodge data and cusp weight $U$ on the open part which satisfies
positivity. Then for every $B\ge4|\Psi|+2$ for which the arithmetic limit $\mathrm{ar}_X(B)$ exists, the Stokes
estimate holds at $B$.
\formallabel{GlobalTheory.T.intersection_final_estimate}
\formallabel{GlobalTheory.CompanionTheory.Relay.intersection_final_estimate}
\formallabel{ModelDimension.abelian_cm_stokes_identity}
\formallabel{ModelDimension.OrbitMargin.stokes_of_model}
\formallabel{ModelDimension.OrbitMargin.stokes_of_statements}
\formallabel{GroundTheory.Heights₂.Heights.LowerBoundFree.StokesLimitPos.proof}
\end{proposition}

\begin{proof}
For $1\le a\le b$ and $t>0$ let $\bar N_{a,b}(t)$ be the line with the multidegree of $\mathcal M_{a,b}(B,t)$ whose metric
is built in the same way from the log-log regularisation of the Hodge line instead of the bounded regularisation of
the corrected weight, and put $\mathrm{top}_N(t)=\volthmII(P_X;\bar N_{a,b}(t)^{d_X+1})/b^{d_X+1}$ and
$\deg_N(t)=\deg(P_X;\bar N_{a,b}(t))/b^{d_X}$. Let $I_0(t,t')$ be the integral over $P_X$ of a Stokes weight
against the form of $\bar N_{a,b}(t)$, and $V_b(t)$ the integral of $U/2$ against the same form. We use the following
results of the development, which we do not reproduce; they are derived from the analytic inputs listed in the
statement.
\begin{enumerate}[label=(\roman*)]
\item \emph{Semipositivity} (\lref{GroundTheory.Heights₂.Heights.LowerBoundFree.NbarSemipos.proof}): for small $t$,
$\bar N_{a,b}(t)$ is semipositive.
\item \emph{Telescoping} (\lref{GroundTheory.Heights₂.Heights.LowerBoundFree.StokesTelescope.proof}, obtained by
repeated integration by parts, each step of which contributes a term of the right sign by
\lref{GroundTheory.Heights₂.Heights.LowerBoundFree.StokesStep.proof}): if $\mathcal M_{a,b}(B,t')$ and
$\bar N_{a,b}(t)$ are semipositive, then
$\volthmII(P_X;\mathcal M_{a,b}(B,t')^{d_X+1})/b^{d_X+1}-\mathrm{top}_N(t)\le(d_X+1)I_0(t,t')$.
\item \emph{Equal degrees} (\lref{GroundTheory.Heights₂.Heights.HilbertSamuel.GenDegGeomDeg.proof}, which states that the
geometric degree of a smooth semipositive line is the degree of its multidegree on the generic fibre): in
particular $\deg(P_X;\mathcal M_{a,b}(B,t'))/b^{d_X}=\deg_N(t)$ when both lines are semipositive.
\item \emph{Limit in $t'$} (\lref{GroundTheory.Heights₂.Heights.LowerBoundFree.StokesLimitTp.proof}): for small $t$,
$I_0(t,t')$ converges as $t'\to0^+$ to a limit $V(t)\le V_b(t)$.
\item \emph{Limit in $t$} (\lref{GroundTheory.Heights₂.Heights.LowerBoundFree.StokesLimitT.proof}):
$\mathrm{top}_N(t)$, $V_b(t)$ and $\deg_N(t)$ converge as $t\to0^+$, to limits $\mathrm{top}_N^*$, $V_b^*$, $\deg_N^*$.
\item \emph{Limit in the correction} (\lref{GroundTheory.Heights₂.Heights.LowerBoundFree.StokesLimitEps.proof}):
for $(a,b)=(1,m+1)$ and $m\to\infty$, these limits converge to $\mathrm{ar}_X(B)$, $\bar U\deg\mu/2$ and $\deg\mu$.
\end{enumerate}
Now let $e$ satisfy the hypothesis \eqref{pf:volthm:eq:stokeshyp} of the Stokes estimate, let $m\ge0$ and $(a,b)=(1,m+1)$.
For $t$ and $t'$ small, (i), positivity of $\mathcal M_{1,m+1}(B,t')$, the hypothesis, (ii) and (iii) give
\[
(d_X+1)\Bigl(e-\frac{K_A}{m+1}\Bigr)\deg_N(t)-(d_X+1)I_0(t,t')\le\mathrm{top}_N(t).
\]
Letting $t'\to0^+$ and using (iv), then $t\to0^+$ and (v),
$(d_X+1)(e-\frac{K_A}{m+1})\deg_N^*-(d_X+1)V_b^*\le\mathrm{top}_N^*$ for the pair $(1,m+1)$. Letting $m\to\infty$ and
using (vi) gives $(d_X+1)e\deg\mu-(d_X+1)\bar U\deg\mu/2\le\mathrm{ar}_X(B)$, which is the conclusion of the Stokes
estimate.
\end{proof}

% ---------------------------------------------------------------------------
\subsection{Equality of volumes and flat data}\label{pf:volthm:sec:voleq}

\begin{lemma}[Degree of the Hodge line by sheets]\label{pf:volthm:lem:degmu}
There are reals $\kappa(f,j,j')>0$ with the following property. Let $D_c$ be a component datum with Hodge data and a
component origin over $Y$ whose small side is the second copy of $\Psi$, with $j=p$ and $j'=N-j=q$, $1\le q$ and
$3q+1\le p+q$, and let $\mathcal B$ be a uniformisation of $Y$ by matrix balls $D_k$ ($k<n_{\mathcal B}$) with arithmetic
groups $\Gamma_k$ as in $(\mathrm{a})$ of the proof of \cref{pf:volthm:prop:voleqproof}, with fundamental sets $F_k$ of
common volume $v$ for the invariant densities. Then
$\deg\mu=\kappa(|\Psi|,p,q)\cdot\deg\mathrm{Or}\cdot[c.F:\Q]\cdot v$, where $\deg\mathrm{Or}$ is the degree of the
component origin.
\formallabel{SheetDecomposition.degMu_eq}
\end{lemma}

\begin{proof}
We use the following results of the development, which we do not reproduce. By
\lref{SheetDecomposition.exists_sheets}, there is an open set $U$ of $Y$ containing the generic point of the Hodge
component such that the complex points of $P_X$ above $U$ decompose, up to null sets, into \emph{sheets}: over each
fundamental set $F_k$ whose image meets $U$ lie exactly $\deg\mathrm{Or}$ holomorphic sheets (the roots of a separable
monic polynomial in a coordinate of the cover), and none over the others. By
\lref{SheetDecomposition.transferU} (a change of variables for Monge--Amp\`ere densities: on each sheet the form of
the Hodge line is $\kappa(f,j,j')$ times the invariant density of the ball), the integral defining $\deg\mu$ is
$\kappa(|\Psi|,p,q)\sum_k n_k\int_{F_k}\mathrm{dens}_k$, where $n_k$ is the number of sheets over $F_k$; here the
invariant densities are integrable on the fundamental sets (first part of
\lref{GlobalTheory.T.cusp_weight_integrable_on_fundamental_set}). By
\lref{SheetDecomposition.onePiece}, for every embedding $\kappa\colon c.F\to\C$ exactly one piece $k$ with $\kappa_k=\kappa$
has an image meeting $U$; as $c.F$ has $[c.F:\Q]$ complex embeddings, $\sum_kn_k\int_{F_k}\mathrm{dens}_k=
\deg\mathrm{Or}\cdot[c.F:\Q]\cdot v$.
\end{proof}

\begin{proposition}[Equality of volumes]\label{pf:volthm:prop:voleqproof}
Assume $(\mathrm{H}_9)$ and $(\mathrm{H}_0)$; let $E$ be abelian over $\Q$, $Y$ global, $o$ a good origin of $T$ on $Y$
satisfying $(\mathrm{X})$, and $\volthmII$ full. If $q\ge1$ and $3q+1\le p+q$, then $T$ satisfies $\mathrm{EV}$.
\formallabel{GlobalTheory.T.good_tower_intersection_extra}
\formallabel{GlobalTheory.CompanionTheory.Relay.good_tower_intersection_extra}
\formallabel{SheetDecomposition.volEq_of_isFullIntersectionSystem}
\formallabel{SheetDecomposition.volEq_of_degPush}
\formallabel{SheetDecomposition.volEq_of}
\end{proposition}

\begin{proof}
The volumes of the components are compared through the hyperbolic volumes of the ball quotients which uniformise
them. We use the following results of the development, which we do not reproduce.
\begin{enumerate}[label=(\alph*)]
\item \emph{Uniformisation} (\lref{SheetDecomposition.baseUnifEx2}): there are finitely many lattices $L_k$ of
signature $(p,q)$ over $\OO_E$, embeddings $\kappa_k\colon c.F\to\C$, the arithmetic groups $\Gamma_k$ of level-$l$
automorphisms of $L_k$, and holomorphic maps from the matrix balls $D_k$ of dimension $n=|\Psi|pq$ to the complex
points of $Y$ above $\kappa_k$, which together cover the abelian locus, have disjoint images for distinct $k$, and
identify exactly the points equivalent under $\Gamma_k$.
\item \emph{Volumes of the ball quotients}: every $\Gamma_k$ has a fundamental set $F_k$
(\lref{GlobalTheory.T.exists_fundamental_set}), and the volumes $\int_{F_k}\mathrm{dens}_k$ for the invariant
densities do not depend on $k$ (\lref{SheetDecomposition.ballVolEq}; this uses $q\ge1$ and $3q+1\le p+q$). Call
the common value $v$.
\item \emph{Degree push} (\lref{SheetDecomposition.degPush}): there are reals $c_0(f,p,q)\ne0$ such that, given
$(\mathrm{AG})$, $(\mathrm{HM})$, $(\mathrm{H}_9)$, $(\mathrm{H}_0)$ and $(\mathrm{H}_1)$, for $Y$ global, a good origin $o$
with $(\mathrm{X})$, $\volthmII$ full and $1\le q\le p$, every component $a$, and every good component datum with Hodge data and a component
origin at $\mathrm{bad}$ with big side the first copy, carrying vector bundles, with $S=\varnothing$ and with the
generic point of the fixed data of $o$ at $a$: $\deg\mu=c_0(|\Psi|,p,q)\cdot\deg\mathrm{Or}\cdot[c.F:\Q]\cdot\vol(a)$.
\end{enumerate}
Fix a component $a$ and apply \cref{pf:volthm:prop:generic} at the generic point of the Hodge component of $o$ at
level $2$ and $a$, with $S_0=\mathrm{bad}$, $S_1=\varnothing$ and $j=p$; the resulting component origin has big side
the first copy and small side the second, carries vector bundles and is good. By (c) (whose hypotheses $(\mathrm{AG})$,
$(\mathrm{HM})$ and $(\mathrm{H}_1)$ are results of the development: \lref{GlobalTheory.T.agreement_positive},
\lref{GlobalTheory.CompanionTheory.Ar.PB.proportional} and
\lref{GlobalTheory.T.points_integrable_of_resolution_statement}, the last derived from the resolution statement in
$(\mathrm{H}_0)$) and \cref{pf:volthm:lem:degmu},
\[
c_0(|\Psi|,p,q)\,\deg\mathrm{Or}\,[c.F:\Q]\,\vol(a)=\deg\mu=\kappa(|\Psi|,p,q)\,\deg\mathrm{Or}\,[c.F:\Q]\,v .
\]
As $\deg\mathrm{Or}\ge1$, $\vol(a)=\kappa(|\Psi|,p,q)v/c_0(|\Psi|,p,q)$ for every $a$, which is $\mathrm{EV}$.
\end{proof}

\begin{definition}[Flat data and near records]\label{pf:volthm:def:flat}
Let $M$ be a root model of type $(\Psi,N)$ and $\tau\in\Psi$. A \emph{flat datum} $R$ of $M$ at $\tau$ consists of a
number field $F$ with an embedding into $\Omega_M$ which is integral over the coordinates, a smooth Hermitian line
$A$ of multidegree $\sum_ia(\tau,i)$ with a flat Hermitian structure over $F$ (whose arithmetic degree is the
\emph{height} $\mathrm{ht}(R)$), and smooth torus-invariant functions $v_t$ ($t\in\R$) such that
$\prod_ih_{x_t(\tau,i)}=h_A\,e^{-v_t}$ on the complex points of $P_M$ for $t>0$ and the integrals of $v_t$ against the
forms of all top-degree monomials in the Chern-root lines tend to $0$ as $t\to0^+$. For $N$, a finite set $\mathrm{bad}$
of primes and reals $r,L$, a real $h$ is a \emph{near record} if $h=h_{\mathrm{rec}}+\sum_{p\in\mathrm{bad}}d_p$ with
$|h_{\mathrm{rec}}|\le N(\frac32\log r+\frac{|L|}2)$, $|d_p|\le(\frac N{p-1}+Nv_p)\log p$, $v_p\ge0$ and
$\sum_{p\in\mathrm{bad}}v_p\log p\le\log r$.
\formallabel{GlobalTheory.CompanionTheory.Ar.Flat.RootFlat}
\formallabel{GlobalTheory.CompanionTheory.Ar.Flat.NearRecord}
\end{definition}

\begin{proposition}[Flat data with near-record heights]\label{pf:volthm:prop:flatdata}
Assume $(\mathrm{H}_6)$ and $(\mathrm{H}_0)$, let $E$ be abelian over $\Q$, $Y$ global and $p,q\ge1$.
\begin{enumerate}[label=(\alph*)]
\item Let $M$ be a root model with a root origin over $Y$ at the set $\mathrm{bad}$ of primes dividing
$2\,\mathrm{disc}(E)M'_l$, whose Hodge component carries vector bundles. For every $\tau\in\Psi$ there is a flat datum
of $M$ at $\tau$ whose height is a near record for $(p+q,\mathrm{bad},\mathrm{rd}_E,L_H)$.
\item If $o$ is a good origin of $T$ on $Y$ (satisfying $(\mathrm{X})$, and $\volthmII$ full), then for every component
$a$ and $\tau\in\Psi$ there is a flat datum of $M_{2,a}$ at $\tau$ whose height is a near record for
$(p+q,\mathrm{bad},\mathrm{rd}_E,L_H)$.
\end{enumerate}
\formallabel{GlobalTheory.T.good_tower_real_extra}
\formallabel{GlobalTheory.CompanionTheory.Relay.good_tower_real_extra}
\formallabel{RootFieldNumerics.global_model_tower_final_condition}
\formallabel{RootFieldNumerics.rootFlatOnHodge}
\formallabel{RootFieldNumerics.RF.rootFlatOnHodge_of}
\formallabel{GlobalTheory.CompanionTheory.Ar.Flat.RootFlatOnHodge}
\formallabel{GlobalTheory.CompanionTheory.Ar.Flat.OnHodge.tower}
\end{proposition}

\begin{proof}
(b) follows from (a) applied to $M_{2,a}$ and the root origin of $o$ at level $2$: the set of primes of $o$ at level
$2$ is $S_2=\mathrm{bad}$, since $2\in\mathrm{bad}$ and the primes of $\mathrm{bad}$ are those dividing
$2\,\mathrm{disc}(E)M'_l$, and the Hodge components of $o$ carry vector bundles. (The hypotheses $(\mathrm{X})$ and the
fullness of $\volthmII$ are not used.)

For (a) we use the following results of the development, which we do not reproduce.
\lref{RootFieldNumerics.RF.cptEx} provides an open set of $Y$ and a component through the generic point of the
Hodge component, and \lref{RootFieldNumerics.RF.exists_norms} (from $(\mathrm{H}_6)$, the period and character
statements \lref{RootFieldNumerics.RF.periodOf}, \lref{RootFieldNumerics.RF.chiEx} and $(\mathrm{H}_0)$)
provides the norms of the flat generators of the line bundles of the flag. The core construction
\lref{RootFieldNumerics.RF.core} then produces a number field $F'$, a smooth line $A$ of multidegree
$\sum_ia(\tau,i)$ with a flat Hermitian structure whose metric equals $\prod_ih_{\mathrm{sing}(\tau,i)}$ at the interior
points (the product of the singular Hodge metrics of the roots of $\tau$), and whose height is a near record; it uses a
smoothness statement (\lref{RootFieldNumerics.RF.smoothLine}), a local comparison
(\lref{RootFieldNumerics.RF.localCmp_of}), the numerical estimates \lref{RootFieldNumerics.RF.numerics}, a part
of $(\mathrm{H}_0)$, and, at the primes outside $\mathrm{bad}$, the splitting of the local systems on the fibres
(\lref{GlobalTheory.T.local_system_split_fibres}). Finally \lref{GlobalTheory.T.flat_herm_root_line_degree}
turns $A$ into a flat datum of $M$ at $\tau$ with the same height. It uses the realisation of the origin: its
fundamental line of index $(\tau,N-1)$ has multidegree $-\sum_ia(\tau,i)$, its regularised metric is the inverse of
the product of the metrics of the Chern-root lines $x_t(\tau,i)$, and its singular metric is related in the same way
to the Hodge lines $\mathrm{sing}(\tau,i)$.
\end{proof}

% ---------------------------------------------------------------------------
\subsection{Heights of points and the antisymmetry identity}\label{pf:volthm:sec:deep}

This subsection proves the height bounds $\mathrm{I}_5$, $\mathrm{I}_6$, the antisymmetry identity $\mathrm{I}_8$ and the
semipositivity condition $\mathrm{A}_1$ for good towers. In both cases the proof reduces the conditions on the tower to
statements about the reference root models and the real tower of the good origin; these statements are the main
geometric results of the development, and we state precisely what is used of them.

\begin{definition}[Reference models]\label{pf:volthm:def:reference}
Let $T$ be a root tower of type $(\Psi,N)$ and $R$ a family of root models $R(c)$, $c\in\mathrm{Comp}$.
\begin{enumerate}[label=(\alph*)]
\item $R$ is a \emph{reference} for $T$ if there are $C$ and $s_0>0$ such that for every $B\ge2$ and every $c$ the point
of $M_{B,c}$ and that of $R(c)$ have a common realisation (the \emph{join}) on which, for every line index $r$, the
difference of the multidegrees of the lines of $M_{B,c}$ and $R(c)$ is realised by a generic unit $u_r$ with a vertical
bound $C\log B$ (that is, $m u_r^{\pm q}$ are integral in the corresponding charts for some $q,m\ge1$ with $\log m\le
qC\log B$); and for every dominant $\mu$ with $\|\mu\|\le2$ and $0<s\le s_0$, at every complex point of the join and for
all values of the units $u_r$ there, the sum over $r$ of $\volthmcoef(\mu,s)(r)$ times the difference of the logarithms of
the metrics of the lines of $M_{B,c}$ at $\mathrm{scale}(s)$ and of $R(c)$ at $1$, corrected by $\log|u_r|$, lies in
$[-C\log(1/s)-C,\,C]$.
\item A root model $M$ has a \emph{flat unit} if for every $\tau$ there are $k\ge1$ and a generic unit realising the
multidegree $k\sum_ia(\tau,i)$, with a vertical bound.
\item A root model $M$ of type $(\Psi,p+q)$ has the \emph{fibre input} if there is $C\ge0$ such that for every number field
$L$ and every $L$-point $w$ of $P_M$ there are a finite extension $L'$, a Hermitian vector bundle $\mathcal E$ over
$\OO_{L'}$ and an orthogonal full flag pair in $\mathcal E$ (for each $\tau$ flags of lengths $p$ and $q$) such that the
heights at $w$ of the Chern-root lines $x_1(\tau,i)$, $i<p$, differ by at most $C$ from the degrees of the graded pieces of
the first flag, those of $x_1(\tau,p+i)$ differ by at most $C$ from minus the degrees of the graded pieces of the second
flag (in reversed order), and the maximal slope of $\mathcal E$ is at most $C\log(2+C+\max(-\widehat\deg_n\mathcal E,0))+C$,
$\widehat\deg_n$ the normalised degree.
\end{enumerate}
\formallabel{GlobalTheory.CompanionTheory.Ar.Ht.HasRef}
\formallabel{GlobalTheory.CompanionTheory.Ar.Ht.FlatUnit}
\formallabel{GlobalTheory.CompanionTheory.Ar.Ht.FibreInput}
\end{definition}

\begin{proposition}[Heights of points]\label{pf:volthm:prop:I5}
Assume $(\mathrm{H}_0)$; let $E$ be abelian over $\Q$, $Y$ global, $o$ a good origin of $T$ on $Y$ satisfying
$(\mathrm{X})$, $\volthmII$ full and $p,q\ge1$. Then $T$ satisfies $\mathrm{I}_5$ and $\mathrm{I}_6$.
\formallabel{GlobalTheory.T.good_tower_I5}
\formallabel{GlobalTheory.CompanionTheory.Ar.Top.i5_and_i6_of_isGlobal_of_goodOrigin}
\formallabel{RootFieldNumerics.global_model_tower_fifth_condition}
\formallabel{GlobalTheory.T.tower_intersection_of_sandwich_split}
\formallabel{GlobalTheory.CompanionTheory.Ar.TopGood.tower_good_of_sandwich_split}
\formallabel{GlobalTheory.CompanionTheory.Ar.TopGood.tower_intersection_of_ample_flat}
\end{proposition}

\begin{proof}
Let $R(a)$ be the reference root models of $o$ (with root origins at the empty set of primes and flags of degree $1$).
We use the following results; except for the first two, they are results of the development which we do not
reproduce.
\begin{enumerate}[label=(\arabic*)]
\item \emph{Sandwich} (\lref{GlobalTheory.T.exists_sandwich_of_hodge_tate_free}, which states, given the
Hodge--Tate statement of $(\mathrm{H}_0)$, that a Hodge component which is Hodge--Tate free at a prime $p$ of its set of
primes is \emph{sandwiched} at $p$: there is $\kappa$, equal to $4$ if $p=2$ and to $\zeta-1$ for a primitive $p$-th
root of unity $\zeta$ otherwise, such that at every point of the fibre over $p$ the Hodge--Tate and the de Rham lattices
and their duals lie between each other up to multiplication by $\kappa$): since the Hodge components of $o$ carry vector bundles, they are sandwiched at
all primes of $S_o(B)$, for all $B$ and $a$.
\item \emph{Split components} (\lref{RootFieldNumerics.splitComps_of_origin}): the components of $Y$ are split, i.e.\
the algebraic constants at their generic points lie in $c.F$; and a flat generator of the line bundles of the flag
exists, unique up to constants (\lref{GlobalTheory.T.flat_zero_from_resolution_and_gaga}, from resolution in projective
coordinates and GAGA for sections in $(\mathrm{H}_0)$).
\item \emph{Reference} (\lref{GlobalTheory.T.has_reference_of_sandwich_all_primes}): if the Hodge components of $o$ are
sandwiched at all primes of $S_o(B)$ for all $B$ and $a$, then $R$ is a reference for $T$.
\item \emph{Flat units} (\lref{GlobalTheory.T.flat_unit_of_split_root_lines}): given the flat generator of (2) and
split components, each $R(a)$ has a flat unit.
\item \emph{Fibre input} (\lref{GlobalTheory.T.root_model_flag_height}): from parts of the moduli inputs, the
extension of homomorphisms of semi-abelian schemes, resolution in projective coordinates and the bound for the maximal
slope of Lie algebras of $(\mathrm{H}_0)$, each $R(a)$ has the fibre input with $(p,q)$.
\item \emph{Lower bound} (\lref{GlobalTheory.T.tower_i5_of_fibre_input}): if $R$ is a reference for $T$ and each $R(a)$
has an ample reference line (generated by global sections at all points, a property of good origins), a flat unit and the
fibre input with $(p,q)$, then $T$ satisfies $\mathrm{I}_5$.
\item \emph{Upper bound} (\lref{GlobalTheory.T.tower_i6_of_reference}): if $R$ is a reference for $T$, then $T$ satisfies
$\mathrm{I}_6$.
\end{enumerate}
By (1) and (3), $R$ is a reference for $T$; by (2), (4), (5), the ampleness of the reference lines of a good origin and
(6), $T$ satisfies $\mathrm{I}_5$; by (7), $T$ satisfies $\mathrm{I}_6$. (The full intersection system is only used through
the axioms of an intersection family; the hypotheses that $Y$ is global and that $o$ satisfies $(\mathrm{X})$ are not
used.)
\end{proof}

\begin{definition}[Antisymmetry and squares]\label{pf:volthm:def:antisym}
Let $1\le j<N$.
\begin{enumerate}[label=(\alph*)]
\item A polynomial $P$ on $\R^{\Psi\times[N]}$ is \emph{antisymmetric at $j$} if $P\circ s_{\tau,a,a+1}=-P$ for all $\tau$ and
$a$ with $a+1\ne j$ ($s_{\tau,a,a+1}$ the simple reflection).
\item For a realisation $\mathrm{Re}$ of an analytic system $S$ with index set $\Psi\times[N]$ on an arithmetic variety,
$P_{\mathrm{Re}}$ is the polynomial in the weights whose value at $\lambda$ is the polarised intersection number of the
realisation at $\lambda$, $\deg_S$ is the degree polynomial of $S$ and $\vol(S)$ its volume. A \emph{flag system} for $j$
on $\mathrm{Re}$ consists of the conditions $j\le N$, that $S$ has dimension $n$ and is uniformised for $j$ (as in the
definition of uniformised real towers in \cref{pf:volthm:sec:conditions}), and of a \emph{square system} for every
simple reflection $s_{\tau,a,a+1}$ with $a+1<j$ or $j\le a$ (the data of the development from which the antisymmetry
under that reflection is derived).
\item The \emph{antisymmetry statement} $(\mathrm{AS})$ is: given $(\mathrm{AI})$, for every full $\volthmII$, nonempty $\Psi$,
$N\ge2$ and every realisation $\mathrm{Re}$ of an analytic system $S$ of positive dimension with index set $\Psi\times[N]$
which carries a flag system for $j$, the polynomial $P_{\mathrm{Re}}-\vol(S)H_{\mathrm{Fl}}$ is antisymmetric at $j$.
\item A real tower $R$ of $T$ (systems $S_c$ and realisations $\mathrm{Re}_{B,c}$ of $P_{M_{B,c}}$) \emph{has squares} if
for every $B$ and $c$ there are $m>0$, an analytic system $S'$ and a realisation $\mathrm{Re}'$ of it carrying a flag
system for $j$, with $P_{\mathrm{Re}'}=m\,P_{\mathrm{Re}_{B,c}}$ and $\deg_{S'}=m\deg_{S_c}$.
\end{enumerate}
\formallabel{GlobalTheory.CompanionTheory.Ar.PB.Antisymmetry}
\formallabel{SquareStatements.SquaresS}
\formallabel{GlobalTheory.CompanionTheory.Ar.PB.FlagSys}
\end{definition}

\begin{lemma}[Antisymmetry from squares]\label{pf:volthm:lem:squares}
Assume $(\mathrm{AS})$, $(\mathrm{AG})$ and $(\mathrm{AI})$. Let $\volthmII$ be full, $\Psi$ nonempty and $T$ a root tower with a
real tower $R$ which has squares. Then $T$ satisfies $\mathrm{I}_8(\volthmhh^{\volthmII})$.
\formallabel{SquareStatements.AntisymmetryFibreSquares}
\formallabel{SquareStatements.antisymmetry_transposition_factor}
\end{lemma}

\begin{proof}
Fix $B$, $c$, $\tau$, $b=a+1\ne j$ and a weight $\lambda$; let $m,S',\mathrm{Re}'$ be as in the definition of squares, and $n_c\ge1$
the integer of the real tower. The following are results of the development which we do not reproduce. By the
arithmetic agreement (\lref{GlobalTheory.T.real_tower_arith_agreement}, with the convergence given by
$(\mathrm{AG})$), $n_c\,\volthmPP_{B,c}(\volthmlin(\mu)^{n+1})=P_{\mathrm{Re}_{B,c}}(\mu)$ for every weight $\mu$; by the
geometric agreement (\lref{GlobalTheory.T.real_tower_geometric_agreement}) at $\rho$, together with
$\deg_{S'}=m\deg_{S_c}$ and the facts that $\deg_{S'}=\vol(S')\volthmfd$ (\lref{GlobalTheory.CompanionTheory.Ar.PB.proportional},
applicable because $\mathrm{Re}'$ carries a flag system) and $\volthmfd(\rho)\ne0$, one gets $m\,n_c\vol(c)=\vol(S')$. Finally
$\volthmhh(\volthmlin(\mu)^{n+1})=H_{\mathrm{Fl}}(\mu)$ (\cref{pf:volthm:prop:hh}(a)). Hence
\[
G=P_{\mathrm{Re}'}-\vol(S')H_{\mathrm{Fl}}\quad\text{satisfies}\quad
G(\mu)=m\,n_c\bigl(\volthmPP_{B,c}(\volthmlin(\mu)^{n+1})-\vol(c)\volthmhh(\volthmlin(\mu)^{n+1})\bigr)
\]
for all $\mu$. By $(\mathrm{AS})$, $G(s_{\tau,a,b}\lambda)=-G(\lambda)$; as $m\,n_c>0$, this is \eqref{pf:volthm:eq:I8}.
\end{proof}

\begin{proposition}[The antisymmetry identity]\label{pf:volthm:prop:I8}
Assume $(\mathrm{H}_9)$, $(\mathrm{H}_0)$ and $(\mathrm{H}_2)$; let $o$ be a good origin of $T$ on $Y$ satisfying
$(\mathrm{X})$, $\volthmII$ full and $p,q\ge1$. Then $T$ satisfies $\mathrm{I}_8(\volthmhh^{\volthmII})$.
\formallabel{GlobalTheory.T.full_intersection_tower_I8_identity}
\formallabel{GlobalTheory.CompanionTheory.Ar.Top.full_intersection_tower_I8_relation}
\formallabel{SquareStatements.good_tower_antisymmetry}
\formallabel{SquareStatements.antisymmetry_of_singular_statement}
\formallabel{SquareStatements.good_tower_antisymmetry_of}
\formallabel{SquareStatements.sq_of}
\end{proposition}

\begin{proof}
$\Psi$ is nonempty (\lref{GlobalTheory.CompanionTheory.Ar.Top.HostNumbers.host_numbers}). The antisymmetry
statement $(\mathrm{AS})$ is \lref{GlobalTheory.T.antisymmetry'}, and $(\mathrm{AG})$ is
\lref{GlobalTheory.T.agreement_positive}; both are results of the development which we do not reproduce, the first
being the main analytic result behind $\mathrm{I}_8$. By \cref{pf:volthm:lem:squares} it suffices to show that the real
tower $R_o$ of $o$ has squares, with $j=p$. Fix $B$ and a component $c$; the primes dividing $\mathrm{disc}(E)$ lie in
$S_o(B)$. We use two further results of the development:
\begin{itemize}
\item \emph{Ties} (\lref{SquareStatements.tie_of}): given $(\mathrm{H}_0)$ and $(\mathrm{H}_2)$, at the generic point of the
Hodge component of $o$ at level $B$ and $c$ there are a root model with a Hodge tie over $Y$ at $S_o(B)$, of realisation
degree $1$, with ranks $4$ at the primes of $S_o(B)$ and a flag of positive degree, and a realisation $\mathrm{Re}'$ of an
analytic system $S'$ carrying a flag system for $p$ whose polynomials $P_{\mathrm{Re}'}$ and $\deg_{S'}$ are those of the tie.
\item \emph{Comparison of ties} (\lref{SquareStatements.stCmp}): given $(\mathrm{AI})$ and $(\mathrm{H}_9)$, for two Hodge
ties at the same generic point, with equal ranks at the primes of $S$, the first of realisation degree $1$ and the
second of realisation degree equal to the degree of its flag, the polynomials of intersection numbers and the degree
polynomials of the first are those of the second multiplied by the degree of the flag of the first.
\end{itemize}
The root origin of $o$ at level $B$ and $c$ defines a Hodge tie of its root model whose realisation degree is the degree
of its flag (\lref{GlobalTheory.T.root_model_origin_compatibility}; this uses the first part of $(\mathrm{X})$ and the
vector bundles of $o$), with ranks $4$ at the primes of $S_o(B)$ by the second part of $(\mathrm{X})$. Comparing it with the
tie of the first item gives $P_{\mathrm{Re}'}=m\,P_{\mathrm{Re}_{B,c}}$ and $\deg_{S'}=m\deg_{S_c}$ with $m>0$ the degree of the
flag of that tie. So $R_o$ has squares.
\end{proof}

\begin{proposition}[Semipositivity of shifted lines]\label{pf:volthm:prop:A1}
Let $\volthmII$ be full (and assume $(\mathrm{AI})$), and let $T$ be a root tower with a real tower $R$ which is
\emph{smooth} (every realisation $\mathrm{Re}_{B,c}$ is a holomorphic resolution without boundary, of the smooth model of
$P_{M_{B,c}}$), \emph{flat} (for every $c$ and $\tau$ the potentials of the fundamental lines of index $(\tau,N-1)$ of $S_c$
and their negatives are plurisubharmonic, for all $t$) and \emph{positive} (the reference potential of $S_c$ is
positive). Then $T$ satisfies $\mathrm{A}_1$.
\formallabel{GlobalTheory.T.real_tower_A1}
\formallabel{GlobalTheory.CompanionTheory.Ar.Top.tower_A1_of_real_tower}
\formallabel{RootFieldNumerics.real_tower_mean_slope_condition}
\formallabel{RootFieldNumerics.A1.comp_psh}
\end{proposition}

\begin{proof}
For a weight $\mu$ and $s\in\R$ let $c_{\mu,s}$ be the coefficients $c(\infty)=s/k_A$ and $c(\tau,i)=\mu_{\tau,i}-\mu_{\tau,i+1}$
(with $\mu_{\tau,N}=0$), and for an analytic system $S$ with potentials $\mathrm{pot}_r(t)$ let
$\Phi_{\mu,s}(t)=\sum_rc_{\mu,s}(r)\,\mathrm{pot}_r(t)$.

\emph{Step 1: plurisubharmonicity on the analytic systems.} Fix $b>0$ and a component $c$, and let $k=k_A\ge1$. We claim:
there is $s_0(c)>0$ such that for $0<s\le s_0(c)$, $\mathrm{scale}(s)$ is at most the regularisation threshold of $S_c$ and
$\Phi_{\mu,s}(\mathrm{scale}(s))$ is plurisubharmonic on $S_c$ for every dominant $\mu$ with $|\mu_{\tau,i}|\le b$. Indeed,
the reference potential $\theta_\infty$ is positive and the potential of the reference line does not depend on $t$, so
by compactness there is $c_0>0$ with $\theta_\infty-c_0\Theta$ plurisubharmonic, $\Theta$ a fixed positive form
(\lref{RootFieldNumerics.A1.compare_global}). The coefficients $k\,c_{\mu,s}(r)$ are nonnegative except possibly at
the indices $(\tau,N-1)$, whose potentials are pluriharmonic by flatness: for dominant $\mu$,
$\mu_{\tau,i}-\mu_{\tau,i+1}\ge0$ for $i+1<N$, and $k\,c(\infty)=s>0$. Their sum of absolute values is at most
$1+\sum_r2kb$. By \lref{RootFieldNumerics.A1.psh_sum_potentials_small_parameter} (which states that, for such
coefficient vectors $c'$ with $c'(\infty)=s$, the combination $\sum_rc'(r)\mathrm{pot}_r(\mathrm{scale}(s))-\frac{sc_0}2\Theta$
is plurisubharmonic for all small $s$), $k\,\Phi_{\mu,s}(\mathrm{scale}(s))-\frac{sc_0}2\Theta$ is plurisubharmonic, hence
so is $\Phi_{\mu,s}(\mathrm{scale}(s))$, as $\Theta$ is positive and $k>0$.

\emph{Step 2.} Let $a>0$ and $b=a+\|\nu\|_1$, and let $s_0=\min(1,\min_cs_0(c))>0$ (finitely many components). Let $B\ge2$,
$c$, $\|\lambda\|_1\le a$, $0<s\le s_0$ with $\mu=\lambda_s=\lambda-s\nu$ dominant. Then $|\mu_{\tau,i}|\le\|\lambda\|_1+s\|\nu\|_1\le b$, so
$\Phi_{\mu,s}(\mathrm{scale}(s))$ is plurisubharmonic on $S_c$, and $k_A$ does not depend on $B$. By summation by parts,
$\sum_i\mu_{\tau,i}y_i=\sum_i(\mu_{\tau,i}-\mu_{\tau,i+1})\sum_{k\le i}y_k$; hence the weight
$\sum_r\volthmcoef(\mu,s)(r)\log h_{\ell_{\mathrm{scale}(s)}(r)}$ of $M_{B,c}$ equals $\sum_{\tau,i}c_{\mu,s}(\tau,i)\log
h_{\mathrm{fund}(\tau,i)}+c_{\mu,s}(\infty)\log h_A$, where $\log h_{\mathrm{fund}(\tau,i)}=-\sum_{k\le i}\log h_{x_{\mathrm{scale}(s)}(\tau,k)}$
is the regularised fundamental line of the realisation. In the charts of the resolution, the logarithms of the regularised
fundamental lines and of $A$ are the pull-backs of the potentials $\mathrm{pot}_r(\mathrm{scale}(s))$ of $S_c$ along the
holomorphic map of the realisation, so the weight is the pull-back of $\Phi_{\mu,s}(\mathrm{scale}(s))$ there. Since the
realisation is smooth and holomorphic, this gives the nonnegativity of the Levi form of the weight on the tangent cones at all complex points
of $P_{M_{B,c}}$ (\lref{RootFieldNumerics.Levi.levi_chart_nonneg} and \lref{RootFieldNumerics.Levi.levi_nonneg}),
which is $\mathrm{SP}(\mathrm{scale}(s),\lambda_s,s)$. (The inputs $(\mathrm{AI})$ and the fullness of $\volthmII$ are not used.)
\end{proof}

% ---------------------------------------------------------------------------
\subsection{Proof of the volume theorem}\label{pf:volthm:sec:proof}

\begin{proof}[Proof of \cref{pf:volthm:main}]
Let $C_h,C_u$ be as in the hypothesis, and $C=2\max(C_h+C_u/2,0)+3\cdot4+5\ge0$. Let the data be as in
\cref{pf:volthm:def:VolX} with $Y$ global, $o$ satisfying $(\mathrm{X})$, $\volthmII$ full, $q\ge1$, $3q+1\le p+q$ and
$|\mathrm{bad}|\le4|\Psi|N^2$. Then $p\ge2q+1\ge1$ and $q\le p$. By
\lref{GlobalTheory.CompanionTheory.Ar.Top.HostNumbers.host_numbers}, $\mathrm{rd}_E\ge2$ and $\Psi$ is nonempty.
From $(\mathrm{H}_0)$ we use the components $(\mathrm{AI})$, the Hilbert--Samuel package (which contains
$(\mathrm{Z}^-)$ and $(\mathrm{Z}^+)$),
the height package (which contains $(\mathrm{Z}^-_{\mathrm{rel}})$ and
$(\mathrm{GH})$) and $(\mathrm{AP})$. We check the hypotheses of
\cref{pf:volthm:thm:tower} for $T$, the full system $\volthmII$, $C_T=\max(C_h+C_u/2,0)$, $C'=4$,
$r=\mathrm{rd}_E\ge1$ and $L=L_H$.
\begin{enumerate}[label=(\arabic*)]
\item $(\mathrm{FH})$ holds: this is \lref{GlobalTheory.T.flag_height}, which proves the homogeneity, the
  vanishing on walls and the positivity of $H_{\mathrm{Fl}}$ stated in $(\mathrm{FH})$ (not reproduced here).
\item $R_o$ is uniformised, so $T$ satisfies $\mathrm{Pr}$ (\cref{pf:volthm:prop:prop}, with
  $(\mathrm{HM})$).
\item $\mathrm{Cv}$, $\mathrm{I}_2$, $\mathrm{I}_1$ hold by \cref{pf:volthm:prop:I12} ($(\mathrm{AG})$,
  $(\mathrm{CO})$, and $R_o$ reference-linked).
\item $\mathrm{I}_5$ and $\mathrm{I}_6$ hold by \cref{pf:volthm:prop:I5}, $\mathrm{I}_8(\volthmhh)$ by \cref{pf:volthm:prop:I8}
  (using $(\mathrm{H}_9)$ and $(\mathrm{H}_2)$), $\mathrm{A}_1$ by \cref{pf:volthm:prop:A1} ($R_o$ is smooth, flat and
  positive), and $\mathrm{M}_1$, $\mathrm{M}_2$, $\mathrm{M}_0$, $\mathrm{I}_7$ by \cref{pf:volthm:prop:nef}.
\item \Cref{pf:volthm:prop:prop6} gives $\phi$ with $\mathrm{CS}(\mathrm{frob})$ and $\mathrm{EP}(\mathrm{frob})$ for
  $\mathrm{frob}(p,\tau,m)=\phi_p^{m+1}(\tau^+)$ (using $(\mathrm{H}_4)$ and $(\mathrm{H}_5)$), and
  \cref{pf:volthm:prop:propA} gives $\mathrm{GI}$ (using $(\mathrm{H}_7)$ and $(\mathrm{H}_8)$).
\item The elements of $\mathrm{bad}$ are prime (by the description of $\mathrm{bad}$ for tower origins), so
  \cref{pf:volthm:prop:propF} and \cref{pf:volthm:prop:flatdata}(b) give $\mathrm{FL}(4,\mathrm{rd}_E,L_H)$ (using
  $(\mathrm{H}_1)$ and $(\mathrm{H}_6)$).
\item $\mathrm{EV}$ holds by \cref{pf:volthm:prop:voleqproof}.
\item Since $j=p$, $3(N-j)+1=3q+1\le N$. By \cref{pf:volthm:prop:slope}(c) and (a), $T$ satisfies
  $\mathrm{SL}(C_h+C_u/2,\mathrm{rd}_E,L_H)$, hence by (b) $\mathrm{SL}(C_T,\mathrm{rd}_E,L_H)$.
\end{enumerate}
\Cref{pf:volthm:thm:tower} now gives $\mathrm{T}_1$, $\mathrm{AF}(\volthmhh)$, $\mathrm{T}_2(2C_T+3\cdot4+5)$ and
$\mathrm{MS}(2C_T+3\cdot4+5)$, with $2C_T+17=C$. Together with $\mathrm{Pr}$ and $\mathrm{EV}$ this is
the conclusion of $\sfname{VolumeTheoremX}$.
\end{proof}

% ---------------------------------------------------------------------------
\subsection{Remarks on the formal statements}\label{pf:volthm:sec:remarks}

\begin{remark}\label{pf:volthm:rem:formal}
\begin{enumerate}[label=(\arabic*)]
\item The constant of $\sfname{VolumeTheoremX}$ depends only on the constants $C_h,C_u$ of $(\mathrm{H}_3)_{C_h,C_u}$:
  $C=2\max(C_h+C_u/2,0)+17$; the remaining constants ($C'=4$ in $\mathrm{FL}$, the factor $1.6$ of
  \cref{pf:volthm:prop:remtest}, the factor $2$ of \cref{pf:volthm:prop:remlin}) are absolute.
\item The constants of $\mathrm{T}_1$, $\mathrm{AF}$ (the constant $C'$, the threshold $B_1$, and the parameters
  $\alpha,\beta$ for each $B$) are not uniform: they depend on the tower. Only $\mathrm{T}_2$ and $\mathrm{MS}$
  carry the uniform constant $C$.
\item The hypothesis $3q+1\le p+q$ is used for $\mathrm{EV}$ and for the slope condition (through
  $\min(j,N-j)=N-j=q$); the bound $|\mathrm{bad}|\le4|\Psi|(p+q)^2$ is used for the base degree settings.
\item The set $\mathrm{bad}$ attached to a root tower is only required to be a finite set of natural numbers containing $2$;
  for tower origins it is the set of primes dividing $2\,\mathrm{disc}(E)M'_l$, which is used in
  \cref{pf:volthm:prop:propF} (primality) and \cref{pf:volthm:prop:nef}(b).
\item $\volthmPP_M$ is defined by a choice of limit (\cref{pf:volthm:rem:limunder}), and the sums $R_j$ and the series in
  \eqref{pf:volthm:eq:delta} are formal sums which are $0$ for non-summable families (\cref{pf:volthm:rem:sums}); in
  all uses the limits exist and the families are absolutely summable.
\item The third of the five components of the height package in $(\mathrm{H}_0)$ is the true proposition and
  plays no role.
\item The hypothesis $(\mathrm{X})$ of $o$ consists of a local holomorphy condition for the point maps of the analytic
  systems and the condition that all Hodge components have rank $R_p=4$ at the primes of $S_o(B)$; the latter is used
  only through the equality of these ranks (\cref{pf:volthm:prop:prop6,pf:volthm:prop:propA,pf:volthm:prop:I8}).
\item Some formal statements carry hypotheses which their proofs do not use: in \cref{pf:volthm:prop:I5}, that $Y$ is
  global and that $o$ satisfies $(\mathrm{X})$; in \cref{pf:volthm:prop:flatdata}(b), $(\mathrm{X})$ and the fullness of
  $\volthmII$; in \cref{pf:volthm:prop:A1}, $(\mathrm{AI})$ and the fullness of $\volthmII$.
\item The bound for the remainder of the test polynomial is first proved for decompositions with rational coefficients;
  the passage to real decompositions uses that bound itself, applied to the decompositions $Q+tS$ for a decomposition
  $S$ of zero (\cref{pf:volthm:lem:syzygy}).
\end{enumerate}
\end{remark}

